\documentclass[11pt]{article}
\usepackage[margin=1in]{geometry}
\usepackage{amsmath,amssymb,amsthm}
\usepackage{mathrsfs}
\usepackage[hidelinks]{hyperref}
\usepackage{enumitem}
\setlength{\emergencystretch}{3em}

\newtheorem{theorem}{Theorem}
\newtheorem{lemma}[theorem]{Lemma}
\newtheorem{proposition}[theorem]{Proposition}
\newtheorem{corollary}[theorem]{Corollary}
\newtheorem{conjecture}[theorem]{Conjecture}
\newtheorem*{lemmaA}{Lemma A}
\theoremstyle{definition}
\newtheorem{definition}[theorem]{Definition}
\theoremstyle{remark}
\newtheorem{remark}[theorem]{Remark}

\newcommand{\R}{\mathbb{R}}
\newcommand{\Z}{\mathbb{Z}}
\newcommand{\N}{\mathbb{N}}
\newcommand{\Q}{\mathbb{Q}}
\newcommand{\ee}{\mathrm{e}}
\newcommand{\dist}{\operatorname{dist}}
\newcommand{\dimH}{\dim_{\mathrm{H}}}
\newcommand{\Lam}{\Lambda}

\title{Complete spectrum discretization, exact Hausdorff gauges,\\
and dimension phase transitions for Diophantine approximation\\
along the convergent denominators of \texorpdfstring{$e$}{e}}
% TODO before submission: replace with the real author name(s) and affiliation(s).
\author{[Author name]\\ \relax[Affiliation]}
\date{Integrated draft, November 23, 2026}

\begin{document}
\maketitle

\begin{abstract}
Let $(Q_j)_{j\ge1}$ be the selected convergent denominators of $\ee$
immediately preceding the partial quotients $2j$
($Q_1=1,Q_2=7$, $Q_{j+1}=(4j+2)Q_j+Q_{j-1}$), let $\|y\|=\dist(y,\Z)$, and
define the weighted approximation exponent
$\Lam(x)=\limsup_{j\to\infty} j\,\|Q_j x\|$. We prove that the full spectrum is
completely discrete:
$\Lam(\R)=\frac14\Z_{\ge0}\cup\{\infty\}=\{0,1/4,1/2,3/4,1,\dots\}\cup\{\infty\}$,
so every open interval $(k/4,(k+1)/4)$ ($k\in\Z_{\ge0}$) is an unconditional
empty spectral gap, and $\Lam(x)=0$ exactly on the countable group
$L=\Z+\Z\frac{\ee+1}{2}$.

Our main metric results measure the critical levels and their fine structure.
For an integer $k\ge1$ let $h_k$ be the gauge with $h_k(1/Q_n)=(2k+1)^{-n}$,
interpolated linearly. On the circle $\R/\Z$, the exact level $\{\Lam=k/4\}$
has infinite but $\sigma$-finite $\mathcal H^{h_k}$ measure, while
$\{\Lam<k/4\}$ is $\mathcal H^{h_k}$-null. The gauge $h_k$ is exact: the level
has infinite $\mathcal H^{h_k}$ measure in every nonempty open arc, every gauge
of smaller order gives $\{\Lam\le k/4\}$ measure zero, and for every gauge of
larger order $\{\Lam=k/4\}$ is not $\sigma$-finite. For arbitrary gauges,
including oscillating ones, the three measure regimes are determined by
$\liminf_n (2k+1)^n g(1/Q_n)$, or by $\liminf_n 2^n g(1/Q_n)$ for the
unsigned exact-limit set. For finite-state tail constraints the base is
the largest Perron root among components supporting infinitely many
extreme defects. Furthermore, at each
positive level $c=k/4$ ($k\ge1$) we establish a sharp three-way splitting
between the limsup level (digit alphabet of size $2k+1$, critical gauge base
$2k+1$), the unsigned exact limit set $\{x:j\|Q_jx\|\to k/4\}$ (universal binary
alphabet $\{\pm k\}$, critical gauge base $2$ \emph{independently} of
$k\ge1$), and each signed exact limit set
$\{x:j(Q_jx-\mathrm{round}(Q_jx))\to\pm k/4\}$ (which collapses to a single
countable coset of $L$). More generally, for any polynomial-growth recurrence
$Q_{j+1}=(cj^r+d)Q_j+Q_{j-1}$ ($r,c\in\Z_{\ge1}$, $d\in\Z$), the critical
exponent $\Lam^{(r,c,d)}(x)=\limsup_{j\to\infty}j^r\|Q_jx\|$ has full
spectrum $\frac1c\Z_{\ge0}\cup\{\infty\}$, intrinsic rank-$2$ zero-level
lattice $L^{(r,c,d)}$, critical gauges $h_k^{(r,c,d)}(1/Q_n)=(2k+1)^{-n}$
with $1/r$-scaled logarithmic order and positive-level three-way splitting, and
subcritical Hausdorff dimension phase transition $1-\alpha/r$ ($0\le\alpha<r$).
Even more generally, for \emph{any} growing-quotient recurrence
$Q_{j+1}=A_jQ_j+Q_{j-1}$ ($A_j\in\Z_{\ge2}$, $A_j\to\infty$), the normalized
approximation functional $\Lam_G(x)=\limsup_{j\to\infty}A_j\|Q_jx\|$ agrees
identically with the Kraaikamp--Liardet error ratio, has full discrete spectrum
$\Z_{\ge0}\cup\{\infty\}$ with zero-level lattice $L_G=\frac{1}{g}(\Z+\Z\theta)$
($g=\gcd(Q_{J_0-1},Q_{J_0})$), universal critical and binary gauges on $\R$ and
$\R/\Z$, subcritical Hausdorff dimension $(1-\tau)/(1+\tau\gamma_{\sup})$ for
$\Lam_{G,\tau}(x)=\limsup_{j\to\infty}A_j^\tau\|Q_jx\|$ ($0\le\tau<1$, where
$\gamma_{\sup}=\limsup_{j\to\infty}\frac{\log A_j}{\log Q_j}\in[0,\infty]$),
and a supercritical ($\tau>1$) secondary phase transition at
$\tau_{\mathrm{crit}}=1+1/\gamma_{\sup}$ where finite sublevel sets jump from
the rank-$2$ lattice $L_G$ to the strict rank-$1$ rational sublattice
$\frac{1}{g}\Z\subsetneq L_G$.

Next, every finite level set $\{\Lam\le c\}$ has Hausdorff dimension $0$.
Hence $\{\Lam<\infty\}$ has dimension $0$, and $\Lam=\infty$ Lebesgue-almost
everywhere (with Hausdorff dimension $1$). For $0\le\alpha<1$, the sublevel set $\{\Lam_\alpha\le c\}$ for $c\ge0$
has Hausdorff dimension $1-\alpha$. The same holds for the set where
$j^\alpha\|Q_jx\|\to c$ for $c\ge0$ if $\alpha>0$, and for
$c\in[0,1/2]$ if $\alpha=0$. The unweighted case also follows from work
of Bugeaud, Liao and Rams. For ordinary continued fractions,
we determine the Hausdorff dimensions of Hermite window-count constraints,
show that eventual periodicity of the partial quotients modulo $m$ ($m\ge2$) has
the same Hausdorff dimension as the fixed residue alphabet $F_{m,1}$ (strictly
between $357/512$ and $179/256$ for $m=4$), prove the global strict separation
$\dimH F_{m,m}<\dimH F_{m,1}=\dimH E_m$ for \emph{all} moduli $m\ge2$ via an
exact resolvent identity, and establish the two-term asymptotic expansion
$\dimH F_{m,r}=\frac12+\frac{C_r}{2m}-\frac{C_r^2\log m}{2m^2}+O(m^{-2})$ as
$m\to\infty$ for fixed $r\ge1$ (where $C_r=\sum_{k=1}^\infty 1/q_k(r)$ is the
reciprocal metallic-continuant series and $C_1=\psi_{\mathrm{Fib}}=\sum_{k=1}^\infty 1/F_k\approx3.3598856662$ is the reciprocal Fibonacci constant), while
$\dimH F_{m,m}=\frac12+\frac{1}{2m}-\frac{\log m}{2m^2}+O(m^{-2})$. More generally,
for finite multi-residue alphabets $F_{m,S,B}$ ($S\subset\Z_{\ge1}$,
$B\subset\Z_{\ge0}$, $d=|S|+|B|\ge1$), Theorem~\ref{thm:multi-residue-asymptotic}
gives a self-contained, Bowen-free proof of the classical identity
$\dimH E_S=s_0(S)$ and establishes the universal phase transition
$\dimH F_{m,S,B}\to\max(1/2,s_0(S))$, a three-term
Euler--Mascheroni subcritical expansion to $O((\log m)^2/m^3)$ with coefficient
$\frac12 d^2 C_S(\gamma C_S-C_S')$, and the supercritical power-law convergence
rate $\Theta(m^{-2s_0(S)})$ for $s_0(S)>1/2$.

The classification of $\{\Lam=0\}$, the initial and full spectral-gap ladder,
Theorem~\ref{thm:gauge},
Corollaries~\ref{cor:exact-gauge} and~\ref{cor:liminf-gauges},
Remark~\ref{rem:gauge-scale}, all three parts (i)--(iii) of
Theorem~\ref{thm:trichotomy} on both $\R$ and $\R/\Z$ (\texttt{UnitAddCircle},
including the binary exact-limit gauge regime and the signed-coset collapse
$E_{\mathrm{lim}}^\pm(k/4)=x_k^\pm+L$), the dimension-zero theorem for all
finite sublevels (and full Lebesgue measure of $\{\Lam=\infty\}$), both the upper
and lower bounds for the subcritical dimension $1-\alpha$, the complete
general polynomial-growth recurrence theory ($r\ge1$)---Lemmas~\ref{lem:poly-growth}--\ref{lem:poly-lattice},
Theorems~\ref{thm:general-rec}--\ref{thm:poly-phase-transition},
Corollary~\ref{cor:poly-exact-gauge}, and Remark~\ref{rem:poly-gauge-scale}---and
its universal extension to arbitrary growing-quotient recurrences
$Q_{j+1}=A_jQ_j+Q_{j-1}$ (Theorems~\ref{thm:growing-rec-spectrum}
and~\ref{thm:growing-rec-dim-phase}, and
Corollary~\ref{cor:growing-rec-three-regimes}), together
with all Hausdorff-dimension statements of Theorem~\ref{thm:windows} (both the
exact $\dimH=1$ prescribed-parity statement and the exact $\dimH=1/2$ upper and lower bounds for eventually bounded and prescribed
increments), together with all of Theorems~\ref{thm:short-ae} and~\ref{thm:no-short-dim}
(both the almost-everywhere logarithmic short-window frequency, the
Hausdorff-dimension identity $\dimH=1$, and the Lebesgue-nullity deduction
pointwise for mathlib's regular continued fractions), all parts~(1)--(3) of
Theorem~\ref{thm:finite-state-tail} (under distinct outgoing edge labels at each
vertex and given explicit positive right Perron--Frobenius eigenvectors on marked
strongly connected components), as well as the general-modulus ($m\ge2$)
uniform $24$-Lipschitz coordinate-shift sandwich, eventual-periodicity dimension
equality, unconditional transfer-operator Hausdorff-dimension criterion
(bypassing Bowen-pressure axioms), strict lower bound $\dimH F_{m,r}>1/2$, upper
bound $\dimH F_{4,1}\le 4/5$, step-function certificate soundness for
Theorems~\ref{thm:residues} and~\ref{thm:residues-general}, and the complete
resolvent identity, global strict separation $\dimH F_{m,m}<\dimH F_{m,1}=\dimH E_m$
for all $m\ge2$, and two-term reciprocal-Fibonacci/metallic asymptotic expansions of
Theorem~\ref{thm:residues-asymptotic}, as well as the multi-residue phase
transition, Euler--Mascheroni three-term asymptotics, and supercritical
$\Theta(m^{-2s_0(S)})$ rate of Theorem~\ref{thm:multi-residue-asymptotic}
(\texttt{ResidueMultiAsymptotic.lean}, $24$ declarations audited in
\texttt{AuditResidueMultiAsymptotic.lean}) and the closed brackets of
Table~\ref{tab:residues-general} (\texttt{ResidueCertTable.lean}, $27$ declarations
audited in \texttt{AuditResidueCert.lean})---are machine-checked in
Lean~4 using only the base axioms \texttt{propext}, \texttt{Classical.choice}
and \texttt{Quot.sound}.
Only the finite-dimensional linear-algebraic Perron--Frobenius eigenvector
existence input in Theorem~\ref{thm:finite-state-tail} and the Python-executed
open-interval strictness verification of the narrow $m\in\{2,3,4,5,6\}$ brackets
via the pressure formalism remain outside the Lean kernel
(Section~\ref{sec:verification}, $937$ main audited declarations across $134$
Lean files in \texttt{Audit.lean} plus standalone audits across $344$ total Lean files).
The classification and the first gap are a subsequence analogue of a criterion
of Kraaikamp and Liardet. The work is a spin-off of a project on the
irrationality of $e+\pi$ and does \emph{not} resolve that open problem.
\end{abstract}

\section{Introduction}

This article studies the metric theory of Diophantine approximation along a
\emph{sparse, super-exponentially growing} subsequence of denominators: the
selected convergent denominators $(Q_j)_{j\ge1}$ of $\ee$ (notation in the
Setup below). We first show in Theorem~\ref{thm:full-spectrum} that the integer
three-term recurrence for $(Q_j)$ forces complete discretization of the weighted
approximation spectrum:
\[
\Lam(\R)=\frac14\Z_{\ge0}\cup\{\infty\}
=\left\{0,\frac14,\frac12,\frac34,1,\dots\right\}\cup\{\infty\},
\]
creating an infinite family of equidistant empty spectral gaps
$(k/4,(k+1)/4)$ of width $1/4$ and forcing the sublevel sets
$E_c=\{x:\Lam(x)\le c\}$ to be constant across each gap
($E_c=E_{k/4}$ for $c\in[k/4,(k+1)/4)$, Corollary~\ref{thm:sub}).

Our main metric result at the critical weight, Theorem~\ref{thm:gauge},
measures each positive quantized level of $\Lam$: for each integer $k\ge1$
there is an explicit Hausdorff gauge $h_k$ with $h_k(1/Q_n)=(2k+1)^{-n}$, built
from the recurrence itself, for which the exact level $\{\Lam=k/4\}$ has
infinite but $\sigma$-finite Hausdorff measure $\mathcal H^{h_k}$ on the circle
$\R/\Z$, while the strict sublevel $\{\Lam<k/4\}$ is $\mathcal H^{h_k}$-null.
This differs both from the $h$-sets of Cabrelli, Mendivil, Molter and
Shonkwiler \cite[\S5, Theorem~3]{CMMS}, whose gauge gives a finite positive
measure, and from the classical Lagrange level at value $3$ studied by Moreira,
Erazo and Angelini \cite{MEA}, which has either zero or non-$\sigma$-finite
measure for every dimension function. Corollary~\ref{cor:exact-gauge} shows
that $h_k$ is the exact critical gauge of the level $k/4$: smaller gauges give
measure zero even on the non-strict sublevel $\{\Lam\le k/4\}$, larger gauges
destroy $\sigma$-finiteness on $\{\Lam=k/4\}$, and the $\mathcal H^{h_k}$
measure of $\{\Lam=k/4\}$ is infinite in every nonempty open arc.
Corollary~\ref{cor:liminf-gauges} classifies every nondecreasing gauge by
its lower limit along the recurrence scales, even when the normalized
cost oscillates. Theorem~\ref{thm:finite-state-tail} shows how this
classification changes under constraints on consecutive defects.
The largest Perron root is taken only over components in which a path
can revisit an extreme defect infinitely often.
Furthermore, Theorem~\ref{thm:trichotomy} establishes a sharp three-way
splitting at each positive spectral level $k/4$ ($k\ge1$) comparing the limsup
level $\{\Lam=k/4\}$ (gauge base $2k+1$), the unsigned exact limit set
$\{x:j\|Q_jx\|\to k/4\}$ (universal critical gauge base $2$ for \emph{all}
$k\ge1$), and the signed exact limit sets
$\{x:j(Q_jx-\mathrm{round}(Q_jx))\to\pm k/4\}$ (each of which collapses to a
single countable $L$-coset). Theorems~\ref{thm:general-rec}--\ref{thm:poly-phase-transition}
extend the quantization, reference gauges, three-way splitting and
dimension results to general polynomial-growth recurrences
$Q_{j+1}=(cj^r+d)Q_j+Q_{j-1}$ ($r,c\in\Z_{\ge1}$, $d\in\Z$), replacing the
$\ee$-convergent lattice $L$ by an intrinsic rank-$2$ lattice $L^{(r,c,d)}$
spanned by fundamental solutions and establishing the discretized spectrum
$\frac1c\Z_{\ge0}\cup\{\infty\}$, exact gauges with $1/r$-scaled logarithmic
asymptotics, positive-level trichotomy, and subcritical dimension phase
transition $1-\alpha/r$ ($0\le\alpha<r$). In Section~\ref{sec:growing-rec}
(Theorems~\ref{thm:growing-rec-spectrum} and~\ref{thm:growing-rec-dim-phase},
and Corollary~\ref{cor:growing-rec-three-regimes}), we remove the polynomial
growth hypothesis altogether and establish the complete spectral, gauge, and
Hausdorff-dimension phase diagram for arbitrary integer recurrences
$Q_{j+1}=A_jQ_j+Q_{j-1}$ with $A_j\ge2$ and $A_j\to\infty$, uncovering both the
universal subcritical distortion law $\dimH=(1-\tau)/(1+\tau\gamma_{\sup})$ and
a secondary supercritical phase transition at $\tau_{\mathrm{crit}}=1+1/\gamma_{\sup}$
where the level sets collapse from the rank-$2$ lattice $L_G$ to the rank-$1$
rational sublattice $\frac{1}{g}\Z$.

On the coarse scale of Hausdorff dimension, Theorem~\ref{thm:finite} shows that
every finite level set $E_c$ has Hausdorff dimension $0$, so $\{\Lam<\infty\}$
has dimension $0$ while its complement $\{\Lam=\infty\}$ has full Lebesgue
measure and Hausdorff dimension $1$. At subcritical weights
$\Lam_\alpha(x)=\limsup_{j\to\infty}j^\alpha\|Q_jx\|$ ($0\le\alpha<1$),
Theorems~\ref{thm:upper} and~\ref{conj:weight} give exact Hausdorff dimension $1-\alpha$ for
$\{x:\Lam_\alpha(x)\le c\}$ for every $c\ge0$, and for
$\{x\in[0,1]:j^\alpha\|Q_jx\|\to c\}$ for every admissible target $c$.
For $\alpha=0$ such targets are $0\le c\le 1/2$.
The proof uses a covering bound and a shifted-grid Moran construction;
the case $\alpha=0$ also follows from Bugeaud, Liao and Rams \cite{BLR}, and
both bounds for $0\le\alpha<1$ are machine-checked in Lean~4. Finally, in
Section~\ref{sec:windows-residues}, Theorems~\ref{thm:windows}--\ref{thm:no-short-dim}
determine the Hausdorff dimensions and almost-sure frequency of Hermite
continued-fraction window constraints, Theorem~\ref{thm:residues} proves
that eventual periodicity of ordinary continued-fraction digits modulo $4$ has
the same Hausdorff dimension as the fixed alphabet $1\bmod 4$, bounded strictly
between $357/512$ and $179/256$, Theorem~\ref{thm:residues-general} extends this
to every modulus $m\ge2$ (via a uniform $24$-Lipschitz sandwich, an unconditional
transfer-operator Hausdorff-dimension framework, and exact integer certificates
for $m\in\{2,3,4,5,6\}$), Theorem~\ref{thm:residues-asymptotic} proves the
global strict separation $\dimH F_{m,m}<\dimH F_{m,1}=\dimH E_m$ for all $m\ge2$
together with the exact two-term large-$m$ asymptotic expansion governed by the
reciprocal Fibonacci constant $\psi_{\mathrm{Fib}}=\sum_{k=1}^\infty 1/F_k$ and
its metallic-continuant generalizations $C_r=\sum_{k=1}^\infty 1/q_k(r)$, and
Theorem~\ref{thm:multi-residue-asymptotic} establishes the multi-residue
phase transition $\lim_{m\to\infty}\dimH F_{m,S,B}=\max(1/2,s_0(S))$, a
self-contained Bowen-free proof of the classical Poincar\'e-abscissa dimension
formula for finite alphabets $E_S$,
three-term Euler--Mascheroni subcritical asymptotics to $O((\log m)^2/m^3)$,
and the two-sided supercritical power-law rate $\Theta(m^{-2s_0(S)})$.

These results build on the zero-level classification and initial spectral gap
(Section~\ref{sec:class}), which arose as machinery in a separate formalization
project on the (open) irrationality of $e+\pi$ \cite{Phase1}.

\subsection{Setup}

The continued fraction of $\ee$ is $[2;1,2,1,1,4,1,1,6,\dots]$. Let $P_j/Q_j$
be the convergent immediately preceding the partial quotient $2j$; the first
few are $3/1,19/7,193/71,2721/1001$. The denominators satisfy
\begin{equation}\label{eq:rec}
Q_{j+1}=(4j+2)Q_j+Q_{j-1}\quad(j\ge2),\qquad Q_1=1,\ Q_2=7,
\end{equation}
so $Q_{j+1}/Q_j\sim 4j$ and $\log Q_j\sim j\log j$ (super-exponential growth).
Write $\|y\|=\dist(y,\Z)$ and define
\[
\Lam(x):=\limsup_{j\to\infty} j\,\|Q_j x\|\in[0,\infty].
\]
More generally, for $\alpha\ge0$ put
$\Lam_\alpha(x):=\limsup_{j\to\infty} j^\alpha\|Q_j x\|\in[0,\infty]$,
so that $\Lam=\Lam_1$.
Let $L:=\Z+\Z\frac{\ee+1}{2}$, a countable dense subgroup of $\R$.

\subsection{Relation to prior work}\label{sec:prior}

The \emph{phenomenon} (a spectral gap plus zero-dimensional sub-threshold level
sets) is classical: in the classical Lagrange spectrum, the spectrum below
$\sqrt5$ (and below $3$) is countable (Hausdorff dimension $0$), and Hall's ray
gives a half-line. Moreira \cite{Moreira} studies the Hausdorff dimension of
portions of the classical Markov and Lagrange spectra below a threshold. Lima
and Moreira \cite{LimaMoreira} prove phase transitions for the \emph{dynamical}
Markov and Lagrange spectra of generic horseshoes; the corresponding question
for the classical spectra is left open there. Note that our level sets $E_c$ are
subsets of the space of inputs $x\in\R$.

The closest prior continued-fraction result is a theorem of Kraaikamp and
Liardet \cite{KL}, quoted as Lemma~4.1 in \cite{BDS}. Let $\alpha$ be irrational
with convergents $p_n/q_n$. Then a real $\beta$ lies outside
$\Z+\Z\alpha$ if and only if
$\|q_n\beta\|>\tfrac14\,q_n|q_n\alpha-p_n|$ for infinitely many $n$. This
is a classification together with a gap, with constant $1/4$, for the
normalized quantity $\|q_n\beta\|/(q_n|q_n\alpha-p_n|)$ along \emph{all}
convergent denominators. Theorems~\ref{thm:class} and~\ref{thm:gap} can be
compared with this criterion in two complementary ways:
\begin{itemize}[leftmargin=1.8em]
\item \emph{As a subsequence along the convergents of $\ee$:} Passing from all
convergent denominators of $\ee$ to the subsequence $(Q_j)$ of every third
convergent denominator changes two things. First, the zero group grows from
$\Z+\Z\ee$ to $L$: for $x=(\ee+1)/2$ the quantity $j\|Q_jx\|$ tends to $0$, so
the infinitely many $n$ supplied by Kraaikamp and Liardet for $\alpha=\ee$ lie
outside the subsequence $(Q_j)$. Second, the normalization changes. Since
$2j+1<1/(Q_j|Q_j\ee-P_j|)<2j+2$, one has
$\limsup_{j\to\infty}\|Q_jx\|/(Q_j|Q_j\ee-P_j|)=2\Lam(x)$. Thus the gap
$\Lam\ge1/4$ reads $\limsup\ge1/2$ in their normalization, and the value
$\Lam(x_0)=1/4$ is attained by the explicit point $x_0$ in
Remark~\ref{rem:unweighted}. The two constants $1/4$ therefore refer to
different normalizations.
\item \emph{As the full convergent sequence of $\beta=(\ee+1)/2$:} Put
$R_j=(P_j+Q_j)/2\in\Z$. The fractions $R_j/Q_j$ are the convergents of
$\beta=(\ee+1)/2=[1;1,6,10,14,\dots]$, whose initial convergent pairs are
$(1,1)$ and $(2,1)$ and whose remaining partial quotients are $4j+2$ for
$j\ge1$. Applying the Kraaikamp--Liardet criterion \cite{KL} (in the form
recorded on p.~106 of \cite{BDS}) directly to $\beta$ gives
$\|Q_jx\|>\tfrac14 Q_j|Q_j\beta-R_j|$ infinitely often for $x\notin L$. Since
the continued-fraction error formula gives $jQ_j|Q_j\beta-R_j|\to1/4$, this
black-box application yields $\Lam(x)\ge1/16$ outside $L$, which already
recovers the zero-level classification $\{\Lam=0\}=L$, whereas the direct
three-term recurrence analysis in Theorems~\ref{thm:gap}
and~\ref{thm:full-spectrum} yields the sharp initial gap $1/4$ and the complete
quantization $\Lam(\R)=\frac14\Z_{\ge0}\cup\{\infty\}$.
\end{itemize}
For a general irrational $\alpha$ whose partial quotients tend to infinity, we
First treat polynomial-growth denominator recurrences
$Q_{j+1}=(cj^r+d)Q_j+Q_{j-1}$ in
Theorems~\ref{thm:general-rec}--\ref{thm:poly-phase-transition} and then, in
Section~\ref{sec:growing-rec} (Theorems~\ref{thm:growing-rec-spectrum}
and~\ref{thm:growing-rec-dim-phase}, and
Corollary~\ref{cor:growing-rec-three-regimes}), establish the complete unified
theory for arbitrary growing-quotient recurrences $Q_{j+1}=A_jQ_j+Q_{j-1}$ with
$A_j\in\Z_{\ge2}$ and $A_j\to\infty$. In particular, Theorem~\ref{thm:growing-rec-spectrum}(i)
proves the exact asymptotic identity
$\limsup_{j\to\infty}\|Q_jx\|/(Q_j|Q_j\theta-P_j|)=\limsup_{j\to\infty}A_j\|Q_jx\|$
matching the Kraaikamp--Liardet error normalization across all growing-quotient
recurrences, while Theorem~\ref{thm:growing-rec-dim-phase} reveals how the upper
logarithmic growth exponent $\gamma_{\sup}=\limsup_{j\to\infty}\frac{\log A_j}{\log Q_j}\in[0,\infty]$
governs both the subcritical Hausdorff-dimension distortion
$(1-\tau)/(1+\tau\gamma_{\sup})$ ($0\le\tau<1$) and the supercritical secondary
collapse from $L_G$ to $\frac{1}{g}\Z$ across $\tau_{\mathrm{crit}}=1+1/\gamma_{\sup}$.
For the modular continued-fraction alphabets $F_{m,r}=\{[0;a_1,a_2,\dots]:a_n\equiv r\pmod m\}$,
the underlying infinite-alphabet dimension theory is that of Mauldin and
Urba\'nski \cite{MU96} (see \cite[Theorem~2.3]{CLU}), and Chousionis, Leykekhman
and Urba\'nski \cite{CLU} proved that every arithmetic-progression alphabet has
full dimension spectrum. Asymptotic expansions of the Hausdorff dimension when
generators with small derivative are added to or removed from a conformal IFS
were developed by Das, Fishman, Simmons and Urba\'nski \cite{DFSU}, who
strengthened Hensley's expansion \cite{Hensley92} for partial quotients $\le N$
and Good's estimates \cite{Good41} for partial quotients $\ge N$. The families
$F_{m,S,B}$ fall within the scope of their perturbation theorem
\cite[Theorem~4.1]{DFSU} (fixed finite $S$, added generators $\ge m-\max B$);
in their notation the leading constant is $1/\bar c=\sum_{n\ge0}\sum_{w\in S^n}q(w)^{-1}=C_S$
when $s_0(S)<1/2$. Arithmetic-progression alphabets are not treated in
\cite{DFSU}, and Theorems~\ref{thm:residues-asymptotic}
and~\ref{thm:multi-residue-asymptotic} make the first terms explicit for them;
the $\log m/m^2$ coefficient has the same structure as Hensley's
$72/\pi^4=2(6/\pi^2)^2$ (cf.\ \cite[(1.4)]{DFSU}). Part~(i) of
Theorem~\ref{thm:multi-residue-asymptotic} is the classical Bowen formula
written as the abscissa of a Poincar\'e series; our contribution there is a
proof that avoids the pressure formalism, together with its formalization.
The upper bound in part~(v), and hence the supercritical case of part~(ii),
follow from the Heinemann--Urba\'nski estimate \cite{HU02} (in the form of
\cite[Theorem~2.4]{CLU}) combined with the Lyapunov bound
\cite[Proposition~2.6]{CLU}, and the case $S=\emptyset$ of part~(ii) follows
from Good's theorem by monotonicity. Part~(iii) is a quantitative form of the
Laurent expansion of $\zeta$ at $1$ \cite[(25.2.4)]{DLMF}. A targeted negative
search (arXiv, Crossref and web search; coverage recorded in the verification
archive) found no prior appearance of $C_r$, of $\psi_{\mathrm{Fib}}$, or of the
coefficient $\frac12 d^2C_S(\gamma C_S-C_S')$ as asymptotic dimension constants
of residue-class digit sets; this does not exclude such an appearance.
% Theorem/equation numbers cited for CLU and DFSU above are those of arXiv:1805.11904v1 and
% arXiv:2007.10554v2; the numbering of the published versions has not been verified
% (verification/paper2_novelty_review_20261125.md).
Accordingly, we
regard Theorems~\ref{thm:class} and~\ref{thm:gap} as a variant of known results.
The main new contributions are the complete spectrum discretization and
three-way positive-level splitting (Theorems~\ref{thm:full-spectrum},
\ref{thm:trichotomy}, \ref{thm:general-rec}--\ref{thm:poly-phase-transition},
and~\ref{thm:growing-rec-spectrum}--\ref{thm:growing-rec-dim-phase}), the critical gauge theorems
(Theorem~\ref{thm:gauge} and Corollary~\ref{cor:exact-gauge}), the machine-checked
dimension-zero and subcritical dimension theorems (Theorems~\ref{thm:finite},
\ref{thm:upper}, and~\ref{conj:weight}), the continued-fraction window and
residue theorems (Theorems~\ref{thm:windows}--\ref{thm:residues-general} and the
strict separation in Theorem~\ref{thm:residues-asymptotic}(ii)), and explicit
residue-class instances of known asymptotic methods
(Theorems~\ref{thm:residues-asymptotic}(iii)--(iv)
and~\ref{thm:multi-residue-asymptotic}).

Related metric questions for Cantor series expansions, where the approximation
is measured by the tail $x-\omega_n(x)$ of the partial sums, were studied by
Cheng and Zhang~\cite{CZ}, who determine Hausdorff dimensions of the
corresponding limsup sets; spectra and exact gauges are not considered there.
For the selected sequence $(Q_j)$, Corollary~8 of Bugeaud, Liao and Rams
\cite{BLR} applies directly since $Q_{j+1}/Q_j\to\infty$. Theorems~\ref{thm:poly-phase-transition} and
\ref{thm:growing-rec-dim-phase} are compared with \cite{BLR} in
Remarks~\ref{rem:poly-BLR} and~\ref{rem:growing-BLR}. Corollary~8 gives dimension
one for every constant target $c\in[0,1/2]$ at weight $\alpha=0$ in
Theorem~\ref{conj:weight} (in particular for $\{x:\|Q_jx\|\to0\}$), by taking
the centres equal to $c$ and using $|\|y\|-\|c\||\le\|y-c\|$; integer
translation invariance transfers full dimension between $\R$ and $[0,1]$. For
nonzero weights $0<\alpha<1$, the logarithmic growth condition (1.6) of their
Theorem~7 fails for the linear maps $x\mapsto Q_jx$ because
$\log Q_{j+1}/\log Q_j\to1$, and classical mass-transference methods for
$\limsup$ shrinking-target sets \cite{BV} do not apply to the eventual
(liminf-type) bounds defining $E_c$ and the prescribed-limit sets. We establish
the exact dimension $1-\alpha$ here by an elementary covering upper bound and a
shifted-grid Moran construction.

Regarding Hausdorff gauges at critical levels, Cabrelli, Mendivil, Molter and
Shonkwiler \cite[\S5, Theorem~3]{CMMS} construct a gauge giving positive
\emph{finite} measure to a Cantor set from its gap sequence, whereas Moreira,
Erazo and Angelini \cite{MEA} show that the classical Lagrange level at value
$3$ has either zero or non-$\sigma$-finite measure for every dimension function.
Here the recurrence specifies the gauge $h_k(1/Q_n)=(2k+1)^{-n}$, and the entire
critical level $\{\Lam=k/4\}$ on $\R/\Z$ has \emph{infinite, $\sigma$-finite}
$\mathcal H^{h_k}$ measure, proved by combining a countable decomposition by
initial recurrence data with translation invariance under the dense subgroup
$L$.

Relative to earlier literature, our setting has four distinguishing features:
(i) a \emph{subsequence} of convergent denominators of $\ee$ (every third),
(ii) the critical \emph{weight} $j$, (iii) the lattice
$L=\Z+\Z\frac{\ee+1}{2}$ (note $(\ee+1)/2$, not $\ee$), and (iv) the
\emph{sharp} quantum $1/4$. The weight $j$ is \emph{critical} for the covering
argument: it makes the per-parent branching number uniformly bounded, yielding
dimension $0$ for \emph{all} finite level sets (Theorem~\ref{thm:finite}),
whereas without the weight the branching grows linearly and the unweighted
zero-set $\{x:\|Q_jx\|\to0\}$ has full dimension $1$, while the weighted
zero-set $\{x:j\|Q_jx\|\to0\}$ equals the countable group $L$.

\section{Classification and initial spectral gap}\label{sec:class}

The two results of this section were first proved, and Lean-kernel-checked at
Layer~A, in the Phase-1 project \cite{Phase1}. Here Layer~A means that the
audited declaration in Lean~4 depends only on Lean's three standard base axioms
\texttt{propext}, \texttt{Classical.choice} and \texttt{Quot.sound}. For
completeness we give short self-contained proofs. They use three facts about the
selected convergents $P_j/Q_j$:
\begin{enumerate}[label=(\arabic*),leftmargin=2em]
\item $P_j$ and $Q_j$ are odd, and $P_j$ satisfies the same recurrence
\eqref{eq:rec} as $Q_j$.
\item $P_jQ_{j+1}-P_{j+1}Q_j=2(-1)^{j+1}$ for all $j\ge1$.
\item The approximation error $\delta_j:=Q_j\ee-P_j$ satisfies $\delta_j\to0$
and $|\delta_j|<1/(2jQ_j)$ for all $j\ge1$.
\end{enumerate}
The last estimate is not in mathlib v4.35, and earlier drafts of the Phase-1
formalization assumed it as an axiom. Both $\delta_j\to0$ and
$|\delta_j|<1/(2jQ_j)$ are now kernel-checked Lean theorems
(\texttt{deltaSel\_tendsto\_zero} and \texttt{abs\_deltaSel\_lt}, proved via
\texttt{DeltaSign.lean}), derived from the Hermite--Pad\'e integral
representation
\[
(-1)^n n!\,\delta_n=\int_0^1 \ee^t t^n(1-t)^n\,dt>0.
\]

\paragraph{Digits.} For $x\in\R$ and $j\ge1$ let $M_j=\mathrm{round}(Q_jx)\in\Z$
be a nearest integer to $Q_jx$, put $u_j=Q_jx-M_j\in[-1/2,1/2]$ (so
$|u_j|=\|Q_jx\|\le1/2$), and for $j\ge2$ define the integer recurrence defect
(digit)
\begin{equation}\label{eq:defect-id}
 d_j:=M_{j+1}-A_jM_j-M_{j-1}=A_ju_j+u_{j-1}-u_{j+1}\in\Z,\qquad A_j:=4j+2,
\end{equation}
where the second equality follows immediately from \eqref{eq:rec}.

\begin{lemmaA}
If $d_j=0$ for all $j\ge N$, then $x\in L$.
\end{lemmaA}

\begin{proof}
Put $v_j=(P_j+Q_j)/2$, an integer since $P_j$ and $Q_j$ are odd. The
vectors $(v_N,v_{N+1})$ and $(Q_N,Q_{N+1})$ have determinant
$\tfrac12(P_NQ_{N+1}-P_{N+1}Q_N)=\pm1$, so they form a basis of $\Z^2$.
Choose integers $a,b$ with $M_N=av_N+bQ_N$ and $M_{N+1}=av_{N+1}+bQ_{N+1}$.
The sequences $(M_j)$, $(v_j)$ and $(Q_j)$ all satisfy \eqref{eq:rec}
from index $N$ on, hence $M_j=av_j+bQ_j$ for all $j\ge N$. With
$y:=x-b-a\frac{\ee+1}{2}$ this gives
\[
 u_j=Q_jx-M_j=Q_jy+\tfrac a2\,\delta_j .
\]
Since $|u_j|\le1/2$, $\delta_j\to0$ and $Q_j\to\infty$, we get $y=0$, that
is, $x=b+a\frac{\ee+1}{2}\in L$.
\end{proof}

\begin{theorem}[Classification, {\cite[Thm.~2]{Phase1}}]\label{thm:class}
For $x\in\R$, $\quad j\|Q_jx\|\to0\iff x\in L.$
\end{theorem}

\begin{proof}
If $x=m+n\frac{\ee+1}{2}$ with $m,n\in\Z$, then
$Q_jx=Q_jm+nv_j+\tfrac n2\delta_j$ with $v_j$ as above, so
$\|Q_jx\|\le|n|\,|\delta_j|/2<|n|/(4jQ_j)$ and $j\|Q_jx\|\to0$.
Conversely, if $j|u_j|\to0$, then
$|d_j|\le A_j|u_j|+|u_{j-1}|+|u_{j+1}|\to0$. The $d_j$ are integers, so
$d_j=0$ eventually, and $x\in L$ by Lemma~A.
\end{proof}

\begin{theorem}[Gap, {\cite[Cor.~3]{Phase1}}]\label{thm:gap}
If $x\notin L$ and $c<1/4$, then it is \emph{not} the case that
$j\|Q_jx\|\le c$ eventually. Equivalently $\Lam(x)\ge1/4$ for $x\notin L$.
\end{theorem}

\begin{proof}
Suppose $j|u_j|\le c<1/4$ for all $j\ge J$. For $j\ge J+1$,
\[
 |d_j|\le\frac{(4j+2)c}{j}+\frac{c}{j-1}+\frac{c}{j+1}
 \le4c+\frac{4c}{j-1},
\]
which is less than $1$ for all large $j$. Hence $d_j=0$ eventually, and
$x\in L$ by Lemma~A.
\end{proof}

\paragraph{Lean status.} Theorems~\ref{thm:class} and~\ref{thm:gap} are
kernel-checked at Layer~A in the module
\texttt{EpiFormal/\allowbreak Classification.lean} as
\texttt{theorem2\_\allowbreak classification} and
\texttt{corollary3\_\allowbreak gap}. Their axiom audit reports
only \texttt{propext}, \texttt{Classical.choice} and \texttt{Quot.sound}.

\section{Complete spectrum discretization and step-function collapse}\label{sec:spectrum}

\begin{corollary}[Initial spectral gap, Layer~A in Lean]\label{thm:spectral}
For every $x\in\R$, either $\Lam(x)=0$ or $\Lam(x)\ge1/4$. Thus
$\Lam(\R)\subseteq\{0\}\cup[1/4,\infty]$, and $\Lam(x)=0\iff x\in L$.
\end{corollary}

This is an immediate consequence of Theorems~\ref{thm:class}
and~\ref{thm:gap}; we state it separately because it is the form used
below and the form directly machine-checked in \texttt{SpectralGap.lean}.

\begin{proof}
If $x\in L$, Theorem~\ref{thm:class} gives $j\|Q_jx\|\to0$, hence $\Lam(x)=0$.
If $x\notin L$, Theorem~\ref{thm:gap} gives $\Lam(x)\ge1/4$. The dichotomy and
the characterization of the zero value follow.
\end{proof}

\begin{remark}\label{rem:linear-indep}
In particular $\Lam(\ee)=0$, since $\ee=2\cdot\frac{\ee+1}{2}-1\in L$; $\ee$ is
\emph{not} a generic point for this exponent. For $x=B\pi$ ($B=1,\dots,10$),
$\sqrt2$, the golden ratio and random reals, the values $j\|Q_jx\|$ for
$j\le60$ are large (Section~\ref{sec:num}), but such finite data do not
determine $\Lam(x)$. By Corollary~\ref{thm:spectral}, the statement
``$\Lam(B\pi)>0$ for every integer $B\ge1$'' is equivalent to
$B\pi\notin L$ for all $B\ge1$, i.e.\ to the linear independence of
$1,\ee,\pi$ over $\Q$, which is open. Nothing in this article decides it.
\end{remark}

We now show that the initial gap $(0,1/4)$ is merely the bottom rung of an
infinite ladder of equidistant spectral gaps: the entire finite spectrum of
$\Lam$ is quantized in multiples of $1/4$.

\begin{theorem}[Complete discretization of the spectrum, Layer~A in Lean]\label{thm:full-spectrum}
For every $x\in\R$ with $\Lam(x)<\infty$, one has $u_j=O(1/j)\to0$ and
\begin{equation}\label{eq:digit-asymp}
\lim_{j\to\infty}\bigl(d_j-4ju_j\bigr)=0,\qquad
\lim_{j\to\infty}\bigl(|d_j|-4j\|Q_jx\|\bigr)=0.
\end{equation}
Consequently $4\Lam(x)=\limsup_{j\to\infty}|d_j|\in\Z_{\ge0}$ whenever
$\Lam(x)<\infty$, and the full spectrum of $\Lam$ is given by the exact identity
\begin{equation}\label{eq:full-spectrum}
\Lam(\R)=\frac14\Z_{\ge0}\cup\{\infty\}
=\left\{0,\frac14,\frac12,\frac34,1,\dots\right\}\cup\{\infty\}.
\end{equation}
In particular, for every integer $k\ge0$, the open interval
$\bigl(\frac{k}{4},\frac{k+1}{4}\bigr)$ is an unconditional empty spectral gap
for $\Lam$.
(Lean, module \texttt{FullSpectrumDiscrete}:
\texttt{euler\_\allowbreak signed\_\allowbreak defect\_\allowbreak approx},
\texttt{euler\_\allowbreak abs\_\allowbreak defect\_\allowbreak approx}, and
\texttt{spectral\_\allowbreak gap\_\allowbreak open\_\allowbreak interval}.)
\end{theorem}

\begin{proof}
If $\Lam(x)<\infty$, then $\|Q_jx\|=|u_j|\le(\Lam(x)+1)/j$ for all large $j$,
so $u_j\to0$. Rearranging \eqref{eq:defect-id} gives
$d_j-4ju_j=2u_j+u_{j-1}-u_{j+1}\to0$, and the reverse triangle inequality
$\bigl||d_j|-4j|u_j|\bigr|\le|d_j-4ju_j|$ yields \eqref{eq:digit-asymp}.
Since $|d_j|\in\Z_{\ge0}$, the sequence $(|d_j|)$ can only have a nonnegative
integer or $\infty$ as its limit superior. By \eqref{eq:digit-asymp},
$4\Lam(x)=\limsup_{j\to\infty}4j\|Q_jx\|=\limsup_{j\to\infty}|d_j|\in\Z_{\ge0}\cup\{\infty\}$,
which excludes every open interval $(k/4,(k+1)/4)$ from $\Lam(\R)$.
Conversely, $\Lam(0)=0$, $\Lam=\infty$ holds on a set of full Lebesgue measure
(Section~\ref{sec:finite-levels}), and for each integer $k\ge1$ choosing a
sufficiently large starting index $N$ and constant digits $d_j\equiv k$ for
$j\ge N$ (as in the proof of Theorem~\ref{thm:gauge} below) produces a point
$x_k\in\R$ with $|u_j|<1/8$ for $j\ge N$, $M_j=\mathrm{round}(Q_jx_k)$, and
$4j\|Q_jx_k\|\to k$, so $\Lam(x_k)=k/4$. Thus every value in
$\frac14\Z_{\ge0}\cup\{\infty\}$ is attained.
\end{proof}

For $c\ge0$ write $E_c:=\{x\in\R:\Lam(x)\le c\}$.

\begin{corollary}[Sub-threshold level sets and gap constancy, Layer~A in Lean]\label{thm:sub}
For every integer $k\ge0$ and every $c\in\bigl[\frac{k}{4},\frac{k+1}{4}\bigr)$,
one has $E_c=E_{k/4}=\{x\in\R:u_j(x)\to0\text{ and }\limsup_{j\to\infty}|d_j(x)|\le k\}$.
In particular, if $c<1/4$ ($k=0$), then $E_c=E_0=L$, so $E_c$ is countable and
$\dimH E_c=0$.
(Lean, modules \texttt{FullSpectrumDiscrete} and \texttt{SpectralGap}:
\texttt{eventually\_\allowbreak abs\_\allowbreak defect\_\allowbreak le\_\allowbreak of\_\allowbreak apprExponent\_\allowbreak le},
\texttt{levelSet\_\allowbreak eq\_\allowbreak on\_\allowbreak gap},
\texttt{levelSet\_\allowbreak below\_\allowbreak eq\_\allowbreak lattice}, and
\texttt{dimH\_\allowbreak levelSet\_\allowbreak below}.)
\end{corollary}

\begin{proof}
Clearly $E_{k/4}\subseteq E_c$ for $c\ge k/4$. Conversely, if $x\in E_c$ with
$c<(k+1)/4$, then $\Lam(x)\le c<(k+1)/4$ cannot lie in the empty spectral gap
$(k/4,(k+1)/4)$ by Theorem~\ref{thm:full-spectrum}, so $\Lam(x)\le k/4$ and
eventually $|d_j(x)|\le k$. For $k=0$, $\Lam(x)=0\iff x\in L$ by
Corollary~\ref{thm:spectral}, and $L$ is a countable image of $\Z^2$, hence has
Hausdorff dimension $0$ \cite{Falconer}.
\end{proof}

\paragraph{Lean status.} Corollaries~\ref{thm:spectral} and~\ref{thm:sub} (for
$c<1/4$), together with the countability of $L$ and $\dimH E_c=0$ for $c<1/4$,
are kernel-checked at Layer~A in \texttt{EpiFormal/\allowbreak SpectralGap.lean}
(five declarations: \texttt{apprExponent\_eq\_zero\_iff},
\texttt{spectral\_gap}, \texttt{levelSet\_below\_eq\_lattice},
\texttt{latticeL\_countable}, and \texttt{dimH\_levelSet\_below}; no
\texttt{sorryAx}). The full spectrum discretization, open-interval spectral-gap
exclusion $(k/4,(k+1)/4)$, and the step-function constancy $E_c=E_{k/4}$ across
each gap are kernel-checked at Layer~A in
\texttt{EpiFormal/\allowbreak FullSpectrumDiscrete.lean} (15 declarations, base
axioms only, no \texttt{sorryAx}). All of these modules---along with
\texttt{Gauge\allowbreak Liminf\allowbreak Complete.lean},
\texttt{Finite\allowbreak State\allowbreak Gauge.lean},
\texttt{Finite\allowbreak State\allowbreak Positive.lean},
\texttt{Finite\allowbreak State\allowbreak Complete.lean},
\texttt{Window\allowbreak Short\allowbreak AeRealize.lean},
\texttt{Residue\allowbreak ModGeneral.lean},
\texttt{Transfer\allowbreak Operator\allowbreak Dim.lean}, and
\texttt{Trichotomy\allowbreak And\allowbreak Poly\allowbreak Extension.lean}
(which completes Theorem~\ref{thm:trichotomy}(iii) and Section~\ref{sec:poly-rec})---are
audited together in \texttt{Audit.lean} ($872$ declarations across $131$ files,
$9176$ build jobs) with zero \texttt{sorryAx} and only the three Lean base axioms.

\section{All finite level sets have dimension zero, and the complementary infinite level}\label{sec:finite-levels}

The sub-threshold result ($c<1/4$) uses the sharp classification. For
$c\ge1/4$ the level set $E_c$ is uncountable (Theorem~\ref{thm:gauge}), but we
show it still has Hausdorff dimension $0$. The entire covering argument ---
factorial growth of $Q_j$, uniform branching bound, surviving-center recursion,
periodicity reduction, Hausdorff-content estimate, and the countable-union
assembly --- is \emph{machine-checked at Layer~A}
(\texttt{EpiFormal/DimHCore.lean} + \texttt{EpiFormal/DimHFinite.lean}, base
axioms only, no \texttt{sorryAx}, and \emph{no} continued-fraction axioms).
The elementary proof below mirrors the formalization.

\begin{lemma}[Factorial growth, Layer~A in Lean]\label{lem:factorial}
For $j\ge1$, $Q_j\ge j!$. (Lean: \texttt{DimHCore.\allowbreak Qsel\_ge\_factorial}.)
\end{lemma}

\begin{proof}
We proceed by induction on $j\ge1$. For $j=1$, $Q_1=1=1!$. For $j=2$,
$Q_2=7\ge2=2!$. For $j\ge2$, every term $Q_i$ is positive, so the recurrence
\eqref{eq:rec} gives
\[
Q_{j+1}=(4j+2)Q_j+Q_{j-1}>(j+1)Q_j\ge(j+1)j!=(j+1)!,
\]
completing the induction.
\end{proof}

\begin{lemma}[Uniform branching bound, Layer~A in Lean]\label{lem:branch}
Fix $t\ge0$, $j\ge1$ and a parent center $c\in\Z$. The number of child centers
$k'\in\{0,1,\dots,Q_{j+1}\}$ with
\[
\left|\frac{c}{Q_j}-\frac{k'}{Q_{j+1}}\right|\le \frac{t}{jQ_j}+\frac{t}{(j+1)Q_{j+1}}
\]
is at most $B:=\lfloor16t\rfloor+1$, \emph{uniformly in $j$ and $c$}.
(Lean: \texttt{DimHCore.\allowbreak fiber\_card\_le}.)
\end{lemma}

\begin{proof}
Put $R_j:=t/(jQ_j)+t/((j+1)Q_{j+1})$. Multiplying the distance inequality
$|c/Q_j-k'/Q_{j+1}|\le R_j$ by $Q_{j+1}>0$ shows that every qualifying integer
$k'$ lies in the closed real interval $[A,A+W]$ with left endpoint
$A:=Q_{j+1}(c/Q_j-R_j)$ and width
\[
W:=2R_jQ_{j+1}=2t\left(\frac{Q_{j+1}}{jQ_j}+\frac{1}{j+1}\right).
\]
Since $Q_0=Q_1=1$ and $Q_{i-1}\le Q_i$ for all $i\ge1$, the recurrence
\eqref{eq:rec} implies $Q_{j+1}=(4j+2)Q_j+Q_{j-1}\le(4j+3)Q_j\le7jQ_j$ for all
$j\ge1$, while $1/(j+1)\le1$. Hence
\[
W\le 2t(7+1)=16t.
\]
For any qualifying $k'$, we have $0\le k'-A\le W\le16t$, so the map
$k'\mapsto\lfloor k'-A\rfloor$ takes values in $\{0,1,\dots,\lfloor16t\rfloor\}$
and is injective on integers (if $k'_1<k'_2$, then $k'_2-A\ge(k'_1-A)+1$, so
$\lfloor k'_2-A\rfloor\ge\lfloor k'_1-A\rfloor+1$). Therefore the number of
qualifying integers $k'$ is at most $\lfloor16t\rfloor+1=B$.
\end{proof}

\begin{theorem}[Finite level sets, Layer~A in Lean]\label{thm:finite}
For every $c\in\R$, $\dimH E_c=0$. Consequently
$F:=\{x\in\R:\Lam(x)<\infty\}=\bigcup_{n\ge1}E_n$ has $\dimH F=0$.
(Lean: \texttt{DimHFinite.\allowbreak dimH\_levelSet\_finite}; \texttt{\#print axioms}
returns \texttt{[propext, Classical.choice, Quot.sound]} only.)
\end{theorem}

\begin{proof}
Fix $c\in\R$ and choose a rational or real threshold $t>\max(c,0)$ (in Lean,
$t=\max(c,0)+1>0$). First note that for every $m\in\Z$ and $j\ge1$, since
$Q_j\in\Z$, we have $\|Q_j(x+m)\|=\|Q_jx+Q_jm\|=\|Q_jx\|$, and hence
$\Lam(x+m)=\Lam(x)$. Thus $E_c$ is invariant under integer translations and
decomposes as the countable union
\[
E_c=\bigcup_{m\in\Z}\bigl(m+(E_c\cap[0,1])\bigr).
\]
Because integer translations are isometries and Hausdorff dimension is stable
under countable unions ($\dimH(\bigcup_m X_m)=\sup_m\dimH X_m$), it suffices to
prove $\dimH(E_c\cap[0,1])=0$.

For each $j\ge1$, define $r_j:=t/(jQ_j)$ and
$A_j(t):=\{x\in\R:\|Q_jx\|\le t/j\}$. If $x\in E_c\cap[0,1]$, then
$\limsup_{j\to\infty}j\|Q_jx\|=\Lam(x)\le c<t$, which by definition of the limit
superior means that $j\|Q_jx\|\le t$ \emph{eventually} (a liminf-type tail
condition, \emph{not} a limsup set). Therefore there exists an integer $J\ge1$
such that $x\in A_j(t)$ for all $j\ge J$. Consequently,
\[
E_c\cap[0,1]\subseteq\bigcup_{J=1}^\infty C_J,\qquad
C_J:=[0,1]\cap\bigcap_{j=J}^\infty A_j(t).
\]
By countable stability again, it suffices to fix $J\ge1$ and prove that
$\dimH C_J=0$, i.e.\ $\mathcal H^s(C_J)=0$ for every $s>0$.

For $j\ge J$ and $k\in\Z$, let $I_j(k):=[k/Q_j-r_j,k/Q_j+r_j]$ be the closed
interval of center $k/Q_j$ and radius $r_j$ (so $\operatorname{diam}I_j(k)=2r_j$).
For any $x\in C_J$ and $j\ge J$, let $k_j(x):=\mathrm{round}(Q_jx)\in\Z$ be a
nearest integer to $Q_jx$. Since $x\in[0,1]$ and $Q_j\ge1$, we have
$0\le Q_jx\le Q_j$, so $k_j(x)\in\{0,1,\dots,Q_j\}$; moreover, $x\in A_j(t)$
gives $|Q_jx-k_j(x)|=\|Q_jx\|\le t/j$, and dividing by $Q_j$ yields
$|x-k_j(x)/Q_j|\le r_j$, that is, $x\in I_j(k_j(x))$. Furthermore, at two
consecutive levels $j$ and $j+1$, the same point $x\in C_J$ belongs to both
$I_j(k_j(x))$ and $I_{j+1}(k_{j+1}(x))$, so the triangle inequality gives the
parent--child center compatibility condition
\begin{equation}\label{eq:parent-child-tri}
\left|\frac{k_j(x)}{Q_j}-\frac{k_{j+1}(x)}{Q_{j+1}}\right|
\le\left|\frac{k_j(x)}{Q_j}-x\right|+\left|x-\frac{k_{j+1}(x)}{Q_{j+1}}\right|
\le r_j+r_{j+1}=\frac{t}{jQ_j}+\frac{t}{(j+1)Q_{j+1}}.
\end{equation}

Define the surviving center sets $S_n\subseteq\{0,1,\dots,Q_{J+n}\}$ for
$n\ge0$ inductively by setting $S_0:=\{0,1,\dots,Q_J\}$ and, for $n\ge0$,
letting $S_{n+1}$ be the set of all $k'\in\{0,1,\dots,Q_{J+n+1}\}$ for which
there exists $k\in S_n$ satisfying
$|k/Q_{J+n}-k'/Q_{J+n+1}|\le r_{J+n}+r_{J+n+1}$. By induction on $n\ge0$ using
\eqref{eq:parent-child-tri}, every $x\in C_J$ satisfies $k_{J+n}(x)\in S_n$ for
all $n\ge0$, and hence
\[
C_J\subseteq\bigcup_{k\in S_n} I_{J+n}(k)\qquad\text{for all }n\ge0.
\]
By Lemma~\ref{lem:branch}, for each parent $k\in S_n$ the number of children
$k'\in S_{n+1}$ is at most $B=\lfloor16t\rfloor+1$. Since $|S_0|=Q_J+1$,
induction gives $|S_n|\le(Q_J+1)B^n$ for all $n\ge0$. Each interval
$I_{J+n}(k)$ has diameter
\[
\operatorname{diam}I_{J+n}(k)=2r_{J+n}=\frac{2t}{(J+n)Q_{J+n}}
\le\frac{2t}{(J+n)!}\longrightarrow0\qquad(n\to\infty),
\]
where we used $Q_{J+n}\ge(J+n)!$ from Lemma~\ref{lem:factorial} and $J+n\ge1$.

Now fix any $s>0$ and any gauge scale $\delta>0$. For all sufficiently large
$n$, the intervals $\{I_{J+n}(k):k\in S_n\}$ form a finite $\delta$-cover of
$C_J$, and by Lemma~\ref{lem:factorial} ($Q_{J+n}\ge(J+n)!$) their $s$-dimensional
Hausdorff covering cost satisfies
\[
\sum_{k\in S_n}\bigl(\operatorname{diam}I_{J+n}(k)\bigr)^s
\le |S_n|\left(\frac{2t}{(J+n)Q_{J+n}}\right)^s
\le(Q_J+1)(2t)^s B^{-J}\left(\frac{(B^{1/s})^{J+n}}{(J+n)!}\right)^s.
\]
Using the elementary factorial lower bound $m!\ge(m/\ee)^m$ for $m=J+n\ge1$
(or the convergent ratio $(B^{1/s})^{m+1}/(m+1)!\div((B^{1/s})^m/m!)=B^{1/s}/(m+1)<1/2$
for $m+1>2B^{1/s}$, formalized in mathlib as
\texttt{FloorSemiring.tendsto\_pow\_div\_factorial\_atTop}), we have
\[
0\le\frac{(B^{1/s})^m}{m!}\le\left(\frac{\ee B^{1/s}}{m}\right)^m\longrightarrow0
\qquad(m=J+n\to\infty),
\]
and raising to the positive power $s>0$ shows that the covering cost tends to
$0$ as $n\to\infty$. Thus $\mathcal H_\delta^s(C_J)=0$ for every $\delta>0$,
so $\mathcal H^s(C_J)=\lim_{\delta\downarrow0}\mathcal H_\delta^s(C_J)=0$ for
every $s>0$, which gives $\dimH C_J=0$. Taking the countable union over
$J\ge1$ and $m\in\Z$ yields $\dimH E_c=0$. Finally,
$F=\{x\in\R:\Lam(x)<\infty\}=\bigcup_{n=1}^\infty E_n$ is a countable union of
sets of Hausdorff dimension $0$, so $\dimH F=0$.
\end{proof}

\begin{corollary}[The complementary infinite level, Layer~A in Lean]\label{cor:infinite-level}
The set $F=\{x\in\R:\Lam(x)<\infty\}$ has Lebesgue measure zero, and its
complement $F_\infty:=\{x\in\R:\Lam(x)=\infty\}=\R\setminus F$ has full
Lebesgue measure on $\R$. Moreover, for every nondegenerate interval
$[a,b]\subset\R$ ($a<b$), one has
$\operatorname{Leb}(F_\infty\cap[a,b])=b-a>0$ and
\[
\dimH(F_\infty\cap[a,b])=\dimH F_\infty=1.
\]
(Lean, module \texttt{EpiFormal/\allowbreak GaugeExact.lean}: Layer~A, base axioms only.)
\end{corollary}

\begin{proof}
That $\operatorname{Leb}(F)=0$ follows in two independent ways. First, by
Theorem~\ref{thm:finite}, $\dimH F=0<1$, which implies $\mathcal H^1(F)=0$;
since one-dimensional Hausdorff measure $\mathcal H^1$ on $\R$ coincides with
outer Lebesgue measure $\operatorname{Leb}$, we obtain $\operatorname{Leb}(F)=0$.

Second, one can see $\operatorname{Leb}(F)=0$ directly by an elementary
measure-preserving calculation without invoking Hausdorff dimension. For an
integer $m\ge1$ and $j>2m$, the circle endomorphism $T_{Q_j}:\R/\Z\to\R/\Z$,
$x\mapsto Q_jx\bmod1$, preserves normalized Lebesgue measure on $\R/\Z$ (each
arc of length $\ell$ has $Q_j$ disjoint preimage arcs of length $\ell/Q_j$).
The set $\bar A_j(m):=\{x\in\R/\Z:\|Q_jx\|\le m/j\}=T_{Q_j}^{-1}([-m/j,m/j]\bmod1)$
therefore has circle Lebesgue measure $2m/j$. For any fixed starting index
$J\ge1$ and every $j\ge\max(J,2m+1)$, monotonic inclusion gives
\[
\operatorname{Leb}_{\R/\Z}\left(\bigcap_{i=J}^\infty \bar A_i(m)\right)
\le\operatorname{Leb}_{\R/\Z}\bigl(\bar A_j(m)\bigr)=\frac{2m}{j}\longrightarrow0
\qquad(j\to\infty),
\]
so $\operatorname{Leb}_{\R/\Z}(\bigcap_{i\ge J}\bar A_i(m))=0$. Every point
$x\in\R/\Z$ with $\Lam(x)<\infty$ belongs to $\bigcap_{i\ge J}\bar A_i(m)$ for
some integers $m>\Lam(x)$ and $J\ge1$. Countable subadditivity over $m\ge1$ and
$J\ge1$ yields $\operatorname{Leb}_{\R/\Z}(\{\Lam<\infty\})=0$ on the circle,
and countable union over integer translates $[k,k+1]$ ($k\in\Z$) yields
$\operatorname{Leb}(F)=0$ on $\R$.

Now let $[a,b]\subset\R$ be any nondegenerate closed interval ($a<b$). Since
$[a,b]=(F_\infty\cap[a,b])\cup(F\cap[a,b])$ and $\operatorname{Leb}(F\cap[a,b])\le\operatorname{Leb}(F)=0$,
additivity of Lebesgue measure gives
\[
\mathcal H^1(F_\infty\cap[a,b])=\operatorname{Leb}(F_\infty\cap[a,b])
=\operatorname{Leb}([a,b])-\operatorname{Leb}(F\cap[a,b])=b-a>0.
\]
By the definition of Hausdorff dimension, $\mathcal H^1(F_\infty\cap[a,b])>0$
implies $\dimH(F_\infty\cap[a,b])\ge1$. On the other hand, any subset of $\R$
has Hausdorff dimension at most $\dimH\R=1$, so
\[
1\le\dimH(F_\infty\cap[a,b])\le\dimH F_\infty\le\dimH\R=1.
\]
Hence $\dimH(F_\infty\cap[a,b])=\dimH F_\infty=1$. Together with
Theorem~\ref{thm:finite}, this establishes the sharp metric dichotomy: $\Lam$
is infinite Lebesgue-almost everywhere (with local Hausdorff dimension~$1$ on
every interval), while the set where $\Lam$ is finite has Hausdorff
dimension~$0$.
\end{proof}

\paragraph{Lean status of the assembly.} All steps of Theorem~\ref{thm:finite}
are kernel-checked at Layer~A. The surviving centers are encoded as a
\texttt{Finset} recursion (\texttt{survSet}), not as paths:
$S_0=\{0,\dots,Q_J\}$ and $S_{n+1}=\bigcup_{k\in S_n}\mathrm{branchSet}(k)$,
with $|S_n|\le(Q_J+1)B^n$ by \texttt{Finset.card\_biUnion\_le\_card\_mul} and
Lemma~\ref{lem:branch} (\texttt{card\_survSet\_le}). Membership of the center
sequence $\mathrm{round}(Q_{J+n}x)$ in $S_n$ is proved by induction via the
triangle inequality (\texttt{centerOf\_mem\_survSet}), giving the cover
(\texttt{tailSet\_cover}). The content bound
$|S_n|(2t/((J+n)Q_{J+n}))^s\to0$ uses $Q_j\ge j!$ and
\texttt{FloorSemiring.\allowbreak tendsto\_pow\_div\_factorial\_atTop} with
$c=B^{1/s}$ (\texttt{tendsto\_contentBound}); no asymptotic notation is used.
Feeding this into \texttt{hausdorffMeasure\_le\_liminf\_sum} yields
$\mu H^s=0$ for all $s>0$ (\texttt{hausdorffMeasure\_tailSet\_zero}), hence
$\dimH=0$ (\texttt{dimH\_tailSet\_zero}). Periodicity
$\|Q_j(x+m)\|=\|Q_jx\|$ (\texttt{apprSeq\_add\_intCast}) reduces general $x$ to
$[0,1]$; integer shifts are $1$-Lipschitz and countable unions preserve
dimension $0$ (\texttt{dimH\_iUnion}). The sub-threshold case $c<1/4$ is
\emph{not} used: a monotonicity reduction to $c'=\max(c,0)$ lets the covering
argument handle all $c$ uniformly, so the theorem needs no continued-fraction
axioms.

\section{Subcritical weights and exact Hausdorff dimensions}\label{sec:weights}

\begin{remark}\label{rem:unweighted}
The inequality $\|Q_jx\|\le j\|Q_jx\|$ gives the inclusion from the
weighted zero-set $\{x:j\|Q_jx\|\to0\}=L$ into the unweighted zero-set
$\{x:\|Q_jx\|\to0\}$. To see strictness explicitly, the point $x_0$ constructed
in \texttt{QuarterPointFinal.lean} (with constant digits $d_j\equiv1$) satisfies
$j\|Q_jx_0\|\to1/4$ (\texttt{QuarterPointFinal.\allowbreak lambda\_x0\_eq\_quarter}).
Its unweighted error $\|Q_jx_0\|$ tends to zero, while $x_0\notin L$ by the
weighted classification (Theorem~\ref{thm:class}). The combined witness is
kernel-checked at Layer~A in \texttt{EpiFormal/\allowbreak UnweightedCorrection.lean}
(\texttt{UnweightedCorrection.\allowbreak witness}) using only the three Lean
base axioms.
\end{remark}

Recall that for $0\le\alpha\le1$ we write
$\Lam_\alpha(x)=\limsup_{j\to\infty}j^\alpha\|Q_jx\|$, so that $\Lam_1=\Lam$.

\begin{theorem}[Weighted upper bound, Layer~A in Lean]\label{thm:upper}
For $0\le\alpha\le1$ and every finite $c\in\R$,
$\dimH\{x\in\R:\Lam_\alpha(x)\le c\}\le1-\alpha$.
(Lean: \texttt{WeightedFinal.\allowbreak dimH\_levelSet\ensuremath{\alpha}\_le};
\texttt{\#print axioms} returns \texttt{[propext, Classical.choice, Quot.sound]} only.)
\end{theorem}

\begin{proof}
Fix $0\le\alpha\le1$ and $c\in\R$, and choose $t>\max(c,0)$ (in Lean,
$t=\max(c,0)+1>0$). Put $E_{\alpha,c}:=\{x\in\R:\Lam_\alpha(x)\le c\}$. Since
$Q_j\in\Z$, we have $\|Q_j(x+m)\|=\|Q_jx\|$ for every $m\in\Z$ and $j\ge1$,
hence $\Lam_\alpha(x+m)=\Lam_\alpha(x)$. Thus
\[
E_{\alpha,c}=\bigcup_{m\in\Z}\bigl(m+(E_{\alpha,c}\cap[0,1])\bigr).
\]
By isometry invariance and countable stability of Hausdorff dimension, it
suffices to show that $\dimH(E_{\alpha,c}\cap[0,1])\le1-\alpha$.

For $j\ge1$, set $r_{j,\alpha}:=t/(j^\alpha Q_j)$ and
$A_{j,\alpha}(t):=\{x\in\R:\|Q_jx\|\le t/j^\alpha\}$. Any $x\in E_{\alpha,c}\cap[0,1]$
satisfies $\limsup_{j\to\infty}j^\alpha\|Q_jx\|\le c<t$, so $j^\alpha\|Q_jx\|\le t$
eventually; hence
\[
E_{\alpha,c}\cap[0,1]\subseteq\bigcup_{J=1}^\infty C_{J,\alpha},\qquad
C_{J,\alpha}:=[0,1]\cap\bigcap_{j=J}^\infty A_{j,\alpha}(t).
\]
Fix $J\ge1$. For each $j\ge J$ and $k\in\Z$, let
$I_{j,\alpha}(k):=[k/Q_j-r_{j,\alpha},k/Q_j+r_{j,\alpha}]$ be the closed
interval of center $k/Q_j$ and radius $r_{j,\alpha}$. For any $x\in C_{J,\alpha}$
and $j\ge J$, the nearest integer $k_j(x):=\mathrm{round}(Q_jx)$ belongs to
$\{0,1,\dots,Q_j\}$ and satisfies $|x-k_j(x)/Q_j|=\|Q_jx\|/Q_j\le r_{j,\alpha}$,
so $x\in I_{j,\alpha}(k_j(x))$. At consecutive levels $j$ and $j+1$, the
triangle inequality gives
\begin{equation}\label{eq:weighted-parent-child}
\left|\frac{k_j(x)}{Q_j}-\frac{k_{j+1}(x)}{Q_{j+1}}\right|
\le\left|\frac{k_j(x)}{Q_j}-x\right|+\left|x-\frac{k_{j+1}(x)}{Q_{j+1}}\right|
\le r_{j,\alpha}+r_{j+1,\alpha}=\frac{t}{j^\alpha Q_j}+\frac{t}{(j+1)^\alpha Q_{j+1}}.
\end{equation}
Since $0\le\alpha\le1$ and $j\ge1$, we have $1-\alpha\ge0$ and
$j=j^\alpha j^{1-\alpha}\le j^\alpha(j+1)^{1-\alpha}$, so
$1/j^\alpha\le(j+1)^{1-\alpha}/j$, while
$1/(j+1)^\alpha=(j+1)^{1-\alpha}/(j+1)$. Putting $t'_j:=t(j+1)^{1-\alpha}\ge0$,
we obtain the pointwise radius comparison (\texttt{WeightedFiber.wrad\_le\_rad})
\[
r_{j,\alpha}+r_{j+1,\alpha}
\le\frac{t'_j}{jQ_j}+\frac{t'_j}{(j+1)Q_{j+1}}.
\]
Applying Lemma~\ref{lem:branch} with threshold $t'_j$ shows that for each fixed
parent center $k\in\{0,1,\dots,Q_j\}$, the number of child centers
$k'\in\{0,1,\dots,Q_{j+1}\}$ satisfying \eqref{eq:weighted-parent-child} is at
most
\[
\lfloor16t'_j\rfloor+1\le 16t(j+1)^{1-\alpha}+1\le K(j+1)^{1-\alpha},
\qquad K:=16t+1,
\]
since $(j+1)^{1-\alpha}\ge1$.

Define the surviving center sets $S_{n,\alpha}\subseteq\{0,1,\dots,Q_{J+n}\}$
($n\ge0$) inductively by $S_{0,\alpha}:=\{0,1,\dots,Q_J\}$ and letting
$S_{n+1,\alpha}$ consist of all $k'\in\{0,1,\dots,Q_{J+n+1}\}$ within
$r_{J+n,\alpha}+r_{J+n+1,\alpha}$ of some $k/Q_{J+n}$ with $k\in S_{n,\alpha}$.
By induction on $n\ge0$, $C_{J,\alpha}\subseteq\bigcup_{k\in S_{n,\alpha}}I_{J+n,\alpha}(k)$
and the cardinality $|S_{n,\alpha}|$ satisfies
(\texttt{WeightedDimH.card\_wsurvSet\_le})
\[
|S_{n,\alpha}|
\le(Q_J+1)\prod_{i=0}^{n-1}\bigl(K(J+i+1)^{1-\alpha}\bigr)
\le(Q_J+1)K^n\bigl((J+n)!\bigr)^{1-\alpha},
\]
because $\prod_{i=0}^{n-1}(J+i+1)=(J+n)!/J!\le(J+n)!$. By
Lemma~\ref{lem:factorial}, $Q_{J+n}\ge(J+n)!$, and $(J+n)^\alpha\ge1$, so each
interval $I_{J+n,\alpha}(k)$ has diameter
\[
\operatorname{diam}I_{J+n,\alpha}(k)=2r_{J+n,\alpha}
=\frac{2t}{(J+n)^\alpha Q_{J+n}}\le\frac{2t}{(J+n)!}\longrightarrow0
\qquad(n\to\infty).
\]
Now fix any exponent $s>1-\alpha$, and put $\beta:=s-(1-\alpha)>0$. For every
$\delta>0$ and all sufficiently large $n$, $\{I_{J+n,\alpha}(k):k\in S_{n,\alpha}\}$
is a finite $\delta$-cover of $C_{J,\alpha}$ whose $s$-dimensional Hausdorff
covering cost satisfies
\[
\sum_{k\in S_{n,\alpha}}\bigl(\operatorname{diam}I_{J+n,\alpha}(k)\bigr)^s
\le |S_{n,\alpha}|\left(\frac{2t}{(J+n)!}\right)^s
\le(Q_J+1)(2t)^s K^{-J}\cdot K^{J+n}\bigl((J+n)!\bigr)^{1-\alpha-s}.
\]
With $m=J+n$ and $K^{1/\beta}>0$, the factor $K^m(m!)^{1-\alpha-s}=K^m(m!)^{-\beta}=\bigl((K^{1/\beta})^m/m!\bigr)^\beta$
satisfies, via the elementary factorial lower bound $m!\ge(m/\ee)^m$,
\[
0\le K^m(m!)^{-\beta}=\left(\frac{(K^{1/\beta})^m}{m!}\right)^\beta
\le\left(\frac{\ee K^{1/\beta}}{m}\right)^{\beta m}\longrightarrow0
\qquad(m=J+n\to\infty)
\]
(formalized in \texttt{WeightedPhase.\allowbreak content\_\allowbreak core\_\allowbreak tendsto} via
\texttt{FloorSemiring.\allowbreak tendsto\_\allowbreak pow\_\allowbreak div\_\allowbreak factorial\_\allowbreak atTop} and continuity of
$u\mapsto u^\beta$ at $0$). Hence $\mathcal H_\delta^s(C_{J,\alpha})=0$ for all
$\delta>0$, so $\mathcal H^s(C_{J,\alpha})=0$ and $\dimH C_{J,\alpha}\le s$.
Since $s>1-\alpha$ was arbitrary, letting $s\downarrow1-\alpha$ gives
$\dimH C_{J,\alpha}\le1-\alpha$ for every $J\ge1$, and countable stability over
$J\ge1$ and $m\in\Z$ yields $\dimH E_{\alpha,c}\le1-\alpha$.
\end{proof}

\paragraph{Lean status.} Theorem~\ref{thm:upper} is kernel-checked at Layer~A
as \texttt{WeightedFinal.\allowbreak dimH\_levelSet\ensuremath{\alpha}\_le},
assembled through \texttt{WeightedFiber.lean}, \texttt{WeightedDimHCore.lean},
\texttt{WeightedUpper.lean}, and \texttt{WeightedFinal.lean}. Its axiom audit
reports only the three base axioms.

\begin{theorem}[Prescribed limits at subcritical weights, Layer~A in Lean]\label{conj:weight}
Let $0\le\alpha<1$. For $\alpha>0$ take any $c\ge0$; for $\alpha=0$
take $0\le c\le1/2$. Then
\[
 \dimH\{x\in[0,1]: j^\alpha\|Q_jx\|\longrightarrow c\}=1-\alpha .
\]
In particular $\dimH\{x\in\R:\Lam_\alpha(x)\le c\}=1-\alpha$ for every
finite $c\ge0$, including $c=0$.
(Lean, module \texttt{WeightedExactDim}:
\texttt{dimH\_\allowbreak exactLimitSet\ensuremath{\alpha}\_eq},
\texttt{dimH\_\allowbreak levelSet\ensuremath{\alpha}\_eq}, and
\texttt{dimH\_\allowbreak apprExponent\ensuremath{\alpha}\_eq\_\allowbreak zero};
\texttt{\#print axioms} returns \texttt{[propext, Classical.choice, Quot.sound]} only.)
\end{theorem}

\begin{proof}
Write $E_{\alpha}^{\mathrm{lim}}(c):=\{x\in[0,1]:j^\alpha\|Q_jx\|\to c\}$. Since
$j^\alpha\|Q_jx\|\to c$ implies $\Lam_\alpha(x)=\limsup_{j\to\infty}j^\alpha\|Q_jx\|=c$,
we have $E_{\alpha}^{\mathrm{lim}}(c)\subseteq\{x\in\R:\Lam_\alpha(x)\le c\}$, so
Theorem~\ref{thm:upper} immediately gives the upper bound
$\dimH E_{\alpha}^{\mathrm{lim}}(c)\le1-\alpha$.

We now give the complete paper proof of the lower bound
$\dimH E_{\alpha}^{\mathrm{lim}}(c)\ge1-\alpha$. For $j\ge1$, put
\[
w_j:=\frac{j^{-\alpha}}{\log(j+2)},\qquad r_j:=\frac{w_j}{Q_j}.
\]
Choose an integer $N\ge3$ large enough that $w_j<1/8$ for all $j\ge N$, and
also $cj^{-\alpha}+w_j\le1/2$ for all $j\ge N$ whenever either ($\alpha>0$,
$c\ge0$) or ($\alpha=0$, $0\le c<1/2$). Define the target fractional shift
$\theta_j\in[w_j,1/2-w_j]$ for $j\ge N$ by clamping $cj^{-\alpha}$ away from the
half-integer boundary:
\[
\theta_j:=\begin{cases}
w_j,&\text{if }c=0,\\[2pt]
cj^{-\alpha},&\text{if }c>0\text{ and }(\alpha>0\text{ or }\alpha=0,c<1/2),\\[2pt]
\tfrac12-w_j,&\text{if }\alpha=0\text{ and }c=\tfrac12.
\end{cases}
\]
(Note why the boundary adjustment $\theta_j=1/2-w_j$ is needed in the endpoint
case $\alpha=0,c=1/2$, and $\theta_j=w_j$ when $c=0$: if one centered directly
at $(k+1/2)/Q_j$, points on the right half of the interval would cross the
half-integer $k+1/2$ where $y\mapsto\|y\|$ folds back, whereas with center
$(k+\theta_j)/Q_j$ and radius $r_j=w_j/Q_j$, every point $x$ in the closed
interval $I_j(k):=[(k+\theta_j-w_j)/Q_j,(k+\theta_j+w_j)/Q_j]$ satisfies
$k\le k+\theta_j-w_j\le Q_jx\le k+\theta_j+w_j\le k+1/2$. On $[k,k+1/2]$ the
distance to the nearest integer is given without folding by
$\|Q_jx\|=Q_jx-k\in[\theta_j-w_j,\theta_j+w_j]$.) In particular, in the case
$\alpha=0,c=1/2$, for $x\in I_j(k)$ we have
$Q_jx\in[k+1/2-2w_j,k+1/2]$, so
\[
\|Q_jx\|=\frac12-\left(k+\frac12-Q_jx\right)\in\left[\frac12-2w_j,\frac12\right]
\quad\Longrightarrow\quad
\bigl|\|Q_jx\|-\tfrac12\bigr|\le2w_j=\frac{2}{\log(j+2)}\longrightarrow0,
\]
and in all three cases every $x\in I_j(k)$ satisfies
\[
\bigl|j^\alpha\|Q_jx\|-c\bigr|
\le j^\alpha\bigl|\|Q_jx\|-\theta_j\bigr|+\bigl|j^\alpha\theta_j-c\bigr|
\le j^\alpha w_j+j^\alpha w_j=\frac{2}{\log(j+2)}\longrightarrow0
\qquad(j\to\infty).
\]
Moreover, for fixed $j\ge N$ and distinct integers $k_1\ne k_2$, the distance
between the centers of $I_j(k_1)$ and $I_j(k_2)$ is
$|k_1-k_2|/Q_j\ge1/Q_j>2w_j/Q_j=2r_j$ (since $w_j<1/8<1/2$), so the level-$j$
intervals $\{I_j(k):k\in\Z\}$ are pairwise disjoint.

Next we bound the number of strictly nested level-$(j+1)$ intervals inside a
given parent interval $I_j(k)$ ($j\ge N$). A child interval
$I_{j+1}(k')=[c_{j+1}(k')-r_{j+1},c_{j+1}(k')+r_{j+1}]$ with center
$c_{j+1}(k')=(k'+\theta_{j+1})/Q_{j+1}$ is contained in
$I_j(k)=[c_j(k)-r_j,c_j(k)+r_j]$ ($c_j(k)=(k+\theta_j)/Q_j$) if and only if
$|c_{j+1}(k')-c_j(k)|\le r_j-r_{j+1}$, which is equivalent to requiring the
integer $k'$ to lie in the closed interval $[A_j(k),B_j(k)]$ where
\[
A_j(k):=Q_{j+1}\bigl(c_j(k)-(r_j-r_{j+1})\bigr)-\theta_{j+1},\qquad
B_j(k):=Q_{j+1}\bigl(c_j(k)+(r_j-r_{j+1})\bigr)-\theta_{j+1}.
\]
The length of $[A_j(k),B_j(k)]$ is independent of the parent index $k$ and of
the shifts $\theta_j,\theta_{j+1}$:
\[
W_j:=B_j(k)-A_j(k)=2(r_j-r_{j+1})Q_{j+1}
=2\left(w_j\frac{Q_{j+1}}{Q_j}-w_{j+1}\right).
\]
Since any closed real interval $[A,B]$ of length $W=B-A\ge1$ contains at least
$\lfloor B\rfloor-\lceil A\rceil+1>W-1$ integers, the number of eligible child
centers $k'\in\Z$ with $I_{j+1}(k')\subseteq I_j(k)$ is at least
$2(r_j-r_{j+1})Q_{j+1}-1$. Using $Q_{j+1}=(4j+2)Q_j+Q_{j-1}>(4j+2)Q_j$ from
\eqref{eq:rec} and $w_{j+1}\le w_j$, we have
\[
2(r_j-r_{j+1})Q_{j+1}-1
=2w_j\frac{Q_{j+1}}{Q_j}-2w_{j+1}-1
\ge 2w_j(4j+2)-2w_j-1
=\frac{8j^{1-\alpha}+2j^{-\alpha}}{\log(j+2)}-1.
\]
Since $0\le\alpha<1$, $j^{1-\alpha}/\log(j+2)\to\infty$ as $j\to\infty$.
Increasing $N$ if necessary so that $7j^{1-\alpha}/\log(j+2)\ge1$ and
$b_j:=\lfloor j^{1-\alpha}/\log(j+2)\rfloor\ge2$ for all $j\ge N$, we obtain
\[
2(r_j-r_{j+1})Q_{j+1}-1\ge\frac{j^{1-\alpha}}{\log(j+2)}\ge b_j\ge2
\qquad\text{for all }j\ge N,
\]
so every parent interval $I_j(k)$ ($j\ge N$) contains at least $b_j$ pairwise
disjoint child intervals $I_{j+1}(k')$.

Choose $N$ large enough that $I_N(1)\subseteq(0,1)$. Construct a homogeneous
Cantor tree by starting at level $N$ with the single root interval $I_N(1)$
($\mathcal E_N:=\{1\}$), and for each $j\ge N$ and each retained parent
$k\in\mathcal E_j$, keeping exactly $b_j$ child indices $k'\in\mathcal E_{j+1}$
with $I_{j+1}(k')\subseteq I_j(k)$ (for instance, the first $b_j$ integers
$\lceil A_j(k)\rceil+m$, $0\le m<b_j$). Then at each level $n\ge N$ there are
exactly
\[
M_n:=\prod_{j=N}^{n-1}b_j
\]
pairwise disjoint closed intervals $\{I_n(k):k\in\mathcal E_n\}$ of radius
$r_n=w_n/Q_n$, and their compact intersection
\[
K:=\bigcap_{n=N}^\infty\bigcup_{k\in\mathcal E_n}I_n(k)\subseteq[0,1]
\]
is nonempty and satisfies $K\subseteq E_{\alpha}^{\mathrm{lim}}(c)$. Let $\mu$
be the unique Borel probability measure on $K$ assigning equal mass
$\mu(I_n(k))=M_n^{-1}$ to each of the $M_n$ retained level-$n$ intervals ($n\ge N$).

To estimate the asymptotic growth of $\log M_n$ and $\log Q_{n+1}$, observe
that for $j\ge N$,
\[
\frac{j^{1-\alpha}}{2\log(j+2)}\le\frac{j^{1-\alpha}}{\log(j+2)}-1
<b_j\le\frac{j^{1-\alpha}}{\log(j+2)},
\]
so $\log b_j=(1-\alpha)\log j-\log\log(j+2)+O(1)$. By integral comparison for
the monotone function $t\mapsto\log t$,
$\sum_{j=N}^{n-1}\log j=\int_1^n\log t\,dt+O(\log n)=n\log n-n+O(\log n)$, and
$\sum_{j=N}^{n-1}\log\log(j+2)\le n\log\log(n+2)$, which gives
\begin{equation}\label{eq:log-Mn-asymp}
\log M_n=\sum_{j=N}^{n-1}\log b_j
=(1-\alpha)n\log n+O(n+n\log\log n)=(1-\alpha)n\log n+o(n\log n).
\end{equation}
Similarly, from $4jQ_j<(4j+2)Q_j<Q_{j+1}\le(4j+3)Q_j$ for all $j\ge1$, taking
logarithms and summing from $j=1$ to $n$ gives
\[
\log Q_{n+1}=\sum_{j=1}^n\log\frac{Q_{j+1}}{Q_j}
=\sum_{j=1}^n\bigl(\log j+\log 4+O(1/j)\bigr)=n\log n+O(n).
\]

Now fix any $0<s<1-\alpha$. By \eqref{eq:log-Mn-asymp} and
$\log Q_{n+1}=n\log n+O(n)$, we have
\begin{align*}
\log\left(\frac{3}{M_n Q_{n+1}^{-s}}\right)
&=s\log Q_{n+1}-\log M_n+\log 3\\
&=-\bigl((1-\alpha)-s\bigr)n\log n+o(n\log n)\longrightarrow-\infty
\qquad(n\to\infty).
\end{align*}
Hence we may choose an index $N_s\ge N$ such that
\begin{equation}\label{eq:mass-scale-bound}
\frac{3}{M_n}\le Q_{n+1}^{-s}\qquad\text{for all }n\ge N_s.
\end{equation}
Let $\{U_i\}_{i\ge1}$ be any countable cover of $K$ by nonempty sets of diameter
$0<h_i:=\operatorname{diam}(U_i)<Q_{N_s}^{-1}$ (if $\operatorname{diam}(U_i)=0$,
then $U_i$ is a singleton and $\mu(U_i)=0$ since $M_n^{-1}\to0$). For each $i$,
since $Q_n$ is strictly increasing to $\infty$, there exists a unique integer
$n_i\ge N_s$ such that
\[
Q_{n_i+1}^{-1}\le h_i<Q_{n_i}^{-1}.
\]
The set $U_i$ is contained in a closed interval $[x_0,x_0+h_i]$ of length
$h_i<Q_{n_i}^{-1}$. If a level-$n_i$ interval $I_{n_i}(k)$ ($k\in\Z$) intersects
$[x_0,x_0+h_i]$, then its center $c_{n_i}(k)=(k+\theta_{n_i})/Q_{n_i}$ must lie
in $[x_0-r_{n_i},x_0+h_i+r_{n_i}]$, which has length
$h_i+2r_{n_i}<Q_{n_i}^{-1}+2w_{n_i}Q_{n_i}^{-1}<(1+1/4)Q_{n_i}^{-1}<2Q_{n_i}^{-1}$.
Since the centers $c_{n_i}(k)$ ($k\in\Z$) form a grid of spacing exactly
$Q_{n_i}^{-1}$, an interval of length $<2Q_{n_i}^{-1}$ can contain at most $3$
grid points $c_{n_i}(k)$ (if it contained $4$ points $k_1<k_2<k_3<k_4$, its
length would be at least $(k_4-k_1)Q_{n_i}^{-1}\ge3Q_{n_i}^{-1}$). Thus $U_i$
intersects at most $3$ retained level-$n_i$ intervals $I_{n_i}(k)$
($k\in\mathcal E_{n_i}$). Because $\mu$ is supported on
$\bigcup_{k\in\mathcal E_{n_i}}I_{n_i}(k)$ and each $I_{n_i}(k)$ has measure
$M_{n_i}^{-1}$, it follows from \eqref{eq:mass-scale-bound} and
$Q_{n_i+1}^{-1}\le h_i$ that
\[
\mu(U_i)\le\frac{3}{M_{n_i}}\le Q_{n_i+1}^{-s}\le h_i^s=\bigl(\operatorname{diam}(U_i)\bigr)^s.
\]
Summing over the cover $\{U_i\}$ yields
\[
\sum_{i}\bigl(\operatorname{diam}(U_i)\bigr)^s
\ge\sum_{i}\frac{3}{M_{n_i}}\ge\sum_{i}\mu(U_i)\ge\mu(K)=1.
\]
Taking the infimum over all countable $\delta$-covers with
$0<\delta<Q_{N_s}^{-1}$ and letting $\delta\downarrow0$ gives
$\mathcal H^s(E_{\alpha}^{\mathrm{lim}}(c))\ge\mathcal H^s(K)\ge1>0$, and therefore
$\dimH E_{\alpha}^{\mathrm{lim}}(c)\ge s$. Letting $s\uparrow1-\alpha$ proves
$\dimH E_{\alpha}^{\mathrm{lim}}(c)\ge1-\alpha$, establishing the exact identity
$\dimH E_{\alpha}^{\mathrm{lim}}(c)=1-\alpha$.

Finally, for any finite $c\ge0$ (including $c=0$), we have the inclusions
\[
E_{\alpha}^{\mathrm{lim}}(0)
\subseteq\{x\in\R:\Lam_\alpha(x)=0\}
\subseteq\{x\in\R:\Lam_\alpha(x)\le c\}.
\]
The lower bound $\dimH E_{\alpha}^{\mathrm{lim}}(0)=1-\alpha$ just proved and
the upper bound $\dimH\{x\in\R:\Lam_\alpha(x)\le c\}\le1-\alpha$ from
Theorem~\ref{thm:upper} squeeze both sets to have exact Hausdorff dimension
$1-\alpha$.

The Lean formalization in \texttt{EpiFormal/\allowbreak WeightedExactDim.lean}
certifies the same lower bound at Layer~A for every $0<s<1-\alpha$ by taking
polynomial window widths $w_j=j^{-\alpha-\varepsilon}$ with
$\varepsilon=(1-\alpha-s)/2>0$, clamping the target shift into
$[w_j,1/2-w_j]$, and applying compactness plus finite-tree counting directly to
the definition of Hausdorff measure to establish
$\mathcal H^s(K)\ge1$ (\texttt{one\_le\_hausdorffMeasure\_cantorSet}).
\end{proof}

For $\alpha=0$ the dimension assertion also follows from Corollary~8 of
\cite{BLR} with constant target $c$, by the distance inequality used
above. For $0<\alpha<1$, condition~(1.6) of their Theorem~7 fails because
$\log Q_{j+1}/\log Q_j\to1$. The lower bound in the paper proof above can also
be read from the homogeneous Moran bounds of Feng, Wen and Wu \cite{FWW}: the
retained level has $\log M_n\sim(1-\alpha)n\log n$ and both bounding logarithmic
length scales are asymptotic to $n\log n$. The independent covering upper bound
(Theorem~\ref{thm:upper}) applies to the full prescribed-limit set and sublevel
set.

\section{Critical levels, exact Hausdorff gauges, the positive-level trichotomy, and general recurrences}\label{sec:gauge}

\subsection{Critical levels and exact gauges}

For an integer $k\ge1$, set $b=2k+1$. Define an increasing continuous
gauge $h_k$ by $h_k(0)=0$, $h_k(1/Q_n)=b^{-n}$ ($n\ge1$), and
linear interpolation on each intervening interval $[1/Q_{n+1},1/Q_n]$ (with
$h_k(t)=h_k(1/Q_1)$ for $t\ge1/Q_1$). Only its values near zero matter. Write
$\mathcal H^{h_k}$ for Hausdorff measure with covering cost
$h_k(\operatorname{diam} U)$.

\begin{theorem}[Size of a critical spectral level, Layer~A in Lean]\label{thm:gauge}
On the circle $\R/\Z$, the exact level $\{x:\Lam(x)=k/4\}$ has
infinite but $\sigma$-finite $\mathcal H^{h_k}$ measure.
The strict sublevel $\{x:\Lam(x)<k/4\}$ has zero $\mathcal H^{h_k}$ measure.
In the family of gauges $g_\beta$ given by $g_\beta(1/Q_n)=\beta^{-n}$
and the same linear interpolation ($\beta>1$), the exact level has zero
$\mathcal H^{g_\beta}$ measure for $\beta>2k+1$, whereas for $1<\beta<2k+1$ its
$\mathcal H^{g_\beta}$ measure is not $\sigma$-finite.
\end{theorem}

\begin{proof}
\emph{Digits.} Use $M_j$, $u_j=Q_jx-M_j$, $A_j=4j+2$ and
$d_j=M_{j+1}-A_jM_j-M_{j-1}=A_ju_j+u_{j-1}-u_{j+1}$ as in
Section~\ref{sec:class}, with a fixed Borel rule at ties (for instance,
$M_j=\lfloor Q_jx+1/2\rfloor$). Ties occur only on the
countable set $\bigcup_{j\ge1}\{x:Q_jx\in\Z+\tfrac12\}$, which is null for every
non-atomic measure and is avoided by the Cantor construction below.
If $\Lam(x)\le k/4$, eventually $|u_j|\le(k+1)/j$, so
$d_j-4ju_j=2u_j+u_{j-1}-u_{j+1}\to0$ and
$\limsup_{j\to\infty}|d_j|=4\Lam(x)\le k$. Since $d_j\in\Z$, the integer digits
eventually satisfy $|d_j|\le k$. For the \emph{strict} sublevel $\Lam(x)<k/4$,
Corollary~\ref{thm:sub} improves this to $|d_j|\le k-1$ eventually.

\emph{Upper bounds.} Partition the non-strict sublevel $\{\Lam\le k/4\}$ on
$[0,1]$ (or $\R/\Z$) into countably many classes
$\mathcal C(N,m_{N-1},m_N)$ indexed by a common starting index $N\ge2$ and the
two starting integers $(M_{N-1},M_N)=(m_{N-1},m_N)\in\Z^2$, requiring
$|u_j|\le(k+1)/j$ and $|d_j|\le k$ for all $j\ge N$. With a fixed Borel rounding
rule these classes are Borel sets. Because $M_{j+1}=A_jM_j+M_{j-1}+d_j$ is
determined inductively from $(M_{N-1},M_N)$ by the $n-N$ digits
$d_N,\dots,d_{n-1}\in\{-k,\dots,k\}$, by level $n>N$ a class has at most
$b^{\,n-N}$ possible values of $M_n$ ($b=2k+1$). Each value of $M_n$ confines
$x$ to the interval $[M_n/Q_n-(k+1)/(nQ_n),M_n/Q_n+(k+1)/(nQ_n)]$ of diameter
$2(k+1)/(nQ_n)\le1/Q_n$ for all $n\ge2(k+1)$ (or directly inside
$[M_n/Q_n-1/(2Q_n),M_n/Q_n+1/(2Q_n)]$ of diameter $1/Q_n$ using $|u_n|\le1/2$).
Since $1/Q_n\to0$, for every $\delta>0$ and all large $n$ this gives a finite
$\delta$-cover of $\mathcal C(N,m_{N-1},m_N)$ with $\mathcal H_\delta^{h_k}$
cost at most
\[
b^{\,n-N}h_k(1/Q_n)=b^{\,n-N}b^{-n}=b^{-N}<\infty.
\]
Letting $\delta\downarrow0$ gives
$\mathcal H^{h_k}(\mathcal C(N,m_{N-1},m_N))\le b^{-N}<\infty$. For the strict
sublevel $\{\Lam<k/4\}$, the allowed tail digits satisfy
$d_j\in\{-(k-1),\dots,k-1\}$, so the count at level $n$ uses at most
$(2k-1)^{n-N}$ values of $M_n$ with $2k-1<b=2k+1$. The $\mathcal H_\delta^{h_k}$
covering cost of each class is therefore at most
\[
(2k-1)^{n-N}h_k(1/Q_n)=b^{-N}\left(\frac{2k-1}{2k+1}\right)^{n-N}\longrightarrow0
\qquad(n\to\infty),
\]
so every class is $\mathcal H^{h_k}$-null, and countable subadditivity yields
\begin{equation}\label{eq:strict-null}
 \mathcal H^{h_k}\bigl(\{x:\Lam(x)<k/4\}\bigr)=0 .
\end{equation}

\emph{Positivity.} Fix a large integer $N\ge2$ (chosen below), set
$M_{N-1}=M_N=0$, and prescribe arbitrary digits
$\omega=(d_N,d_{N+1},\dots)\in\{-k,\ldots,k\}^{\{j\ge N\}}$ via the linear
recurrence
\[
M_{j+1}=A_jM_j+M_{j-1}+d_j\qquad(j\ge N).
\]
Define the determinant sequence $W_j:=M_{j+1}Q_j-M_jQ_{j+1}$ for $j\ge N-1$, so
that $W_{N-1}=M_NQ_{N-1}-M_{N-1}Q_N=0$. Using $M_{j+1}=A_jM_j+M_{j-1}+d_j$ and
$Q_{j+1}=A_jQ_j+Q_{j-1}$ from \eqref{eq:rec}, the $A_jM_jQ_j$ terms cancel and
we obtain the first-order recurrence
\[
W_j=(A_jM_j+M_{j-1}+d_j)Q_j-M_j(A_jQ_j+Q_{j-1})
=d_jQ_j-(M_jQ_{j-1}-M_{j-1}Q_j)=d_jQ_j-W_{j-1}
\]
for all $j\ge N$. We claim by induction on $j\ge N-1$ that
\begin{equation}\label{eq:Wj-bound}
|W_j|\le 2kQ_j\qquad\text{for all }j\ge N-1.
\end{equation}
Indeed, for $j=N-1$, $|W_{N-1}|=0\le2kQ_{N-1}$. For $j\ge N$, since $Q_j\ge2Q_{j-1}$
(in fact $Q_j\ge(4j-2)Q_{j-1}\ge6Q_{j-1}$ for $j\ge2$), the inductive
hypothesis $|W_{j-1}|\le2kQ_{j-1}\le kQ_j$ and $|d_j|\le k$ yield
\[
|W_j|\le|d_j|Q_j+|W_{j-1}|\le kQ_j+2kQ_{j-1}\le kQ_j+kQ_j=2kQ_j,
\]
proving \eqref{eq:Wj-bound}. Dividing $W_j$ by $Q_jQ_{j+1}$ gives the adjacent
rational difference
\[
\frac{M_{j+1}}{Q_{j+1}}-\frac{M_j}{Q_j}=\frac{W_j}{Q_jQ_{j+1}},\qquad
\left|\frac{M_{j+1}}{Q_{j+1}}-\frac{M_j}{Q_j}\right|\le\frac{2k}{Q_{j+1}}.
\]
Because $Q_{j+2}\ge(4j+6)Q_{j+1}\ge2Q_{j+1}$ for all $j\ge1$, we have
$Q_{m+1+r}\ge2^r Q_{m+1}$ for all $m\ge N$ and $r\ge0$. Consequently, the
telescoping series converges uniformly on $\{-k,\dots,k\}^{\{j\ge N\}}$ to a
real limit
\[
x(\omega):=\lim_{j\to\infty}\frac{M_j}{Q_j}
=\frac{M_m}{Q_m}+\sum_{j=m}^\infty\frac{W_j}{Q_jQ_{j+1}}\qquad(m\ge N),
\]
with the uniform tail bound
\begin{equation}\label{eq:telescoping-tail}
\left|x(\omega)-\frac{M_m}{Q_m}\right|
\le\sum_{j=m}^\infty\frac{|W_j|}{Q_jQ_{j+1}}
\le\sum_{j=m}^\infty\frac{2k}{Q_{j+1}}
\le\frac{2k}{Q_{m+1}}\sum_{r=0}^\infty 2^{-r}
=\frac{4k}{Q_{m+1}}\qquad\text{for all }m\ge N.
\end{equation}
Multiplying \eqref{eq:telescoping-tail} by $Q_m$ gives, for every $m\ge N$,
\[
|Q_mx(\omega)-M_m|\le\frac{4kQ_m}{Q_{m+1}}\le\frac{4k}{4m+2}<\frac{k}{m}.
\]
Choose $N>8k$ (so $4k/Q_{N+1}<1/8$). Then for every $j\ge N$,
$|Q_jx(\omega)-M_j|<k/j<1/8<1/2$, which ensures that $M_j$ is the \emph{unique}
nearest integer to $Q_jx(\omega)$, that $u_j(x(\omega))=Q_jx(\omega)-M_j$ avoids
all half-integer ties, and that the true recurrence digit of $x(\omega)$ at
index $j$ is $M_{j+1}-A_jM_j-M_{j-1}=d_j$. Moreover,
$|4ju_j(x(\omega))-d_j|\le2|u_j|+|u_{j-1}|+|u_{j+1}|\to0$, so
$4\Lam(x(\omega))=\limsup_{j\to\infty}|d_j|\le k$. Hence the compact image
$K:=x(\{-k,\dots,k\}^{\{j\ge N\}})$ lies in the \emph{non-strict} sublevel
$\{x\in\R:\Lam(x)\le k/4\}$; because $M_N=0$ and $|x(\omega)|\le4k/Q_{N+1}$ by
\eqref{eq:telescoping-tail} at $m=N$, $K$ is contained in
$[-4k/Q_{N+1},4k/Q_{N+1}]$, an interval of length $8k/Q_{N+1}<1/4<1/2$ on which
the canonical projection $\R\to\R/\Z$ is an isometric embedding.

Equip $\{-k,\dots,k\}^{\{j\ge N\}}$ with the uniform Bernoulli product
probability measure assigning independent weight $b^{-1}=(2k+1)^{-1}$ to each
digit $d_j$, and let $\mu$ be its pushforward probability measure on $K$. Take
any interval $U\subset\R$ of length $h=\operatorname{diam}(U)$ with
$1/Q_{n+1}\le h<1/Q_n$ for some $n>N$. If $x,y\in K\cap U$, then for every index
$j$ with $N\le j<n$, we have $|x-y|\le h<1/Q_n$, $|Q_jx-M_j(x)|<1/8$, and
$|Q_jy-M_j(y)|<1/8$, so the triangle inequality gives
\[
|M_j(x)-M_j(y)|
\le|M_j(x)-Q_jx|+Q_j|x-y|+|Q_jy-M_j(y)|
<\frac18+\frac{Q_j}{Q_n}+\frac18=\frac14+\frac{Q_j}{Q_n}\le\frac14+\frac16<1.
\]
Since $M_j(x)$ and $M_j(y)$ are integers, $M_j(x)=M_j(y)$ for all $N\le j<n$
(and $M_{N-1}=0$ by definition). At level $n$, the same estimate gives
\[
|M_n(x)-M_n(y)|<\frac14+Q_n h<\frac14+1=\frac54<2,
\]
so for all $x\in K\cap U$ the integer $M_n(x)$ can take at most $3$ consecutive
integer values (differing by at most $1$ from $\min_{x\in K\cap U}M_n(x)$).
Because the integer sequence $M_{N-1},M_N,\dots,M_n$ uniquely determines the
first $n-N$ digits $d_N,\dots,d_{n-1}$ via $d_j=M_{j+1}-A_jM_j-M_{j-1}$, the
intersection $K\cap U$ is contained in at most $3$ cylinder sets of depth $n-N$,
each of $\mu$-measure $b^{-(n-N)}$. Since $h_k$ is nondecreasing and
$h\ge1/Q_{n+1}$, we have $h_k(h)\ge h_k(1/Q_{n+1})=b^{-(n+1)}$, and therefore
\[
\mu(U)=\mu(K\cap U)\le 3b^{-(n-N)}=3b^{N+1}b^{-(n+1)}\le 3b^{N+1}h_k(h).
\]
Summing over any countable cover $\{U_i\}$ of $K$ with
$\operatorname{diam}(U_i)<1/Q_{N+1}$ gives
$\sum_i h_k(\operatorname{diam}U_i)\ge(3b^{N+1})^{-1}\sum_i\mu(U_i)\ge(3b^{N+1})^{-1}>0$
(the mass distribution principle for $\mathcal H^{h_k}$), so
$\mathcal H^{h_k}(K)\ge(3b^{N+1})^{-1}>0$. By the isometric embedding of
$[-4k/Q_{N+1},4k/Q_{N+1}]$ into $\R/\Z$, the projection of $K$ to the circle has
the same positive $\mathcal H^{h_k}$ measure. Since $K\subseteq\{\Lam\le k/4\}$
and \eqref{eq:strict-null} shows that the strict sublevel $\{\Lam<k/4\}$ has
$\mathcal H^{h_k}$ measure zero, we conclude that
\[
 \mathcal H^{h_k}\bigl(\{x\in\R/\Z:\Lam(x)=k/4\}\bigr)
 \ \ge\ \mathcal H^{h_k}(K)-\mathcal H^{h_k}\bigl(\{\Lam<k/4\}\bigr)
 \ =\ \mathcal H^{h_k}(K)\ >\ 0 .
\]

\emph{Measurability, infinitude and $\sigma$-finiteness.} The functions
$x\mapsto j\|Q_jx\|$ are continuous on the circle $\R/\Z$. Hence $\Lam$ is Borel
measurable as their pointwise limit superior, every sublevel $\{\Lam\le c\}$ and
superlevel $\{\Lam\ge c\}=\bigcap_{m\ge1}\bigcup_{j\ge m}\{x:j\|Q_jx\|\ge c-1/m\}$
is Borel, and so is the exact level
$S_k:=\{x\in\R/\Z:\Lam(x)=k/4\}=\{\Lam\le k/4\}\cap\{\Lam\ge k/4\}$. By
Corollary~\ref{cor:infinite-level} (or Theorem~\ref{thm:finite}), $S_k$ has
Lebesgue measure zero. For any $y\in L=\Z+\Z\frac{\ee+1}{2}$, the triangle
inequality for $\|\cdot\|$ gives
$\bigl|j\|Q_j(x+y)\|-j\|Q_jx\|\bigr|\le j\|Q_jy\|\to0$ as $j\to\infty$ by
Theorem~\ref{thm:class} (the critical weight $j$ is essential here), so
$\Lam(x+y)=\Lam(x)$ and $S_k$ is invariant under translation by the subgroup
$\bar L:=L/\Z\subset\R/\Z$. Since $(\ee+1)/2$ is irrational, Kronecker's theorem
implies that $\bar L$ is dense in $\R/\Z$.

Suppose for a contradiction that $0<\mathcal H^{h_k}(S_k)<\infty$. Then the
normalized restriction
\[
\nu(B):=\frac{\mathcal H^{h_k}(S_k\cap B)}{\mathcal H^{h_k}(S_k)}
\qquad(B\subseteq\R/\Z\text{ Borel})
\]
is a Borel probability measure on $\R/\Z$ concentrated on $S_k$ ($\nu(S_k)=1$).
Because $\mathcal H^{h_k}$ is translation-invariant on $\R/\Z$ and $S_k+\ell=S_k$
for every $\ell\in\bar L$, $\nu$ is invariant under rotation by every
$\ell\in\bar L$: $\nu(B+\ell)=\nu(B)$ for all Borel $B\subseteq\R/\Z$. We show
that $\nu$ is invariant under rotation by \emph{every} $t\in\R/\Z$. Let
$f:\R/\Z\to\R$ be any continuous function. By uniform continuity on the compact
group $\R/\Z$, for any $t\in\R/\Z$ we may choose a sequence $\ell_m\in\bar L$
with $\ell_m\to t$ in $\R/\Z$, and then
$\sup_{x\in\R/\Z}|f(x-\ell_m)-f(x-t)|\to0$ as $m\to\infty$. By $\bar L$-invariance
of $\nu$ and bounded convergence,
\[
\int_{\R/\Z}f(x-t)\,d\nu(x)
=\lim_{m\to\infty}\int_{\R/\Z}f(x-\ell_m)\,d\nu(x)
=\lim_{m\to\infty}\int_{\R/\Z}f(x)\,d\nu(x)
=\int_{\R/\Z}f(x)\,d\nu(x).
\]
Thus the rotated Borel probability measure $\nu_t(B):=\nu(B+t)$ has the same
integral as $\nu$ on every continuous function $f\in C(\R/\Z)$ (or equivalently
the same Fourier coefficients $\int_{\R/\Z}\ee^{2\pi i m x}\,d\nu(x)=0$ for all
$m\in\Z\setminus\{0\}$), so $\nu_t=\nu$ for all $t\in\R/\Z$ by uniqueness of
Haar probability measure on $\R/\Z$: concretely, for every integer $m\ge1$,
partitioning $\R/\Z$ into $m$ disjoint half-open arcs of equal length $1/m$
gives mass $1/m$ to each arc, Hence every arc of rational length $p/m$ has
$\nu$-measure $p/m$, and monotonicity extends $\nu(I)=\operatorname{Leb}_{\R/\Z}(I)$
to all open arcs $I\subset\R/\Z$, which generate the Borel $\sigma$-algebra.
Thus $\nu=\operatorname{Leb}_{\R/\Z}$, contradicting
$1=\nu(S_k)=\operatorname{Leb}_{\R/\Z}(S_k)=0$. Therefore
$\mathcal H^{h_k}(S_k)=\infty$. Finally, $S_k$ has $\sigma$-finite
$\mathcal H^{h_k}$ measure because $S_k\subseteq\{\Lam\le k/4\}$, which is
covered by the countably many Borel classes $\mathcal C(N,m_{N-1},m_N)$ above,
each of finite $\mathcal H^{h_k}$ measure $\le b^{-N}<\infty$.

For the power-gauge family $g_\beta(1/Q_n)=\beta^{-n}$, on each linear
interpolation interval $[1/Q_{n+1},1/Q_n]$ both $g_\beta$ and $h_k$ are affine
with positive values at the endpoints, so the ratio $g_\beta(r)/h_k(r)$ is
monotone on $[1/Q_{n+1},1/Q_n]$ and therefore lies between the endpoint values
$(\beta/(2k+1))^{-n}$ and $(\beta/(2k+1))^{-n-1}$. Thus
$g_\beta(r)/h_k(r)\to0$ when $\beta>2k+1$ and $h_k(r)/g_\beta(r)\to0$ when
$1<\beta<2k+1$ as $r\downarrow0$, which are special cases of
Corollary~\ref{cor:exact-gauge}(ii)--(iii).

For comparison, Cabrelli, Mendivil, Molter and Shonkwiler
\cite[\S5, Theorem~3]{CMMS} construct, for a Cantor set whose gap
sequence is nonincreasing, a gauge $h$ for which the set is an $h$-set
($1/32\le\mathcal H^h\le1$, a finite positive measure), whereas Moreira, Erazo
and Angelini \cite{MEA} show that the classical Lagrange level at value $3$ is
never $\sigma$-finite with positive measure for any gauge. Here the gauge comes
from the recurrence's own scale, and the entire exact level has \emph{infinite}
but $\sigma$-finite $\mathcal H^{h_k}$ measure.
\end{proof}

\paragraph{Lean status.}
The null part \eqref{eq:strict-null} is kernel-checked at Layer~A in
\texttt{EpiFormal/\allowbreak GaugeNull.lean} as
\texttt{GaugeNull.\allowbreak mkMetric\_\allowbreak strictLevel\_\allowbreak null}
for any nondecreasing function $m:\R_{\ge0\infty}\to\R_{\ge0\infty}$ satisfying
$m(1/Q_n)\le C(2k+1)^{-n}$ ($n\ge1$), without requiring continuity of $m$; the
same module kernel-checks Corollary~\ref{cor:exact-gauge}(ii) on $\R$ as
\texttt{GaugeNull.\allowbreak mkMetric\_\allowbreak levelSet\_\allowbreak null\_\allowbreak of\_\allowbreak subgauge}.
In \texttt{EpiFormal/\allowbreak GaugeExact.lean}, $L$-translation invariance
(\texttt{exactLevel\_\allowbreak add\_\allowbreak latticeL}), the countable
finite-measure covering
(\texttt{exactLevel\_\allowbreak countable\_\allowbreak cover\_\allowbreak finite\_\allowbreak mkMetric}),
and---for arbitrary monotone continuous $m$ satisfying matching two-sided
bounds---the Cantor-tree lower bound
(\texttt{mkMetric\_\allowbreak cantorTree\_\allowbreak pos}), exact-level
positivity (\texttt{mkMetric\_\allowbreak exactLevel\_\allowbreak pos}),
exact-level infinitude on $\R$
(\texttt{mkMetric\_\allowbreak exactLevel\_\allowbreak eq\_\allowbreak top}),
and the non-$\sigma$-finiteness criterion
(\texttt{not\_\allowbreak countable\_\allowbreak union\_\allowbreak finite\_\allowbreak of\_\allowbreak ratio\_\allowbreak tendsto\_\allowbreak zero})
are likewise kernel-checked.

The external continuity hypothesis, the explicit piecewise-linear gauge
construction, and the native quotient-metric circle statements on
\texttt{UnitAddCircle} ($\R/\Z$) are completely discharged in
\texttt{EpiFormal/\allowbreak GaugeCircle.lean}:
the gauges $h_k$ (\texttt{critGauge k}), $g_2$ (\texttt{binaryGauge}), and
$g_\beta$ (\texttt{interpGauge \ensuremath{\beta}}) are constructed explicitly
with kernel-checked monotonicity (\texttt{critGauge\_\allowbreak mono}),
continuity on $\R_{\ge0\infty}$ including $0^+$ and $\infty$
(\texttt{critGauge\_\allowbreak continuous}), and exact endpoint and segment
values (\texttt{critGauge\_\allowbreak at\_\allowbreak inv\_\allowbreak Qsel},
\texttt{critGauge\_\allowbreak eq\_\allowbreak segVal}).
This yields the assumption-free instantiations of Theorem~\ref{thm:gauge},
Corollary~\ref{cor:exact-gauge}(i)--(iii), and Remark~\ref{rem:gauge-scale} both
on $\R$
(\texttt{mkMetric\_\allowbreak critGauge\_\allowbreak strictLevel\_\allowbreak null},
\texttt{mkMetric\_\allowbreak critGauge\_\allowbreak exactLevel\_\allowbreak pos},
\texttt{mkMetric\_\allowbreak critGauge\_\allowbreak exactLevel\_\allowbreak eq\_\allowbreak top},
\texttt{critGauge\_\allowbreak exactLevel\_\allowbreak sigma\_\allowbreak finite\_\allowbreak cover},
\texttt{mkMetric\_\allowbreak critGauge\_\allowbreak exactLevel\_\allowbreak inter\_\allowbreak Ioo\_\allowbreak eq\_\allowbreak top},
\texttt{mkMetric\_\allowbreak powerGauge\_\allowbreak exactLevel\_\allowbreak null\_\allowbreak of\_\allowbreak gt},
\texttt{mkMetric\_\allowbreak powerGauge\_\allowbreak exactLevel\_\allowbreak not\_\allowbreak sigma\_\allowbreak finite\_\allowbreak of\_\allowbreak lt},
and \texttt{critGauge\_\allowbreak adjacent\_\allowbreak level\_\allowbreak separation})
and natively on \texttt{UnitAddCircle} ($\R/\Z$)
(\texttt{mkMetric\_\allowbreak critGauge\_\allowbreak circleStrictLevel\_\allowbreak null},
\texttt{mkMetric\_\allowbreak circleLevelSet\_\allowbreak null\_\allowbreak of\_\allowbreak subgauge},
\texttt{mkMetric\_\allowbreak critGauge\_\allowbreak circleExactLevel\_\allowbreak inter\_\allowbreak open\_\allowbreak eq\_\allowbreak top},
\texttt{mkMetric\_\allowbreak critGauge\_\allowbreak circleExactLevel\_\allowbreak eq\_\allowbreak top},
\texttt{critGauge\_\allowbreak circleExactLevel\_\allowbreak sigma\_\allowbreak finite\_\allowbreak cover},
\texttt{mkMetric\_\allowbreak critGauge\_\allowbreak circleExactLevel\_\allowbreak not\_\allowbreak sigma\_\allowbreak finite\_\allowbreak of\_\allowbreak ratio},
\texttt{mkMetric\_\allowbreak powerGauge\_\allowbreak circleExactLevel\_\allowbreak null\_\allowbreak of\_\allowbreak gt},
and
\texttt{mkMetric\_\allowbreak powerGauge\_\allowbreak circleExactLevel\_\allowbreak not\_\allowbreak sigma\_\allowbreak finite\_\allowbreak of\_\allowbreak lt}),
proved constructively by splitting the Cantor tree into depth-$M$ cylinders,
packing $\Omega\bigl((b-a)Q_{J+M+1}\bigr)$ disjoint $L$-translates using
$\delta_j/2\in L$, and transferring across the local isometry
$\R\to\R/\Z$ on intervals of length $\le1/2$. For the binary exact-limit
regime of Theorem~\ref{thm:trichotomy}(ii), the binary Cantor-tree lower bound,
global and open-set $\mathcal H^{g_2}$-infinitude, $\sigma$-finite covering,
$\mathcal H^{h_k}$-nullity, and the $g_\beta$ threshold dichotomy ($\beta>2$ vs.\
$1<\beta<2$) are likewise kernel-checked on both $\R$ and \texttt{UnitAddCircle}
(\texttt{mkMetric\_\allowbreak binaryGauge\_\allowbreak unsigned\allowbreak ExactLimit\_\allowbreak pos},
\texttt{mkMetric\_\allowbreak binaryGauge\_\allowbreak unsigned\allowbreak ExactLimit\_\allowbreak inter\_\allowbreak Ioo\_\allowbreak eq\_\allowbreak top},
\texttt{mkMetric\_\allowbreak binaryGauge\_\allowbreak circle\allowbreak Unsigned\allowbreak ExactLimit\_\allowbreak inter\_\allowbreak open\_\allowbreak eq\_\allowbreak top},
\texttt{mkMetric\_\allowbreak binaryGauge\_\allowbreak circle\allowbreak Unsigned\allowbreak ExactLimit\_\allowbreak eq\_\allowbreak top},
\texttt{binaryGauge\_\allowbreak circle\allowbreak Unsigned\allowbreak ExactLimit\_\allowbreak sigma\_\allowbreak finite\_\allowbreak cover},
\texttt{critGauge\_\allowbreak circle\allowbreak Unsigned\allowbreak ExactLimit\_\allowbreak null},
\texttt{mkMetric\_\allowbreak powerGauge\_\allowbreak circle\allowbreak Unsigned\allowbreak ExactLimit\_\allowbreak null\_\allowbreak of\_\allowbreak gt},
and
\texttt{mkMetric\_\allowbreak powerGauge\_\allowbreak circle\allowbreak Unsigned\allowbreak ExactLimit\_\allowbreak not\_\allowbreak sigma\_\allowbreak finite\_\allowbreak of\_\allowbreak lt}).
All 53 audited declarations in \texttt{GaugeCircle.lean} report only
\texttt{propext}, \texttt{Classical.choice}, and \texttt{Quot.sound}, with zero
\texttt{sorryAx}.

The gauge $h_k$ is exact for the level $k/4$ in the following sense.
Here a gauge is a nondecreasing function $g:[0,\infty)\to[0,\infty)$
with $g(0)=0$, and $\mathcal H^g$ is the Hausdorff measure on the circle
with covering cost $g(\operatorname{diam}U)$.

\begin{corollary}[Exact critical gauge]\label{cor:exact-gauge}
Let $k\ge1$ and $S_k=\{x\in\R/\Z:\Lam(x)=k/4\}$.
\begin{enumerate}[label=(\roman*)]
\item For every nonempty open arc $I\subseteq\R/\Z$,
$\mathcal H^{h_k}(S_k\cap I)=\infty$.
\item If $g$ is a gauge with $g(1/Q_n)/h_k(1/Q_n)\to0$ as $n\to\infty$, then
$\mathcal H^{g}(\{\Lam\le k/4\})=0$; in particular $\mathcal H^g(S_k)=0$.
(Lean on $\R$ and $\R/\Z$: \texttt{GaugeNull.\allowbreak mkMetric\_\allowbreak levelSet\_\allowbreak null\_\allowbreak of\_\allowbreak subgauge} and \texttt{GaugeCircle.\allowbreak mkMetric\_\allowbreak circleLevelSet\_\allowbreak null\_\allowbreak of\_\allowbreak subgauge}.)
\item If $g$ is a gauge with $h_k(t)/g(t)\to0$ as $t\to0^+$, then $S_k$ is
not $\sigma$-finite for $\mathcal H^g$; in particular
$\mathcal H^g(S_k)=\infty$.
\end{enumerate}
\end{corollary}

\begin{proof}
(i) Hausdorff measures on the circle are invariant under rotations, and
$S_k$ is invariant under rotation by the elements of $L$ (proof of
Theorem~\ref{thm:gauge}). Suppose $\mathcal H^{h_k}(S_k\cap I)<\infty$ for
some nonempty open arc $I$. Since $L$ is dense modulo one, the arcs $I+\ell$
($\ell\in L$) cover the circle, and by compactness finitely many of them,
say $I+\ell_1,\ldots,I+\ell_N$, already do. Then
$\mathcal H^{h_k}(S_k)\le\sum_{i=1}^N\mathcal H^{h_k}(S_k\cap(I+\ell_i))
=N\,\mathcal H^{h_k}(S_k\cap I)<\infty$, contradicting
Theorem~\ref{thm:gauge}.

(ii) In the proof of Theorem~\ref{thm:gauge}, the non-strict sublevel
$\{\Lam\le k/4\}$ is a countable union of classes. For a class with
starting index $N$ and every large $n$, at most $b^{\,n-N}$ intervals of
diameter at most $1/Q_n$ cover it, where $b=2k+1$. Since $g$ is
nondecreasing, the $g$-cost of this cover is at most
\[
 b^{\,n-N}g(1/Q_n)=b^{-N}\,\frac{g(1/Q_n)}{h_k(1/Q_n)}\longrightarrow0 .
\]
So every class is $\mathcal H^g$-null, and so is their countable union.

(iii) Fix $\varepsilon>0$ and choose $\delta>0$ with
$h_k(t)\le\varepsilon g(t)$ for $0\le t<\delta$. Every $\delta$-cover then
satisfies $\sum h_k(\operatorname{diam}U)\le
\varepsilon\sum g(\operatorname{diam}U)$, so
$\mathcal H^{h_k}(E)\le\varepsilon\,\mathcal H^g(E)$ for every set $E$.
Hence $\mathcal H^g(E)<\infty$ implies $\mathcal H^{h_k}(E)=0$. If $S_k$
were a countable union of sets $E_i$ with $\mathcal H^g(E_i)<\infty$, we
would get $\mathcal H^{h_k}(S_k)=0$, contradicting Theorem~\ref{thm:gauge}.
\end{proof}

\begin{corollary}[Arbitrary gauges at the exact and unsigned levels]\label{cor:liminf-gauges}
Let $g$ be a Hausdorff gauge, $r_n=1/Q_n$, and
\[
\rho_b(g)=\liminf_{n\to\infty} b^n g(r_n).
\]
For $S_k=\{\Lam=k/4\}$ take $b=2k+1$, and for
$T_k=\{x:j\|Q_jx\|\to k/4\}$ take $b=2$, where $k\ge1$ and both sets
are taken on the circle. In each case the named set is
$\mathcal H^g$-null if $\rho_b(g)=0$. If $0<\rho_b(g)<\infty$, it is
$\sigma$-finite and its intersection with every nonempty open set has
infinite $\mathcal H^g$ measure. If $\rho_b(g)=\infty$, its intersection
with every nonempty open set is non-$\sigma$-finite. For $b=2k+1$ the
zero assertion applies to the full sublevel $\{\Lam\le k/4\}$, and
$\rho_b(g)<\infty$ implies that $\{\Lam<k/4\}$ is $\mathcal H^g$-null.
\end{corollary}

\begin{proof}
The fixed initial-pair classes of Theorems~\ref{thm:gauge}
and~\ref{thm:trichotomy}(ii) have covers by at most $b^{n-N}$ intervals
of diameter $r_n$ for a fixed starting level $N$ and all large $n$.
Their gauge cost is bounded by $b^{-N}b^n g(r_n)$. Taking the lower
limit along permissible scales gives zero or finite measure for each
class, according as $\rho_b(g)$ is zero or finite, and the classes are
countable. For the strict sublevel the same argument uses $2k-1$
instead of $2k+1$. Bounded cost along a subsequence at base $2k+1$
then tends to zero along that subsequence at base $2k-1$.

Write $h_b(r_n)=b^{-n}$ with linear interpolation. If
$r_{n+1}\le r\le r_n$, then
\[
 h_b(r)\le b^{-n},\qquad g(r)\ge g(r_{n+1}).
\]
For any $0<A<\rho_b(g)$, eventually $b^{n+1}g(r_{n+1})\ge A$,
so $h_b(r)\le(b/A)g(r)$ for all sufficiently small $r$. The critical
reference measure is infinite on the intersection with every nonempty
open set, by Corollary~\ref{cor:exact-gauge}(i) and
Theorem~\ref{thm:trichotomy}(ii), respectively. The comparison gives the
same local infinitude for $g$. When $\rho_b(g)=\infty$, $A$ can be
arbitrarily large, so any set of finite $\mathcal H^g$ measure is
$\mathcal H^{h_b}$-null. A countable union of such sets cannot cover
a nonempty open intersection of the named set.
\end{proof}

\begin{remark}[Oscillating normalized costs]
Let $n_j=2^j$ and prescribe $g(r_n)=b^{-n_j}$ for
$n_j\le n<n_{j+1}$, interpolating linearly in $r$. The lower and upper
limits of $b^ng(r_n)$ are $1$ and $\infty$. Prescribing
$g(r_n)=b^{-(n_{j+1}-1)}$ instead gives lower and upper limits $0$ and $1$.
Both prescriptions are nondecreasing as functions of $r$ and tend to zero
at the origin. They can be made strictly increasing near zero by adding
$h_b^2$, which adds $b^{-n}$ to the normalized costs at the nodes.
\end{remark}

\paragraph{Lean verification.}
The accompanying \texttt{GaugeLiminfComplete.lean} verifies
Corollary~\ref{cor:liminf-gauges} on \texttt{UnitAddCircle}, including
the strict-sublevel assertion and the non-$\sigma$-finite conclusion on
each nonempty open set. Its three principal declarations are
\texttt{circleExactLevel\_\allowbreak liminf\_\allowbreak trichotomy},
\texttt{circleUnsigned\allowbreak ExactLimit\_\allowbreak liminf\_\allowbreak trichotomy}, and
\texttt{mkMetric\_\allowbreak circleStrictLevel\_\allowbreak null\_\allowbreak of\_\allowbreak liminf\_\allowbreak lt\_\allowbreak top}.
All eighteen audited declarations use only the three base axioms.

\begin{remark}[Level separation and interpolation-free logarithmic scale]\label{rem:gauge-scale}
The gauges $h_k$ strictly separate the critical levels. Indeed, for $k\ge2$,
$h_k(1/Q_n)/h_{k-1}(1/Q_n)=\bigl((2k-1)/(2k+1)\bigr)^n\to0$, and on each
interval $[1/Q_{n+1},1/Q_n]$ both gauges are linear, so their ratio lies
between its values at the endpoints. Hence $h_k(t)/h_{k-1}(t)\to0$ as
$t\to0^+$. By Corollary~\ref{cor:exact-gauge}(iii), the level $S_k$ is not
$\sigma$-finite for $\mathcal H^{h_{k-1}}$, whereas $S_{k-1}$ has infinite
but $\sigma$-finite $\mathcal H^{h_{k-1}}$ measure. By
Corollary~\ref{cor:exact-gauge}(ii), the level $S_{k-1}$ is
$\mathcal H^{h_k}$-null, consistent with \eqref{eq:strict-null}.

Moreover, the scale of the critical gauge $h_k$ can be read without reference
to the linear interpolation. For $1/Q_{n+1}\le r\le1/Q_n$ the quantity
$-\log h_k(r)$ lies between $n\log(2k+1)$ and $(n+1)\log(2k+1)$.
By \eqref{eq:rec} both endpoints for $\log(1/r)$ are asymptotic to
$n\log n$, giving
\[
 -\log h_k(r)\sim \log(2k+1)\,
             \frac{\log(1/r)}{\log\log(1/r)}\qquad(r\downarrow0).
\]
In particular, any gauge $g$ whose corresponding logarithmic rate exists and is
strictly greater than $\log(2k+1)$ satisfies $g(1/Q_n)/h_k(1/Q_n)\to0$ and
therefore gives zero measure to $S_k$, while any gauge with a strictly smaller
logarithmic rate satisfies $h_k(t)/g(t)\to0$ and gives a non-$\sigma$-finite
measure on $S_k$.
\end{remark}

\subsection{The three-way splitting at each positive spectral level}

While every subcritical weight $0\le\alpha<1$ assigns the same Hausdorff
dimension $1-\alpha$ to the sublevel set $\{\Lam_\alpha\le c\}$ and to the
exact limit set $\{j^\alpha\|Q_jx\|\to c\}$ (Theorem~\ref{conj:weight}), the
critical weight $\alpha=1$ exhibits a dramatic three-way splitting at each
positive quantized level $c=k/4$ ($k\in\Z_{\ge1}$), governed by the allowed
tail alphabet of the integer recurrence digits $d_j(x)$.

\begin{theorem}[Trichotomy at positive spectral levels, Layer~A in Lean]\label{thm:trichotomy}
Fix an integer $k\ge1$ and let $u_j(x)=Q_jx-\mathrm{round}(Q_jx)$. At the
spectral level $c=k/4$, the sublevel set, the unsigned exact limit set, and
the signed exact limit sets split into three distinct regimes:
\begin{enumerate}[label=(\roman*),leftmargin=2em]
\item \textbf{Limsup level (alphabet of size $2k+1$, gauge base $2k+1$):}
$x\in E_{k/4}=\{\Lam\le k/4\}$ if and only if $u_j(x)\to0$ and eventually
$d_j(x)\in\{-k,-k+1,\dots,k\}$. On $\R/\Z$, both $E_{k/4}$ and the exact
limsup level $\{\Lam=k/4\}$ have infinite, $\sigma$-finite $\mathcal H^{h_k}$
measure for the exact critical gauge $h_k(1/Q_n)=(2k+1)^{-n}$
(Theorem~\ref{thm:gauge} and Corollary~\ref{cor:exact-gauge}).
\item \textbf{Unsigned exact limit set (universal binary alphabet of size $2$, gauge base $2$):}
Define $E_{\mathrm{lim}}(k/4):=\{x\in\R:\lim_{j\to\infty}j\|Q_jx\|=k/4\}$.
Then $x\in E_{\mathrm{lim}}(k/4)$ if and only if $u_j(x)\to0$ and eventually
$|d_j(x)|=k$, i.e.\ $d_j(x)\in\{-k,+k\}$ (exactly two branches per step,
\emph{independently} of $k\ge1$). Consequently, in the interpolated gauge
family $g_\beta(1/Q_n)=\beta^{-n}$, the exact critical gauge for
$E_{\mathrm{lim}}(k/4)$ is $h_{\mathrm{lim}}:=g_2$ ($h_{\mathrm{lim}}(1/Q_n)=2^{-n}$)
for \emph{every} $k\ge1$: on $\R/\Z$, $E_{\mathrm{lim}}(k/4)$ has infinite,
$\sigma$-finite $\mathcal H^{g_2}$ measure (and infinite $\mathcal H^{g_2}$
measure in every nonempty open arc, whereas
$\mathcal H^{h_k}(E_{\mathrm{lim}}(k/4))=0$ since $2<2k+1$), and
$E_{\mathrm{lim}}(c)=\emptyset$ for all $c\notin\frac14\Z_{\ge0}$.
\item \textbf{Signed exact limit sets (single tail branch, single countable $L$-coset):}
For each sign $\pm$, define
$E_{\mathrm{lim}}^\pm(k/4):=\{x\in\R:\lim_{j\to\infty}j\,u_j(x)=\pm k/4\}$.
Then $x\in E_{\mathrm{lim}}^\pm(k/4)$ if and only if $u_j(x)\to0$ and
eventually $d_j(x)=\pm k$. For any $x,y\in E_{\mathrm{lim}}^\pm(k/4)$, their
difference satisfies $x-y\in L=\Z+\Z\frac{\ee+1}{2}$. Hence each signed limit
set is a single countable coset of $L$:
\[
E_{\mathrm{lim}}^+(k/4)=x_k^+ + L,\qquad
E_{\mathrm{lim}}^-(k/4)=x_k^- + L\quad(x_k^-=-x_k^+),
\]
and therefore is countable and has zero measure for every non-atomic Hausdorff
gauge.
\end{enumerate}
(Lean declarations in \texttt{Full\allowbreak Spectrum\allowbreak Discrete.lean},
\texttt{Gauge\allowbreak Circle.lean}, and
\texttt{Trichotomy\allowbreak And\allowbreak Poly\allowbreak Extension.lean}:
the binary digit characterization \texttt{unsigned\_\allowbreak exact\_\allowbreak limit\_\allowbreak iff\_\allowbreak defect\_\allowbreak binary},
the constant digit characterization \texttt{signed\_\allowbreak exact\_\allowbreak limit\_\allowbreak iff\_\allowbreak defect\_\allowbreak const},
the $L$-coset difference theorem \texttt{sub\_\allowbreak mem\_\allowbreak latticeL\_\allowbreak of\_\allowbreak same\_\allowbreak signed\_\allowbreak limit},
countability \texttt{signed\_\allowbreak exact\_\allowbreak limit\_\allowbreak set\_\allowbreak countable},
zero dimension \texttt{dimH\_\allowbreak signed\_\allowbreak exact\_\allowbreak limit\_\allowbreak set\_\allowbreak zero},
the native $\R$ and \texttt{UnitAddCircle} binary-gauge theorems for~(ii) in \texttt{GaugeCircle.lean},
and the full signed-coset collapse for~(iii) via
\texttt{xPlus\_\allowbreak mem\_\allowbreak signed\allowbreak Exact\allowbreak LimitSet},
\texttt{signed\allowbreak Exact\allowbreak LimitSet\_\allowbreak pos\_\allowbreak eq\_\allowbreak coset},
\texttt{signed\allowbreak Exact\allowbreak LimitSet\_\allowbreak neg\_\allowbreak eq\_\allowbreak coset}, and
\texttt{trichotomy\_\allowbreak part\_\allowbreak iii\_\allowbreak complete}.)
\end{theorem}

\begin{proof}
Statement~(i) is immediate from Theorem~\ref{thm:full-spectrum},
Corollary~\ref{thm:sub}, Theorem~\ref{thm:gauge}, and
Corollary~\ref{cor:exact-gauge}.

For~(ii), if $j\|Q_jx\|\to k/4$, then $u_j(x)=O(1/j)\to0$, and
\eqref{eq:digit-asymp} gives $|d_j(x)|\to k$. Since $|d_j(x)|\in\Z_{\ge0}$, we
must have $|d_j(x)|=k$, i.e.\ $d_j(x)\in\{-k,+k\}$, for all sufficiently large
$j$. Conversely, if $u_j(x)\to0$ and eventually $d_j(x)\in\{-k,+k\}$, then
$|d_j(x)|=k$ eventually, so \eqref{eq:digit-asymp} implies $4j\|Q_jx\|\to k$,
that is, $x\in E_{\mathrm{lim}}(k/4)$ (and if $c\notin\frac14\Z_{\ge0}$, no
integer sequence $|d_j|$ can converge to $4c$, so $E_{\mathrm{lim}}(c)=\emptyset$).
We now verify each metric assertion for $E_{\mathrm{lim}}(k/4)$ on $\R/\Z$ under
the binary gauge $g_2(1/Q_n)=2^{-n}$:
\begin{itemize}[leftmargin=1.8em]
\item \emph{$\sigma$-finiteness under $\mathcal H^{g_2}$ and nullity under $\mathcal H^{h_k}$:}
Partition $E_{\mathrm{lim}}(k/4)$ into countably many Borel classes
$\mathcal C_{\mathrm{bin}}(N,m_{N-1},m_N)$ indexed by a starting index $N\ge2$
and starting pair $(M_{N-1},M_N)=(m_{N-1},m_N)\in\Z^2$, requiring
$|u_j|\le(k+1)/j$ and $d_j\in\{-k,+k\}$ for all $j\ge N$. Since at each step
$j\ge N$ there are only $2$ choices for $d_j\in\{-k,+k\}$, at level $n>N$ a
class has at most $2^{n-N}$ possible values of $M_n$, each confining $x$ to an
interval of diameter $2(k+1)/(nQ_n)\le1/Q_n$ for all $n\ge2(k+1)$. Thus the
$\mathcal H_\delta^{g_2}$ covering cost of $\mathcal C_{\mathrm{bin}}(N,m_{N-1},m_N)$
is at most $2^{n-N}g_2(1/Q_n)=2^{-N}<\infty$, proving that
$E_{\mathrm{lim}}(k/4)$ is $\sigma$-finite for $\mathcal H^{g_2}$. For the
larger-base gauge $h_k=g_{2k+1}$, since $2<2k+1$ for every $k\ge1$, the
$\mathcal H_\delta^{h_k}$ covering cost of each class is at most
\[
2^{n-N}h_k(1/Q_n)=(2k+1)^{-N}\left(\frac{2}{2k+1}\right)^{n-N}\longrightarrow0
\qquad(n\to\infty),
\]
so every class is $\mathcal H^{h_k}$-null and
$\mathcal H^{h_k}(E_{\mathrm{lim}}(k/4))=0$. More generally, any gauge $g$ with
$g(1/Q_n)/g_2(1/Q_n)=2^ng(1/Q_n)\to0$ gives
$\mathcal H^g(E_{\mathrm{lim}}(k/4))=0$.
\item \emph{Positivity under $\mathcal H^{g_2}$:}
Fix $N>8k$, set $M_{N-1}=M_N=0$, and restrict the digit sequences in the
\emph{Positivity} step of Theorem~\ref{thm:gauge} to the binary space
$\{-k,+k\}^{\{j\ge N\}}$. By \eqref{eq:telescoping-tail}, every resulting point
$x(\omega)=\lim_{j\to\infty}M_j/Q_j$ satisfies $|Q_jx(\omega)-M_j|<k/j<1/8$ and
has true recurrence digits $d_j(x(\omega))=d_j\in\{-k,+k\}$ for all $j\ge N$, so
$|d_j(x(\omega))|=k$ and $j\|Q_jx(\omega)\|\to k/4$. Hence the compact binary
Cantor image $K_{\mathrm{bin}}:=x(\{-k,+k\}^{\{j\ge N\}})$ lies entirely inside
$E_{\mathrm{lim}}(k/4)\cap[-4k/Q_{N+1},4k/Q_{N+1}]$ (no strict sublevel needs to
be subtracted). Equip $\{-k,+k\}^{\{j\ge N\}}$ with the uniform Bernoulli
product measure assigning weight $1/2$ to $-k$ and $1/2$ to $+k$, and let
$\mu_{\mathrm{bin}}$ be its pushforward to $K_{\mathrm{bin}}$. As in the proof
of Theorem~\ref{thm:gauge}, any interval $U$ of length
$1/Q_{n+1}\le h<1/Q_n$ ($n>N$) determines a unique $M_j$ for each $N\le j<n$ and
at most $3$ values of $M_n$, so $K_{\mathrm{bin}}\cap U$ meets at most $3$ binary
cylinders of depth $n-N$, each of $\mu_{\mathrm{bin}}$-mass $2^{-(n-N)}$. Thus
\[
\mu_{\mathrm{bin}}(U)\le 3\cdot 2^{-(n-N)}=3\cdot 2^{N+1}\cdot 2^{-(n+1)}
\le 3\cdot 2^{N+1}g_2(h),
\]
and the mass distribution principle gives
$\mathcal H^{g_2}(E_{\mathrm{lim}}(k/4))\ge\mathcal H^{g_2}(K_{\mathrm{bin}})\ge(3\cdot 2^{N+1})^{-1}>0$.
\item \emph{Infinitude on every nonempty open arc and non-$\sigma$-finiteness for larger gauges:}
For any $y\in L$, $j\|Q_jy\|\to0$ by Theorem~\ref{thm:class}, so
$\bigl|j\|Q_j(x+y)\|-j\|Q_jx\|\bigr|\le j\|Q_jy\|\to0$; hence
$E_{\mathrm{lim}}(k/4)$ is Borel, Lebesgue-null (since
$E_{\mathrm{lim}}(k/4)\subseteq E_{k/4}$), and invariant under translation by
the dense subgroup $\bar L=L/\Z$ on $\R/\Z$. The Haar-uniqueness argument of
Theorem~\ref{thm:gauge} therefore forces
$\mathcal H^{g_2}(E_{\mathrm{lim}}(k/4))=\infty$ on $\R/\Z$, and the compactness
covering argument of Corollary~\ref{cor:exact-gauge}(i),(iii) with $g_2$ in
place of $h_k$ shows that $\mathcal H^{g_2}(E_{\mathrm{lim}}(k/4)\cap I)=\infty$
for every nonempty open arc $I\subseteq\R/\Z$ and that $E_{\mathrm{lim}}(k/4)$
is not $\sigma$-finite for $\mathcal H^g$ whenever $g_2(t)/g(t)\to0$ as
$t\to0^+$ (in particular for $g_\beta$ with $1<\beta<2$).
\end{itemize}

For~(iii), the signed asymptotic $d_j(x)-4ju_j(x)\to0$ in
\eqref{eq:digit-asymp} shows that $ju_j(x)\to\pm k/4$ if and only if
$u_j(x)\to0$ and the integer $d_j(x)$ is eventually constant and equal to
$\pm k$. The existence of a witness $x_k^+\in E_{\mathrm{lim}}^+(k/4)$ (and
$x_k^-=-x_k^+\in E_{\mathrm{lim}}^-(k/4)$) is given by the constant-digit
construction $d_j\equiv k$ in Theorem~\ref{thm:gauge}. If
$x,y\in E_{\mathrm{lim}}^\pm(k/4)$, then $j(u_j(x)-u_j(y))\to0$; since
$Q_j(x-y)=\bigl(M_j(x)-M_j(y)\bigr)+\bigl(u_j(x)-u_j(y)\bigr)$ with
$M_j(x)-M_j(y)\in\Z$, we have $\|Q_j(x-y)\|\le|u_j(x)-u_j(y)|$, hence
$j\|Q_j(x-y)\|\to0$. By Theorem~\ref{thm:class}, $x-y\in L$. Conversely, for
any $z\in L$, $j\|Q_jz\|\to0$ gives $u_j(z)=o(1/j)\to0$; for
$x\in E_{\mathrm{lim}}^\pm(k/4)$, since $u_j(x)=O(1/j)\to0$, we have
$|u_j(x)+u_j(z)|<1/2$ for all large $j$, so $Q_j(x+z)$ rounds to
$M_j(x)+M_j(z)$ and $u_j(x+z)=u_j(x)+u_j(z)$, whence
$j\,u_j(x+z)=j\,u_j(x)+j\,u_j(z)\to\pm k/4$. Thus $x+z\in E_{\mathrm{lim}}^\pm(k/4)$
and $E_{\mathrm{lim}}^\pm(k/4)=x_k^\pm+L$.
\end{proof}

\paragraph{Lean status for Theorem~\ref{thm:trichotomy}.}
All three parts~(i), (ii), and~(iii) of Theorem~\ref{thm:trichotomy} are now
kernel-checked at Layer~A (using only the three Lean base axioms
\texttt{propext}, \texttt{Classical.choice}, and \texttt{Quot.sound}, with zero
\texttt{sorryAx}). Specifically, the digit characterizations, signed-limit
$L$-difference, countability, and dimension-zero lemmas are proved in
\texttt{EpiFormal/\allowbreak Full\allowbreak Spectrum\allowbreak Discrete.lean};
all metric gauge statements for the limsup regime~(i) and the binary exact-limit
regime $E_{\mathrm{lim}}(k/4)$ in~(ii)---binary Cantor-tree positivity, global
and nonempty-open-subset $\mathcal H^{g_2}$-infinitude, $\sigma$-finite
covering, $\mathcal H^{h_k}$-nullity, and the $g_\beta$ threshold dichotomy
($\beta>2$ vs.\ $1<\beta<2$)---are proved natively on both $\R$ and
\texttt{UnitAddCircle} ($\R/\Z$) in
\texttt{EpiFormal/\allowbreak GaugeCircle.lean} (building on
\texttt{GaugeNull.lean} and \texttt{GaugeExact.lean}); and the full
signed-coset collapse $E_{\mathrm{lim}}^\pm(k/4)=x_k^\pm+L$ in part~(iii) is
proved in
\texttt{EpiFormal/\allowbreak TrichotomyAndPolyExtension.lean} across eight
audited declarations (\texttt{xPlus\_\allowbreak mem\_\allowbreak signedExactLimitSet},
\texttt{neg\_\allowbreak xPlus\_\allowbreak mem\_\allowbreak signedExactLimitSet},
\texttt{add\_\allowbreak latticeL\_\allowbreak mem\_\allowbreak signedExactLimitSet},
\texttt{signedExactLimitSet\_\allowbreak pos\_\allowbreak eq\_\allowbreak coset},
\texttt{signedExactLimitSet\_\allowbreak neg\_\allowbreak eq\_\allowbreak coset},
\texttt{signedExactLimitSet\_\allowbreak zero\_\allowbreak eq\_\allowbreak latticeL},
\texttt{signedExactLimitSet\_\allowbreak empty\_\allowbreak of\_\allowbreak not\_\allowbreak quarter\_\allowbreak int},
and \texttt{trichotomy\_\allowbreak part\_\allowbreak iii\_\allowbreak complete}).

\subsection{Finite-state constraints on defect tails}\label{sec:finite-state-tail}

Let $G$ be a finite directed graph with edge labels in
$\{-k,\ldots,k\}$, with distinct labels on the outgoing edges at each
vertex. Mark the edges labelled $k$ or $-k$. Let $S_k(G)\subset S_k$ be
the points whose eventual defect sequence labels an infinite path in $G$.
A strongly connected component is called marked if it contains an internal
marked edge. Write $\lambda_C$ for the Perron root of the adjacency matrix
of a marked component $C$, counting edge multiplicities.

\begin{theorem}[Finite-state defect tails]\label{thm:finite-state-tail}
If $G$ has no marked component, then $S_k(G)=\varnothing$.
Otherwise put $\lambda_*=\max_C\lambda_C$, where $C$ ranges over
the marked components. If $\lambda_*=1$, then $S_k(G)$ is countable.
If $\lambda_*>1$, it satisfies the three measure conclusions of
Corollary~\ref{cor:liminf-gauges}, with $S_k(G)$ in place of the named
set and with $b=\lambda_*$. In the zero case the null conclusion refers
to $S_k(G)$.
\end{theorem}

\begin{proof}
First consider a strongly connected marked graph with Perron root
$\lambda>1$. Take a positive right Perron vector $w$. The sum of
terminal weights over paths of length $m$ from vertex $v$ equals
$\lambda^m w_v$. Since the coordinates of $w$ are bounded above and
away from zero, the path count is comparable to $\lambda^m$, uniformly
over the initial vertex. The fixed initial-pair covering argument
therefore gives at most $C\lambda^{n-N}$ intervals of diameter $r_n$
for each fixed starting position, pair and vertex. This gives the
zero and finite-cover conclusions when the normalized lower limit is
zero or finite, respectively.

To obtain positive critical measure, give an edge $v\to v'$ probability
$w_{v'}/(\lambda w_v)$. The path-cylinder mass telescopes to
$w_{v_m}/(\lambda^m w_{v_0})\le C\lambda^{-m}$. Fix one marked edge.
From every vertex there is a path crossing it in at most some fixed
number $L_0$ of steps, with probability at least $p>0$. Conditioning
successively in blocks of length $L_0$ bounds the probability of
avoiding the edge for $m$ blocks by $(1-p)^m$. Applying this after each
starting time shows that almost every path crosses a marked edge
infinitely often.

Feed the edge labels into the defect recurrence, with the two
initial integers zero and with starting index $N>8k$ sufficiently
large. The construction in the proof of Theorem~\ref{thm:gauge} gives
a point $x=\lim M_j/Q_j$ with
$|Q_jx-M_j|\le 4kQ_j/Q_{j+1}<1/8$ for $j\ge N$.
Thus the $M_j$ are its nearest integers and the prescribed labels
are its defects. Infinitely many marked labels give $\Lam(x)=k/4$
by the defect asymptotic. The pushforward of the path probability
therefore gives full mass to $S_k(G)$.

If a set of diameter $r$ satisfies $r_{n+1}\le r<r_n$, the difference
between two of its coded values of $M_j$ is less than
$1/4+Q_jr$. For $j<n$ with $j\ge N$ this is less than one, and
at $j=n$ there are at most three possible adjacent integers.
Consequently the set meets at most three label cylinders of length
$n-N$. From a prescribed start vertex each label word specifies at
most one path. Thus its mass is at most $3C\lambda^{-(n-N)}$, which is
bounded by a fixed multiple of $h_\lambda(r)$ for the interpolated gauge
$h_\lambda(r_n)=\lambda^{-n}$. Singleton mass is zero by the same
cylinder estimate. The construction is in a real interval of length
less than $1/2$ if $N$ is large, so projection to the circle is
isometric there. The mass-distribution bound gives
$\mathcal H^{h_\lambda}(S_k(G))>0$.

For Borel measurability, intersect with the Borel level $S_k$. For each
start and vertex, the condition that every finite prefix is admissible
is a countable intersection of Borel conditions on the rounding
functions. An infinite path exists by finite branching. Take the
countable union over starts and the finite union over vertices.
For $y\in L$, the defect asymptotic gives $d_j(y)=0$ eventually.
Because the nearest-integer errors of $x\in S_k$ and of $y$ tend to
zero, their nearest integers add eventually, so $d_j(x+y)=d_j(x)$
eventually. Hence $S_k(G)$ is invariant under the dense group $L$
and is Lebesgue-null. If its positive critical measure were finite,
its normalized restriction would be invariant under this dense group
and, by continuity against continuous test functions, under the
entire circle. It would then be Haar measure, a contradiction.
For a nonempty open set $U$, the family $U+\ell$, $\ell\in L\bmod 1$,
covers the circle by density and has a finite subcover by compactness.
Invariance under those translations would make the whole level have
finite measure if $S_k(G)\cap U$ had finite measure. This proves local
infinitude. The comparison of adjacent scales in
Corollary~\ref{cor:liminf-gauges} now gives the three conclusions for
every gauge. The middle case combines the countable covering by
fixed initial-pair classes with local infinitude on the full level.

For a general finite graph, each infinite path eventually remains
in one strongly connected component, since the component graph is
acyclic. If its point is in $S_k$, its defects have absolute value
$k$ infinitely often. Thus the terminal component is marked, and
\[
 S_k(G)=\bigcup_{C\ {\rm marked}} S_k(C).
\]
For a marked component with $\lambda_C=1$, each vertex has exactly
one internal outgoing edge. Indeed, two outgoing edges at one vertex
give two return loops with distinct first edges. Traversing them in
the two orders gives two distinct loops of equal length. Concatenating
those equal-length loops contradicts the bounded path count at
Perron root one. A marked component is therefore a directed cycle
in this case. Each fixed start index, vertex and integer pair
then specifies a unique defect tail and at most one real point:
the recurrence determines all later $M_j$, and $M_j/Q_j\to x$.
This gives countability.

Finally, when $1<\lambda_C<\lambda_*$ and
$\rho_{\lambda_*}(g)<\infty$, a bounded subsequence at base $\lambda_*$
has normalized cost tending to zero at base $\lambda_C$ by the
factor $(\lambda_C/\lambda_*)^n$. Such components are null.
Countable components are also null. At least one component attains
$\lambda_*$, giving local infinitude or local non-$\sigma$-finiteness
in the corresponding regimes. The finite union of the component
covers gives the asserted global finite-cover conclusion.
\end{proof}

\begin{remark}[Two graph examples]
Forbidding successive marked labels gives the adjacency matrix
\[
 \begin{pmatrix}2k-1&2\\2k-1&0\end{pmatrix}.
\]
For $k\ge2$ its characteristic polynomial has opposite signs at $2$
and $2k+1$, so its Perron root is strictly between those values.
For a reducible example, take a disjoint unmarked vertex with the
$2k-1$ unmarked loops and a marked vertex with the two loops labelled
$\pm k$. For $k\ge2$ the larger adjacency root comes from the unmarked
vertex, while $S_k(G)$ has critical base $2$. Thus the root of the
whole adjacency matrix need not determine the gauge of an exact level.
\end{remark}

For comparison, the graph-directed similarity results of Mauldin
and Williams \cite[Theorems~3 and~4]{MW} use a spectral radius of weighted
similarity matrices and study the full construction. Here the scales
are supplied by $Q_n$, and the insistence on infinitely many marked
edges arises from exact-level membership. The local infinitude uses
the arithmetic translation group of the level.

\paragraph{Lean status.}
All three regimes of Theorem~\ref{thm:finite-state-tail} are now kernel-checked at Layer~A in
\texttt{EpiFormal/\allowbreak Finite\allowbreak State\allowbreak Complete.lean} ($12$ audited declarations
ending in \texttt{thm\_\allowbreak finite\_\allowbreak state\_\allowbreak tail\_\allowbreak complete} on $\R$
and \texttt{thm\_\allowbreak finite\_\allowbreak state\_\allowbreak tail\_\allowbreak circle\_\allowbreak complete}
on \texttt{UnitAddCircle}, building on \texttt{Finite\allowbreak State\allowbreak Gauge.lean},
\texttt{Finite\allowbreak State\allowbreak Tail.lean}, and
\texttt{Finite\allowbreak State\allowbreak Positive.lean}, with base axioms only and zero \texttt{sorryAx}).
Specifically, under the deterministic outgoing-label hypothesis of this subsection, the formalization proves
compactness and subset embedding of the deterministic path tree (\texttt{graphTree}),
direct finite-sum Frostman positivity (\texttt{mkMetric\_\allowbreak graphTree\_\allowbreak pos}),
unmarked-tail nullity via backward-reachability stabilization at $L=|V|$ (\texttt{mkMetric\_\allowbreak unmarkedTail\_\allowbreak null}),
$L$-translation invariance of exact graph tails, and the arbitrary-gauge $\liminf$ comparison;
the positive right Perron--Frobenius eigenvector $(w_v,\lambda_C)$ on each marked strongly connected component is taken as an explicit hypothesis rather than derived via Lean matrix-eigenvalue machinery.

\subsection{General polynomial-growth denominator recurrences}\label{sec:poly-rec}

The quantization, exact gauge, three-way positive-level splitting, and
subcritical dimension phase transition extend from the Euler recurrence
\eqref{eq:rec} to arbitrary polynomial-growth three-term recurrences. Throughout
this subsection, fix a degree $r\in\Z_{\ge1}$, a leading coefficient
$c\in\Z_{\ge1}$, and an integer shift $d\in\Z$. Set $A_j:=cj^r+d$, choose an
integer $J_0\ge1$ such that $A_{J_0}=cJ_0^r+d\ge2$ (so $A_j$ is strictly
increasing and $A_j\ge2$ for all $j\ge J_0$), and let $(Q_j)_{j\ge0}$ be any
sequence of positive integers satisfying
\begin{equation}\label{eq:poly-rec}
Q_{j+1}=(cj^r+d)Q_j+Q_{j-1}\qquad\text{for all }j\ge J_0.
\end{equation}
For $x\in\R$ and $\alpha\ge0$, define the weighted approximation exponents
\[
\Lam^{(r,c,d)}(x):=\limsup_{j\to\infty}j^r\|Q_jx\|\in[0,\infty],\qquad
\Lam_\alpha^{(r,c,d)}(x):=\limsup_{j\to\infty}j^\alpha\|Q_jx\|\in[0,\infty],
\]
so that $\Lam^{(r,c,d)}=\Lam_r^{(r,c,d)}$. Fix the Borel rounding rule
$M_j(x):=\lfloor Q_jx+1/2\rfloor\in\Z$ and signed remainder
$u_j(x):=Q_jx-M_j(x)\in[-1/2,1/2)$ (so $|u_j(x)|=\|Q_jx\|$), and for $j\ge J_0$
define the integer recurrence defect (digit)
\begin{equation}\label{eq:poly-defect-id}
d_j(x):=M_{j+1}(x)-A_jM_j(x)-M_{j-1}(x)=A_ju_j(x)+u_{j-1}(x)-u_{j+1}(x)\in\Z.
\end{equation}

\begin{lemma}[Basic growth and non-asymptotic factorial bound, Layer~A in Lean]\label{lem:poly-growth}
Put $J_1:=\max(J_0+1,\,2|d|+1)$. Then:
\begin{enumerate}[label=(\roman*),leftmargin=2em]
\item For all $j\ge J_0$, $Q_{j+1}>A_jQ_j\ge2Q_j$.
\item For all $j\ge J_0+1$, $0<Q_{j-1}<\frac12Q_j<Q_j$, and hence
\[
cj^r+d<\frac{Q_{j+1}}{Q_j}<cj^r+d+1\le(c+|d|+1)j^r.
\]
\item For all $j\ge J_1$, $A_j\ge\frac{c}{2}j^r$, and for every $m\ge J_1$,
\[
Q_m\ge C_{J_1}\left(\frac{c}{2^{r+1}}\right)^m(m!)^r,\qquad
C_{J_1}:=Q_{J_1}\left(\frac{c}{2}\right)^{-J_1}\bigl((J_1-1)!\bigr)^{-r}>0.
\]
\item As $n\to\infty$, $\log Q_n=rn\log n+(\log c-r)n+O(\log n)$.
\end{enumerate}
\end{lemma}

\begin{proof}
Since $Q_{j-1}\ge1$ and $A_j\ge2$ for $j\ge J_0$, \eqref{eq:poly-rec} yields
$Q_{j+1}>A_jQ_j\ge2Q_j$, proving~(i). If $j\ge J_0+1$, then $j-1\ge J_0$, so~(i)
gives $Q_j>2Q_{j-1}$, i.e.\ $0<Q_{j-1}/Q_j<1/2<1$; adding $A_j$ gives
$cj^r+d<Q_{j+1}/Q_j<cj^r+d+1\le(c+|d|+1)j^r$ (using $j^r\ge1$), proving~(ii). For
$j\ge J_1$, we have $|d|<j/2\le\frac12 cj^r$ and therefore $A_j\ge\frac{c}{2}j^r$.
Multiplying $Q_{j+1}>A_jQ_j\ge\frac{c}{2}j^r Q_j$ over $j=J_1,\dots,m-1$ gives
\[
Q_m\ge Q_{J_1}\left(\frac{c}{2}\right)^{m-J_1}
\left(\frac{(m-1)!}{(J_1-1)!}\right)^r
=C_{J_1}\left(\frac{c}{2}\right)^m\frac{(m!)^r}{m^r}
\ge C_{J_1}\left(\frac{c}{2^{r+1}}\right)^m(m!)^r,
\]
since $m\le2^m$ for all $m\ge1$ implies $m^r\le(2^r)^m$, proving~(iii). Finally,
by~(ii), $\log(Q_{j+1}/Q_j)=r\log j+\log c+O(j^{-r})$ for $j\ge J_1$; summing
over $j=J_1,\dots,n-1$ and using $\sum_{j=1}^{n-1}\log j=n\log n-n+O(\log n)$
together with $\sum_{j=J_1}^{n-1}j^{-r}=O(\log n)$ (for $r=1$) or $O(1)$ (for
$r\ge2$) gives~(iv).
\end{proof}

When no pre-existing convergent sequence of $\ee$ is present, the role of the
group $L=\Z+\Z\frac{\ee+1}{2}$ is played by an intrinsic rank-$2$ lattice built
directly from the homogeneous recurrence $X_{j+1}=A_jX_j+X_{j-1}$ ($j\ge J_0$).
Let $(U_j)_{j\ge J_0-1}$ and $(V_j)_{j\ge J_0-1}$ be the two fundamental integer
solutions specified by
\[
U_{J_0-1}=1,\quad U_{J_0}=0;\qquad V_{J_0-1}=0,\quad V_{J_0}=1.
\]

\begin{lemma}[Intrinsic zero-level lattice and Casorati rigidity, Layer~A in Lean]\label{lem:poly-lattice}
\begin{enumerate}[label=(\roman*),leftmargin=2em]
\item For all $j\ge J_0-1$, $Q_j=Q_{J_0-1}U_j+Q_{J_0}V_j$ and
$U_jV_{j+1}-U_{j+1}V_j=(-1)^{j-J_0+1}$.
\item The limits $\alpha_U:=\lim_{j\to\infty}U_j/Q_j$ and
$\alpha_V:=\lim_{j\to\infty}V_j/Q_j$ exist, satisfy
$Q_{J_0-1}\alpha_U+Q_{J_0}\alpha_V=1$, and obey
\[
\left|\alpha_U-\frac{U_j}{Q_j}\right|<\frac{Q_{J_0}}{Q_jQ_{j+1}},\qquad
\left|\alpha_V-\frac{V_j}{Q_j}\right|<\frac{Q_{J_0-1}}{Q_jQ_{j+1}}
\qquad\text{for all }j\ge J_0.
\]
\item Set $g:=\gcd(Q_{J_0-1},Q_{J_0})\ge1$, write $Q_{J_0-1}=gq_0$ and
$Q_{J_0}=gq_1$ with $\gcd(q_0,q_1)=1$, pick integers $p_0,p_1\in\Z$ with
$p_0q_1-p_1q_0=1$, and put $\tilde Q_j:=Q_j/g=q_0U_j+q_1V_j$ and
$\tilde P_j:=p_0U_j+p_1V_j$. Then
\[
\theta:=\lim_{j\to\infty}\frac{\tilde P_j}{\tilde Q_j}=g(p_0\alpha_U+p_1\alpha_V)\in\R\setminus\Q,
\]
and the intrinsic rank-$2$ lattice
\[
L^{(r,c,d)}:=\Z\alpha_U+\Z\alpha_V=\frac{1}{g}\bigl(\Z+\Z\theta\bigr)
\]
is independent of the choice of $(p_0,p_1)$, contains $\Z$, and projects to a
dense countable subgroup $\bar L^{(r,c,d)}:=L^{(r,c,d)}/\Z$ of $\R/\Z$.
\item For $x\in\R$, $\Lam^{(r,c,d)}(x)=0$ if and only if $u_j(x)\to0$ and
eventually $d_j(x)=0$, if and only if $x\in L^{(r,c,d)}$. Moreover, every
$x\in L^{(r,c,d)}$ satisfies $\|Q_jx\|=O(1/Q_{j+1})$, whence
$\Lam_\alpha^{(r,c,d)}(x)=0$ for all $\alpha\ge0$.
\end{enumerate}
\end{lemma}

\begin{proof}
Part~(i) holds at $j=J_0-1$ by inspection and propagates to all $j\ge J_0-1$
because $(U_j),(V_j),(Q_j)$ satisfy the same recurrence \eqref{eq:poly-rec} whose
transfer matrices $\bigl(\begin{smallmatrix}A_j&1\\1&0\end{smallmatrix}\bigr)$
have determinant $-1$. Expanding $Q_j=Q_{J_0-1}U_j+Q_{J_0}V_j$ in the adjacent
ratio differences yields, for $j\ge J_0$,
\[
\frac{U_{j+1}}{Q_{j+1}}-\frac{U_j}{Q_j}
=\frac{(-1)^{j-J_0}Q_{J_0}}{Q_jQ_{j+1}},\qquad
\frac{V_{j+1}}{Q_{j+1}}-\frac{V_j}{Q_j}
=\frac{(-1)^{j-J_0+1}Q_{J_0-1}}{Q_jQ_{j+1}}.
\]
By Lemma~\ref{lem:poly-growth}(i), $Q_{j+2}>2Q_{j+1}>Q_j$ for all $j\ge J_0$,
so $1/(Q_jQ_{j+1})$ decreases strictly to $0$; the alternating series test
gives convergence and bounds the tail strictly by the magnitude of the first
term, proving~(ii) (and $Q_{J_0-1}\alpha_U+Q_{J_0}\alpha_V=1$ follows by taking
$j\to\infty$ in $Q_{J_0-1}(U_j/Q_j)+Q_{J_0}(V_j/Q_j)=1$).

For~(iii), $\bigl(\begin{smallmatrix}p_0&p_1\\q_0&q_1\end{smallmatrix}\bigr)\in\mathrm{SL}_2(\Z)$
gives $\tilde P_j\tilde Q_{j+1}-\tilde P_{j+1}\tilde Q_j=(-1)^{j-J_0+1}$, so
$\theta=\lim_{j\to\infty}\tilde P_j/\tilde Q_j=\lim_{j\to\infty}g(p_0U_j+p_1V_j)/Q_j=g(p_0\alpha_U+p_1\alpha_V)$
satisfies $0<|\theta-\tilde P_j/\tilde Q_j|<(\tilde Q_j\tilde Q_{j+1})^{-1}$ for
all $j\ge J_0$. If $\theta=p/q\in\Q$ ($q\ge1$), then
$0<|q\tilde P_j-p\tilde Q_j|<q/\tilde Q_{j+1}<1$ for large $j$, contradicting
$|q\tilde P_j-p\tilde Q_j|\in\Z_{\ge1}$; hence $\theta\notin\Q$. Since
$\bigl(\begin{smallmatrix}p_0&p_1\\q_0&q_1\end{smallmatrix}\bigr)\in\mathrm{SL}_2(\Z)$
is unimodular and $q_0\alpha_U+q_1\alpha_V=1/g$, we have
$\Z\alpha_U+\Z\alpha_V=\Z(q_0\alpha_U+q_1\alpha_V)+\Z(p_0\alpha_U+p_1\alpha_V)
=\Z\frac{1}{g}+\Z\frac{\theta}{g}=\frac{1}{g}(\Z+\Z\theta)$, which is thus
independent of $(p_0,p_1)$ and contains $1=g\cdot\frac{1}{g}$. Because
$\theta/g\notin\Q$, Kronecker's theorem shows that $\bar L^{(r,c,d)}$ is dense
in $\R/\Z$.

For~(iv), if $\Lam^{(r,c,d)}(x)=0$, then $j^r|u_j(x)|\to0$, so
$|d_j(x)|\le cj^r|u_j(x)|+|d||u_j(x)|+|u_{j-1}(x)|+|u_{j+1}(x)|\to0$; since
$d_j(x)\in\Z$, eventually $d_j(x)=0$. Conversely, if $u_j(x)\to0$ and $d_j(x)=0$
for all $j\ge N\ge J_0+1$, then by~(i)
$\det\bigl(\begin{smallmatrix}U_{N-1}&V_{N-1}\\U_N&V_N\end{smallmatrix}\bigr)=\pm1$,
so $(M_{N-1},M_N)=a(U_{N-1},U_N)+b(V_{N-1},V_N)$ for unique $a,b\in\Z$, and by
induction $M_j(x)=aU_j+bV_j$ for all $j\ge N-1$; dividing by $Q_j$ and using
$u_j(x)\to0$ yields $x=a\alpha_U+b\alpha_V\in L^{(r,c,d)}$. Finally, if
$x=a\alpha_U+b\alpha_V\in L^{(r,c,d)}$, then~(ii) gives
$\|Q_jx\|\le|Q_jx-(aU_j+bV_j)|\le(|a|Q_{J_0}+|b|Q_{J_0-1})/Q_{j+1}$ for
$j\ge J_0$, and Lemma~\ref{lem:poly-growth}(iii) yields $j^\alpha/Q_{j+1}\to0$
for every $\alpha\ge0$.
\end{proof}

\begin{theorem}[General $(cj^r+d)$ recurrences: complete spectrum discretization,
Layer~A in Lean]\label{thm:general-rec}
For any $r,c\in\Z_{\ge1}$, $d\in\Z$, and $(Q_j)_{j\ge0}$ satisfying
\eqref{eq:poly-rec}:
\begin{enumerate}[label=(\roman*),leftmargin=2em]
\item Whenever $u_j(x)\to0$ (in particular, whenever $\Lam^{(r,c,d)}(x)<\infty$),
the digits $d_j(x)$ satisfy
\[
\lim_{j\to\infty}\bigl(d_j(x)-cj^ru_j(x)\bigr)=0,\qquad
\lim_{j\to\infty}\bigl(|d_j(x)|-cj^r\|Q_jx\|\bigr)=0.
\]
\item The full spectrum of $\Lam^{(r,c,d)}$ is completely discrete with step $1/c$:
\[
\Lam^{(r,c,d)}(\R)=\frac{1}{c}\Z_{\ge0}\cup\{\infty\}
=\left\{0,\frac{1}{c},\frac{2}{c},\frac{3}{c},\dots\right\}\cup\{\infty\},
\]
so every open interval $\bigl(\frac{k}{c},\frac{k+1}{c}\bigr)$ ($k\in\Z_{\ge0}$)
is an empty spectral gap, and for every $t\in\bigl[\frac{k}{c},\frac{k+1}{c}\bigr)$
the sublevel set $E_t^{(r,c,d)}:=\{x\in\R:\Lam^{(r,c,d)}(x)\le t\}$ equals
$E_{k/c}^{(r,c,d)}=\{x\in\R:u_j(x)\to0\text{ and }\limsup_{j\to\infty}|d_j(x)|\le k\}$
(with $E_t^{(r,c,d)}=L^{(r,c,d)}$ for $t<1/c$).
\end{enumerate}
(Lean: for all $r\ge1$, proved at Layer~A in
\texttt{EpiFormal/\allowbreak TrichotomyAndPolyExtension.lean} via
\texttt{PolyRecData.\allowbreak poly\_\allowbreak signed\_\allowbreak defect\_\allowbreak approx},
\texttt{poly\_\allowbreak abs\_\allowbreak defect\_\allowbreak approx},
\texttt{poly\_\allowbreak spectral\_\allowbreak gap\_\allowbreak open\_\allowbreak interval},
\texttt{poly\_\allowbreak levelSet\_\allowbreak eq\_\allowbreak on\_\allowbreak gap},
\texttt{poly\_\allowbreak levelSet\_\allowbreak below\_\allowbreak eq\_\allowbreak lattice}, and
\texttt{poly\_\allowbreak range\_\allowbreak apprExponent}; the $r=1$
specialization also appears in
\texttt{EpiFormal/\allowbreak Full\allowbreak Spectrum\allowbreak Discrete.lean}.)
\end{theorem}

\begin{proof}
Rearranging \eqref{eq:poly-defect-id} gives
$d_j(x)-cj^ru_j(x)=du_j(x)+u_{j-1}(x)-u_{j+1}(x)\to0$ whenever $u_j(x)\to0$
(which holds whenever $\Lam^{(r,c,d)}(x)<\infty$), and the reverse triangle
inequality $\bigl||d_j(x)|-cj^r|u_j(x)|\bigr|\le|d_j(x)-cj^ru_j(x)|$ proves~(i).
Since $|d_j(x)|\in\Z_{\ge0}$, whenever $\Lam^{(r,c,d)}(x)<\infty$ we obtain
$c\Lam^{(r,c,d)}(x)=\limsup_{j\to\infty}|d_j(x)|\in\Z_{\ge0}$, excluding all
open intervals $(k/c,(k+1)/c)$ and forcing $E_t^{(r,c,d)}=E_{k/c}^{(r,c,d)}$ on
$[k/c,(k+1)/c)$. Conversely, $0\in L^{(r,c,d)}$ realizes $0$,
Theorem~\ref{thm:poly-phase-transition}(i) below realizes $\infty$ on a set of
full Lebesgue measure, and for each $k\ge1$ the constant-digit construction
$d_j\equiv k$ in Theorem~\ref{thm:poly-gauge-trichotomy} below produces a point
$x_k^+$ with $|u_j(x_k^+)|<4k/A_j=O(j^{-r})\to0$ and $d_j(x_k^+)=k$ for large
$j$, giving $\Lam^{(r,c,d)}(x_k^+)=k/c$ by~(i). The $t<1/c$ identity follows from
Lemma~\ref{lem:poly-lattice}(iv).
\end{proof}

For each integer $k\ge1$, define the continuous nondecreasing critical gauge
$h_k^{(r,c,d)}:[0,\infty)\to[0,\infty)$ by $h_k^{(r,c,d)}(0)=0$,
$h_k^{(r,c,d)}(1/Q_n)=(2k+1)^{-n}$ for all $n\ge J_0$, linear interpolation on
each $[1/Q_{n+1},1/Q_n]$ ($n\ge J_0$), and constant value $(2k+1)^{-J_0}$ for
$t\ge1/Q_{J_0}$ (well-defined because $Q_n$ is strictly increasing for
$n\ge J_0$ by Lemma~\ref{lem:poly-growth}(i)). Similarly define the universal
binary gauge $g_2^{(r,c,d)}(1/Q_n):=2^{-n}$ ($n\ge J_0$).

\begin{theorem}[Critical gauges and positive-level trichotomy for general $(cj^r+d)$, Layer~A in Lean]\label{thm:poly-gauge-trichotomy}
Fix an integer $k\ge1$. At the positive spectral level $k/c$:
\begin{enumerate}[label=(\roman*),leftmargin=2em]
\item \textbf{Limsup level (alphabet size $2k+1$, gauge base $2k+1$):}
$x\in E_{k/c}^{(r,c,d)}$ if and only if $u_j(x)\to0$ and eventually
$d_j(x)\in\{-k,\dots,k\}$. On $\R/\Z$, the exact level
$S_k^{(r,c,d)}:=\{x\in\R/\Z:\Lam^{(r,c,d)}(x)=k/c\}$ has infinite but
$\sigma$-finite $\mathcal H^{h_k^{(r,c,d)}}$ measure (with infinite measure in
every nonempty open arc), whereas $\{\Lam^{(r,c,d)}<k/c\}$ is
$\mathcal H^{h_k^{(r,c,d)}}$-null.
\item \textbf{Unsigned exact limit set (universal binary alphabet $\{\pm k\}$, gauge base $2$):}
$E_{\mathrm{lim}}^{(r,c,d)}(k/c):=\{x\in\R:\lim_{j\to\infty}j^r\|Q_jx\|=k/c\}$
consists precisely of those $x$ with $u_j(x)\to0$ and eventually
$d_j(x)\in\{-k,+k\}$ (two branches per step, independently of $k\ge1$). On
$\R/\Z$, $E_{\mathrm{lim}}^{(r,c,d)}(k/c)$ has infinite, $\sigma$-finite
$\mathcal H^{g_2^{(r,c,d)}}$ measure (infinite on every nonempty open arc),
while $\mathcal H^{h_k^{(r,c,d)}}(E_{\mathrm{lim}}^{(r,c,d)}(k/c))=0$. For
$t\notin\frac1c\Z_{\ge0}$, $E_{\mathrm{lim}}^{(r,c,d)}(t)=\emptyset$.
\item \textbf{Signed exact limit sets (single tail branch, coset collapse):}
$E_{\mathrm{lim}}^{(r,c,d),\pm}(k/c):=\{x\in\R:\lim_{j\to\infty}j^ru_j(x)=\pm k/c\}$
consists precisely of those $x$ with $u_j(x)\to0$ and eventually $d_j(x)=\pm k$.
With $x_k^+$ built from $d_j\equiv+k$ ($j\ge N$) and $x_k^-:=-x_k^+$, each signed
exact limit set collapses to a single countable coset of the intrinsic lattice
$L^{(r,c,d)}$ of Lemma~\ref{lem:poly-lattice}:
\[
E_{\mathrm{lim}}^{(r,c,d),+}(k/c)=x_k^+ + L^{(r,c,d)},\qquad
E_{\mathrm{lim}}^{(r,c,d),-}(k/c)=x_k^- + L^{(r,c,d)},
\]
and is therefore null for every non-atomic Hausdorff gauge.
\end{enumerate}
\end{theorem}

\begin{proof}
The digit characterizations follow directly from Theorem~\ref{thm:general-rec}.
For the upper bounds and $\sigma$-finiteness on $[0,1]$ (hence on $\R/\Z$ since
the quotient map is $1$-Lipschitz), partition $E_{k/c}^{(r,c,d)}\cap[0,1]$ into
countably many Borel classes $\mathcal C(N,m_{N-1},m_N)$ ($N\ge J_0+1$,
$m_{N-1},m_N\in\Z$) requiring $M_{N-1}=m_{N-1}$, $M_N=m_N$,
$|u_j|\le(k+1)/(cj^r)$, and $d_j\in\Sigma$ for all $j\ge N$, where
$\Sigma=\{-k,\dots,k\}$ ($|\Sigma|=2k+1$) for $E_{k/c}^{(r,c,d)}$,
$\Sigma=\{-(k-1),\dots,k-1\}$ ($|\Sigma|=2k-1$) for $\{\Lam^{(r,c,d)}<k/c\}$,
and $\Sigma=\{-k,+k\}$ ($|\Sigma|=2$) for $E_{\mathrm{lim}}^{(r,c,d)}(k/c)$.
At level $n>N$, the recurrence $M_{j+1}=A_jM_j+M_{j-1}+d_j$ yields at most
$|\Sigma|^{n-N}$ values of $M_n$, each confining $x$ to an interval of diameter
at most $1/Q_n$ (using $|u_n|\le1/2$). Since $n>J_0$, covering with these
intervals gives $\mathcal H_\delta^{h_k^{(r,c,d)}}$ cost at most
$(2k+1)^{-N}\bigl(|\Sigma|/(2k+1)\bigr)^{n-N}$ (bounded by $(2k+1)^{-N}<\infty$
when $|\Sigma|=2k+1$, and $\to0$ when $|\Sigma|<2k+1$), and when $\Sigma=\{-k,+k\}$
gives $\mathcal H_\delta^{g_2^{(r,c,d)}}$ cost at most $2^{n-N}2^{-n}=2^{-N}<\infty$.

For the lower bounds, choose an integer $N\ge J_0+1$ such that
$A_N=cN^r+d>32k$. Since $A_j\ge A_N>32k\ge32>6$ and $Q_N\ge1$, Lemma~\ref{lem:poly-growth}(i)
gives $Q_{N+1}>32kQ_N\ge32k$ and $Q_{j+1}>6Q_j$ for all $j\ge N$. Set
$M_{N-1}=M_N=0$ and, for $\Sigma\in\{\{-k,\dots,k\},\{-k,+k\}\}$ and
$\omega=(d_N,d_{N+1},\dots)\in\Sigma^{\{j\ge N\}}$, define $M_{j+1}=A_jM_j+M_{j-1}+d_j$
($j\ge N$). Then $W_j:=M_{j+1}Q_j-M_jQ_{j+1}$ satisfies $W_{N-1}=0$ and
$W_j=d_jQ_j-W_{j-1}$ ($j\ge N$). Since $j\ge N\ge J_0+1$ implies $Q_j>2Q_{j-1}$,
induction gives $|W_j|\le2kQ_j$ for all $j\ge N-1$, whence
$|M_{j+1}/Q_{j+1}-M_j/Q_j|\le2k/Q_{j+1}$. Using $Q_{m+1+s}\ge2^sQ_{m+1}$,
$x(\omega):=\lim_{j\to\infty}M_j/Q_j$ satisfies $|x(\omega)-M_m/Q_m|\le4k/Q_{m+1}$
for all $m\ge N$ (in particular $|x(\omega)|\le4k/Q_{N+1}$ since $M_N=0$).
Multiplying by $Q_m$ gives $|Q_mx(\omega)-M_m|\le4kQ_m/Q_{m+1}\le4k/A_m\le4k/A_N<1/8<1/2$
for all $m\ge N$, and at $m=N-1$ we also have $Q_{N-1}|x(\omega)|\le4kQ_{N-1}/Q_{N+1}<1/16<1/2$
(since $Q_{N+1}>32kQ_N>64kQ_{N-1}$). Hence $M_m(x(\omega))=M_m$ for all $m\ge N-1$,
no half-integer ties occur, and $d_j(x(\omega))=d_j$ for all $j\ge N$. Moreover,
$K_\Sigma:=x(\Sigma^{\{j\ge N\}})\subseteq[-4k/Q_{N+1},4k/Q_{N+1}]$, an interval
of length $<1/4<1/2$ embedding isometrically into $\R/\Z$. Pushing forward the
uniform Bernoulli product measure on $\Sigma^{\{j\ge N\}}$ yields $\mu$ on $K_\Sigma$.
For any interval $U$ of length $1/Q_{n+1}\le h<1/Q_n$ ($n>N$) and any $x,y\in K_\Sigma\cap U$,
we have $|M_j(x)-M_j(y)|<\frac14+\frac{Q_j}{Q_n}\le\frac14+\frac16<1$ for $N\le j<n$
(so $M_j(x)=M_j(y)$) and $|M_n(x)-M_n(y)|<\frac14+Q_nh<\frac54<2$ (at most $3$ values).
Thus $K_\Sigma\cap U$ meets at most $3$ cylinders of depth $n-N$, so
$\mu(U)\le3|\Sigma|^{-(n-N)}=3|\Sigma|^{N+1}|\Sigma|^{-(n+1)}\le3|\Sigma|^{N+1}g_{|\Sigma|}(h)$
(as $n+1>J_0$), where $g_{2k+1}:=h_k^{(r,c,d)}$ and $g_2:=g_2^{(r,c,d)}$. The
mass distribution principle yields $\mathcal H^{g_{|\Sigma|}}(K_\Sigma)\ge(3|\Sigma|^{N+1})^{-1}>0$.
Taking $\Sigma=\{-k,\dots,k\}$ and subtracting the $\mathcal H^{h_k^{(r,c,d)}}$-null
set $\{\Lam^{(r,c,d)}<k/c\}$ gives $\mathcal H^{h_k^{(r,c,d)}}(S_k^{(r,c,d)})>0$;
taking $\Sigma=\{-k,+k\}$ gives $K_\Sigma\subseteq E_{\mathrm{lim}}^{(r,c,d)}(k/c)$
and $\mathcal H^{g_2^{(r,c,d)}}(E_{\mathrm{lim}}^{(r,c,d)}(k/c))>0$.

By Lemma~\ref{lem:poly-lattice}(iv), $j^r\|Q_jy\|\to0$ for every $y\in L^{(r,c,d)}$,
so both $S_k^{(r,c,d)}$ and $E_{\mathrm{lim}}^{(r,c,d)}(k/c)$ are Borel,
Lebesgue-null (by Theorem~\ref{thm:poly-phase-transition}(i), or directly since
$\operatorname{Leb}_{\R/\Z}(\{\|Q_jx\|\le m/j^r\})=2m/j^r\to0$), and invariant
under translation by the dense subgroup $\bar L^{(r,c,d)}\subset\R/\Z$ of
Lemma~\ref{lem:poly-lattice}(iii). The exact Haar-uniqueness and compactness
covering arguments of Theorem~\ref{thm:gauge} and Corollary~\ref{cor:exact-gauge}(i)
then force both $\mathcal H^{h_k^{(r,c,d)}}(S_k^{(r,c,d)}\cap I)=\infty$ and
$\mathcal H^{g_2^{(r,c,d)}}(E_{\mathrm{lim}}^{(r,c,d)}(k/c)\cap I)=\infty$ on
every nonempty open arc $I\subseteq\R/\Z$.

Finally, for~(iii), take $x_k^+$ from the constant sequence $d_j\equiv+k$ ($j\ge N$).
Since $|u_j(x_k^+)|<1/8<1/2$ for all $j\ge N-1$, $x_k^-=-x_k^+$ satisfies
$M_j(x_k^-)=-M_j(x_k^+)$ and $u_j(x_k^-)=-u_j(x_k^+)$ for $j\ge N-1$, so
$d_j(x_k^-)=-k$ for $j\ge N$ and $x_k^\pm\in E_{\mathrm{lim}}^{(r,c,d),\pm}(k/c)$.
If $x,y\in E_{\mathrm{lim}}^{(r,c,d),\pm}(k/c)$, then
$j^r\|Q_j(x-y)\|\le j^r|u_j(x)-u_j(y)|\to0$, so $x-y\in L^{(r,c,d)}$ by
Lemma~\ref{lem:poly-lattice}(iv). Conversely, if $z\in L^{(r,c,d)}$, then
$u_j(z)=o(j^{-r})\to0$ while $u_j(x_k^\pm)=O(j^{-r})\to0$, so eventually
$|u_j(x_k^\pm)+u_j(z)|<1/2$, forcing $u_j(x_k^\pm+z)=u_j(x_k^\pm)+u_j(z)$ and
$j^ru_j(x_k^\pm+z)\to\pm k/c$. Thus $E_{\mathrm{lim}}^{(r,c,d),\pm}(k/c)=x_k^\pm+L^{(r,c,d)}$.
\end{proof}

\begin{corollary}[Exact critical gauge for $(cj^r+d)$, Layer~A in Lean]\label{cor:poly-exact-gauge}
Let $k\ge1$ and let $g:[0,\infty)\to[0,\infty)$ be any nondecreasing gauge with $g(0)=0$.
\begin{enumerate}[label=(\roman*),leftmargin=2em]
\item If $g(1/Q_n)/h_k^{(r,c,d)}(1/Q_n)\to0$ as $n\to\infty$, then
$\mathcal H^g(\{\Lam^{(r,c,d)}\le k/c\})=0$.
\item If $h_k^{(r,c,d)}(t)/g(t)\to0$ as $t\to0^+$, then $S_k^{(r,c,d)}$ is not
$\sigma$-finite for $\mathcal H^g$ (in particular $\mathcal H^g(S_k^{(r,c,d)})=\infty$).
\end{enumerate}
The exact same dichotomy holds for $E_{\mathrm{lim}}^{(r,c,d)}(k/c)$ with
$g_2^{(r,c,d)}$ in place of $h_k^{(r,c,d)}$.
\end{corollary}

\begin{proof}
This follows by the identical covering comparison and $\sigma$-finiteness
contradiction of Corollary~\ref{cor:exact-gauge}(ii)--(iii), using the upper
covers from Theorem~\ref{thm:poly-gauge-trichotomy}.
\end{proof}

\begin{remark}[Interpolation-free logarithmic scale and the $1/r$ factor]\label{rem:poly-gauge-scale}
For $1/Q_{n+1}\le t\le 1/Q_n$ ($n\ge J_0$), $-\log h_k^{(r,c,d)}(t)$ lies between
$n\log(2k+1)$ and $(n+1)\log(2k+1)$. By Lemma~\ref{lem:poly-growth}(iv),
$\log(1/t)=rn\log n+O(n)$, so $\log\log(1/t)=\log n+O(\log\log n)\sim\log n$,
yielding $\frac{\log(1/t)}{\log\log(1/t)}\sim rn$ and therefore
\[
-\log h_k^{(r,c,d)}(t)\sim\frac{\log(2k+1)}{r}\cdot\frac{\log(1/t)}{\log\log(1/t)},\qquad
-\log g_2^{(r,c,d)}(t)\sim\frac{\log2}{r}\cdot\frac{\log(1/t)}{\log\log(1/t)}
\qquad(t\downarrow0).
\]
Consequently, any gauge $g$ whose logarithmic rate against
$\frac{\log(1/t)}{\log\log(1/t)}$ exceeds $\frac{\log(2k+1)}{r}$ gives
$\mathcal H^g(\{\Lam^{(r,c,d)}\le k/c\})=0$, while any strictly smaller rate
yields non-$\sigma$-finite $\mathcal H^g$ measure on $S_k^{(r,c,d)}$.
\end{remark}

\begin{theorem}[Critical dimension zero, subcritical phase transition $1-\alpha/r$,
and supercritical collapse, Layer~A in Lean]\label{thm:poly-phase-transition}
For any $r,c\in\Z_{\ge1}$, $d\in\Z$, and $(Q_j)_{j\ge0}$ satisfying \eqref{eq:poly-rec}:
\begin{enumerate}[label=(\roman*),leftmargin=2em]
\item \textbf{Critical weight ($\alpha=r$):} For every finite $t\in\R$,
$\dimH\{x\in\R:\Lam^{(r,c,d)}(x)\le t\}=0$. Thus
$F^{(r,c,d)}:=\{x\in\R:\Lam^{(r,c,d)}(x)<\infty\}$ has Hausdorff dimension $0$
and Lebesgue measure $0$, while its complement
$F_\infty^{(r,c,d)}:=\{x\in\R:\Lam^{(r,c,d)}(x)=\infty\}$ has full Lebesgue
measure and $\dimH(F_\infty^{(r,c,d)}\cap[a,b])=1$ for every $a<b$.
\item \textbf{Subcritical weights ($0\le\alpha<r$):} For every $0\le\alpha<r$ and
every $t\ge0$ (with $0\le t\le1/2$ when $\alpha=0$),
\[
\dimH\bigl\{x\in[0,1]:j^\alpha\|Q_jx\|\longrightarrow t\bigr\}
=\dimH\bigl\{x\in\R:\Lam_\alpha^{(r,c,d)}(x)\le t\bigr\}
=1-\frac{\alpha}{r}.
\]
(For $\alpha=0$ and $t>1/2$, the prescribed-limit set is empty while
$\{\Lam_0^{(r,c,d)}\le t\}=\R$ has dimension $1=1-0/r$.)
\item \textbf{Supercritical weights ($\alpha>r$):} For every finite $t\ge0$,
$\{x\in\R:\Lam_\alpha^{(r,c,d)}(x)\le t\}=L^{(r,c,d)}$ (countable, Hausdorff
dimension $0$), while for $t<0$ the sublevel set is empty.
\end{enumerate}
\end{theorem}

\begin{proof}
\emph{Step 1: Upper bound $\dimH\{x\in\R:\Lam_\alpha^{(r,c,d)}(x)\le t\}\le1-\alpha/r$
for all $0\le\alpha\le r$.} Fix $0\le\alpha\le r$ and $t\in\R$, and put
$t_0:=\max(t,0)+1>0$. By $\Z$-periodicity and countable stability, it suffices to
show $\dimH C_{J,\alpha}\le1-\alpha/r$ for each $J\ge J_1$ (from
Lemma~\ref{lem:poly-growth}), where
$C_{J,\alpha}:=[0,1]\cap\bigcap_{j=J}^\infty\{x:\|Q_jx\|\le t_0j^{-\alpha}\}$.
Put $r_{j,\alpha}:=t_0/(j^\alpha Q_j)$ and
$I_{j,\alpha}(k):=[k/Q_j-r_{j,\alpha},k/Q_j+r_{j,\alpha}]$. Any $x\in C_{J,\alpha}$
satisfies $x\in I_{j,\alpha}(k_j(x))$ for $k_j(x)=\mathrm{round}(Q_jx)\in\{0,\dots,Q_j\}$,
with $|k_j(x)/Q_j-k_{j+1}(x)/Q_{j+1}|\le r_{j,\alpha}+r_{j+1,\alpha}$. Multiplying
by $2Q_{j+1}$ and using Lemma~\ref{lem:poly-growth}(ii) together with
$1/(j+1)^\alpha\le1\le j^{r-\alpha}$ ($0\le\alpha\le r$, $j\ge1$), each parent
$k$ yields child centers $k'$ in an interval of length
\[
2(r_{j,\alpha}+r_{j+1,\alpha})Q_{j+1}
=2t_0\left(\frac{Q_{j+1}}{j^\alpha Q_j}+\frac{1}{(j+1)^\alpha}\right)
\le 2t_0(c+|d|+2)j^{r-\alpha},
\]
so each parent has at most $Kj^{r-\alpha}$ children, where $K:=2t_0(c+|d|+2)+1$.
Hence the surviving center set $S_{n,\alpha}\subseteq\{0,\dots,Q_{J+n}\}$ at
level $m:=J+n\ge J_1$ satisfies
$|S_{n,\alpha}|\le(Q_J+1)K^n(m!)^{r-\alpha}$, while
Lemma~\ref{lem:poly-growth}(iii) and $m^\alpha\ge1$ give
$\operatorname{diam}I_{m,\alpha}(k)=2t_0/(m^\alpha Q_m)\le2t_0/(C_{J_1}(c/2^{r+1})^m(m!)^r)\to0$.
For any $s>1-\alpha/r$, set $\beta:=rs-(r-\alpha)=r(s-(1-\alpha/r))>0$ and
$K_2:=K(2^{r+1}/c)^s>0$. Using $m!\ge(m/\ee)^m$, the $s$-dimensional covering
cost satisfies
\begin{align*}
\sum_{k\in S_{n,\alpha}}\bigl(\operatorname{diam}I_{m,\alpha}(k)\bigr)^s
&\le(Q_J+1)K^{-J}\left(\frac{2t_0}{C_{J_1}}\right)^s K_2^m(m!)^{-\beta}\\
&\le(Q_J+1)K^{-J}\left(\frac{2t_0}{C_{J_1}}\right)^s
\left(\frac{\ee K_2^{1/\beta}}{m}\right)^{\beta m}\longrightarrow0
\end{align*}
as $m\to\infty$. Thus $\mathcal H^s(C_{J,\alpha})=0$ for every $s>1-\alpha/r$,
giving $\dimH\{x\in\R:\Lam_\alpha^{(r,c,d)}(x)\le t\}\le1-\alpha/r$. Taking
$\alpha=r$ proves~(i) (with $\operatorname{Leb}(F^{(r,c,d)})=0$ and
$\dimH(F_\infty^{(r,c,d)}\cap[a,b])=1$ exactly as in Corollary~\ref{cor:infinite-level}).

\emph{Step 2: Lower bound $\dimH\{x\in[0,1]:j^\alpha\|Q_jx\|\to t\}\ge1-\alpha/r$
for $0\le\alpha<r$.} Fix $0\le\alpha<r$ and $t\ge0$ (with $0\le t\le1/2$ if
$\alpha=0$). Put $w_j:=j^{-\alpha}/\log(j+2)$ and $r_j:=w_j/Q_j$. Since
$r-\alpha>0$ and $c\ge1$, Lemma~\ref{lem:poly-growth}(ii) gives
$2w_j(cj^r+d-1)-1\sim2c\,j^{r-\alpha}/\log(j+2)$; choose $N\ge\max(J_1,3)$ large
enough that for all $j\ge N$, $w_j<1/8$, $tj^{-\alpha}+w_j\le1/2$ (when
$\alpha>0$ or $\alpha=0,t<1/2$),
\[
2w_j(cj^r+d-1)-1\ge\frac{j^{r-\alpha}}{\log(j+2)}\ge b_j:=\left\lfloor\frac{j^{r-\alpha}}{\log(j+2)}\right\rfloor\ge2,
\]
and $Q_N\ge2$ (so $I_N(1)\subseteq(0,1)$). Defining the clamped shifts
$\theta_j\in[w_j,1/2-w_j]$ and disjoint intervals
$I_j(k):=[(k+\theta_j-w_j)/Q_j,(k+\theta_j+w_j)/Q_j]$ as in
Theorem~\ref{conj:weight}, every $x\in I_j(k)$ obeys
$|j^\alpha\|Q_jx\|-t|\le2/\log(j+2)\to0$. By Lemma~\ref{lem:poly-growth}(ii)
and $w_{j+1}\le w_j$, the nesting count inside each $I_j(k)$ is at least
\[
2(r_j-r_{j+1})Q_{j+1}-1=2\left(w_j\frac{Q_{j+1}}{Q_j}-w_{j+1}\right)-1
>2w_j(cj^r+d-1)-1\ge b_j\ge2.
\]
Retaining $b_j$ children per parent from $I_N(1)$ gives
$M_n:=\prod_{j=N}^{n-1}b_j$ intervals at level $n\ge N$ and a compact subset
$K\subseteq\{x\in[0,1]:j^\alpha\|Q_jx\|\to t\}$ carrying the uniform Bernoulli
probability measure $\mu$ ($\mu(I_n(k))=M_n^{-1}$). Since
$\log b_j=(r-\alpha)\log j-\log\log(j+2)+O(1)$, we have
$\log M_n=(r-\alpha)n\log n+o(n\log n)$, while Lemma~\ref{lem:poly-growth}(iv)
gives $\log Q_{n+1}=rn\log n+O(n)$. Thus for any $0<s<1-\alpha/r$,
\[
\log\left(\frac{3}{M_nQ_{n+1}^{-s}}\right)
=s\log Q_{n+1}-\log M_n+\log3
=-r\left(\left(1-\frac{\alpha}{r}\right)-s\right)n\log n+o(n\log n)\longrightarrow-\infty,
\]
so $3/M_n\le Q_{n+1}^{-s}$ for all $n\ge N_s$. Any set $U_i$ with
$Q_{n_i+1}^{-1}\le h_i:=\operatorname{diam}(U_i)<Q_{n_i}^{-1}$ ($n_i\ge N_s$)
satisfies $h_i+2r_{n_i}<(1+1/4)Q_{n_i}^{-1}<2Q_{n_i}^{-1}$ and thus meets at
most $3$ retained intervals, yielding $\mu(U_i)\le3/M_{n_i}\le Q_{n_i+1}^{-s}\le h_i^s$.
The mass distribution principle gives $\mathcal H^s(K)\ge1$, hence
$\dimH\{x\in[0,1]:j^\alpha\|Q_jx\|\to t\}\ge1-\alpha/r$. Combined with Step~1,
this proves~(ii).

\emph{Step 3: Supercritical collapse ($\alpha>r$).} If $t<0$, the sublevel set
is empty because $\Lam_\alpha^{(r,c,d)}\ge0$. Fix $t\ge0$. If
$\Lam_\alpha^{(r,c,d)}(x)\le t$, then $\|Q_jx\|=O(j^{-\alpha})=o(j^{-r})$, so
$\Lam^{(r,c,d)}(x)=0$ and $x\in L^{(r,c,d)}$ by Lemma~\ref{lem:poly-lattice}(iv);
conversely, every $x\in L^{(r,c,d)}$ satisfies $\Lam_\alpha^{(r,c,d)}(x)=0\le t$
by Lemma~\ref{lem:poly-lattice}(iv), proving~(iii).
\end{proof}

\begin{remark}[Comparison with Bugeaud--Liao--Rams]\label{rem:poly-BLR}
For $\alpha=0$ the dimension assertion in Theorem~\ref{thm:poly-phase-transition}(ii)
also follows from Corollary~8 of \cite{BLR} (stated there without proof), since
$Q_{j+1}/Q_j\to\infty$, by taking
constant centres $y_j=t$ exactly as for the convergent denominators of~$e$.
For $0<\alpha<r$, the right-hand side of the dimension formula in
\cite[Theorem~7]{BLR}, evaluated formally at $f_j(x)=Q_jx$ and $\delta_j=\kappa j^{-\alpha}$ ($\kappa>0$, $j\ge1$),
is $\liminf_n(\log Q_n+\sum_{1\le j<n}\log\delta_j)/(\log Q_n-\log\delta_n)=1-\alpha/r$,
the value in~(ii). That theorem does not apply, however: its condition~(1.6) requires
$\log f'_{j+1}/\log f'_j\to\infty$, whereas here $\log Q_{j+1}/\log Q_j\to1$, and the
proof above does not use that condition. (Hypotheses as stated in arXiv:1710.03596v1, whose proof sketch indicates that
condition~(1.6) may be removable by an interpolation device of Baker; this is not
stated there as a theorem.)
\end{remark}

\paragraph{Lean status for Section~\ref{sec:poly-rec}.}
The entire general-$(r,c,d)$ theory of Section~\ref{sec:poly-rec} for
\emph{all} degrees $r\ge1$, leading coefficients $c\ge1$, and shifts $d\in\Z$
is kernel-checked at Layer~A (base axioms only: \texttt{propext},
\texttt{Classical.choice}, \texttt{Quot.sound}; zero \texttt{sorryAx}) across
four modules totaling $68$ audited declarations under \texttt{PolyRecData}:
\texttt{EpiFormal/\allowbreak TrichotomyAndPolyExtension.lean} ($37$
declarations), \texttt{EpiFormal/\allowbreak PolyGaugeAndPhaseTransition.lean}
($22$ declarations), \texttt{EpiFormal/\allowbreak PolyWeightedUpper.lean} ($1$
declaration), and \texttt{EpiFormal/\allowbreak PolyWeightedLower.lean} ($8$
declarations):
\begin{itemize}[leftmargin=1.8em]
\item Lemma~\ref{lem:poly-growth}(i)--(iv) and decay of $m^\alpha/Q_m$
(\texttt{poly\_\allowbreak Q\_\allowbreak ge\_\allowbreak factorial},
\texttt{poly\_\allowbreak pow\_\allowbreak div\_\allowbreak Q\_\allowbreak tendsto\_\allowbreak zero},
\texttt{poly\_\allowbreak log\_\allowbreak Q\_\allowbreak asymptotic});
\item Lemma~\ref{lem:poly-lattice}(i)--(iv) including the Casorati identity,
fundamental solution error bounds, irrationality $\theta\notin\Q$, and
zero-level lattice equivalence
(\texttt{poly\_\allowbreak casorati},
\texttt{alphaU\_\allowbreak error\_\allowbreak bound},
\texttt{alphaV\_\allowbreak error\_\allowbreak bound},
\texttt{poly\_\allowbreak alpha\_\allowbreak linear\_\allowbreak relation},
\texttt{theta\_\allowbreak irrational},
\texttt{polyApprExponent\_\allowbreak eq\_\allowbreak zero\_\allowbreak tfae},
\texttt{polyApprExponent\_\allowbreak eq\_\allowbreak zero\_\allowbreak iff\_\allowbreak mem\_\allowbreak latticePoly});
\item Theorem~\ref{thm:general-rec}(i)--(ii), establishing the full discrete
spectrum $\Lam^{(r,c,d)}(\R)=\frac{1}{c}\Z_{\ge0}\cup\{\infty\}$ for all
$r\ge1$
(\texttt{poly\_\allowbreak signed\_\allowbreak defect\_\allowbreak approx},
\texttt{poly\_\allowbreak abs\_\allowbreak defect\_\allowbreak approx},
\texttt{polyApprExponent\_\allowbreak xPlusPoly},
\texttt{poly\_\allowbreak spectral\_\allowbreak gap\_\allowbreak open\_\allowbreak interval},
\texttt{poly\_\allowbreak levelSet\_\allowbreak eq\_\allowbreak on\_\allowbreak gap},
\texttt{poly\_\allowbreak levelSet\_\allowbreak below\_\allowbreak eq\_\allowbreak lattice},
\texttt{poly\_\allowbreak range\_\allowbreak apprExponent});
\item Theorem~\ref{thm:poly-gauge-trichotomy}(i)--(iii),
Corollary~\ref{cor:poly-exact-gauge}, and Remark~\ref{rem:poly-gauge-scale} on
both $\R$ and $\R/\Z$ (\texttt{UnitAddCircle}): critical-gauge and
binary-gauge nullity, $\sigma$-finiteness, open-interval/open-arc infinitude,
fast/slow gauge dichotomy, $\liminf$ four-way characterizations, logarithmic
gauge scale (\texttt{polyInterpGauge\_\allowbreak log\_\allowbreak scale\_\allowbreak tendsto},
\texttt{mkMetric\_\allowbreak polyCritGauge\_\allowbreak polyCircleExactLevel\_\allowbreak inter\_\allowbreak open\_\allowbreak eq\_\allowbreak top},
\texttt{mkMetric\_\allowbreak polyBinaryGauge\_\allowbreak polyCircleUnsignedExactLimit\_\allowbreak inter\_\allowbreak open\_\allowbreak eq\_\allowbreak top},
\texttt{polyExactLevel\_\allowbreak fast\_\allowbreak gauge\_\allowbreak null},
\texttt{polyExactLevel\_\allowbreak slow\_\allowbreak gauge\_\allowbreak non\_\allowbreak sigma\_\allowbreak finite},
\texttt{polyExactLevel\_\allowbreak liminf\_\allowbreak four\_\allowbreak way},
\texttt{polyUnsignedExactLimit\_\allowbreak liminf\_\allowbreak four\_\allowbreak way}),
and signed exact-limit coset collapse
$E_{\mathrm{lim}}^{(r,c,d),\pm}(k/c)=x_k^\pm+L^{(r,c,d)}$
(\texttt{poly\_\allowbreak trichotomy\_\allowbreak part\_\allowbreak iii\_\allowbreak complete});
\item Theorem~\ref{thm:poly-phase-transition}(i)--(iii): critical weight
$\alpha=r$ (\texttt{poly\_\allowbreak dimH\_\allowbreak levelSet\_\allowbreak finite},
\texttt{poly\_\allowbreak dimH\_\allowbreak finiteLevel\_\allowbreak zero},
\texttt{poly\_\allowbreak volume\_\allowbreak finiteLevel\_\allowbreak zero},
\texttt{poly\_\allowbreak ae\_\allowbreak apprExponent\_\allowbreak eq\_\allowbreak top},
\texttt{poly\_\allowbreak dimH\_\allowbreak infiniteLevel\_\allowbreak inter\_\allowbreak Icc\_\allowbreak eq\_\allowbreak one},
\texttt{poly\_\allowbreak dimH\_\allowbreak infiniteLevel\_\allowbreak eq\_\allowbreak one}),
subcritical $0\le\alpha<r$ exact Hausdorff dimension $1-\alpha/r$
(\texttt{Poly\allowbreak Weighted\allowbreak Upper.\allowbreak dimH\_\allowbreak poly\allowbreak Weighted\allowbreak Exponent\_\allowbreak le},
\texttt{Poly\allowbreak Weighted\allowbreak Lower.\allowbreak dimH\_\allowbreak poly\allowbreak Exact\allowbreak Limit\_\allowbreak eq},
\texttt{Poly\allowbreak Weighted\allowbreak Lower.\allowbreak dimH\_\allowbreak poly\allowbreak Weighted\allowbreak Sublevel\_\allowbreak eq},
\texttt{Poly\allowbreak Weighted\allowbreak Lower.\allowbreak poly\allowbreak Zero\allowbreak Weight\allowbreak Sublevel\_\allowbreak eq\_\allowbreak univ},
\texttt{Poly\allowbreak Weighted\allowbreak Lower.\allowbreak poly\allowbreak Exact\allowbreak Limit\_\allowbreak zero\_\allowbreak eq\_\allowbreak empty}),
and supercritical collapse $\alpha>r$
(\texttt{poly\_\allowbreak supercritical\_\allowbreak levelSet\_\allowbreak eq\_\allowbreak lattice},
\texttt{poly\_\allowbreak supercritical\_\allowbreak dimH\_\allowbreak zero},
\texttt{poly\_\allowbreak supercritical\_\allowbreak empty\_\allowbreak neg}).
\end{itemize}
These kernel-checked theorems are additionally corroborated numerically by the
$400/50$-digit \texttt{mpmath} checks in
\texttt{verify\_\allowbreak poly\_\allowbreak recurrence\_\allowbreak 2b.py}.

\subsection{Universal extension to arbitrary growing-quotient recurrences}\label{sec:growing-rec}

The polynomial recurrence $Q_{j+1}=(cj^r+d)Q_j+Q_{j-1}$ has $\log A_j=r\log j+O(1)$
and $\log Q_j\sim r\,j\log j$, so $\frac{\log A_j}{\log Q_j}\to 0$. We now drop
all functional-form assumptions on $A_j$ and consider an arbitrary integer
three-term recurrence (\texttt{GrowingRecData})
\begin{equation}\label{eq:growing-rec}
Q_{j+1}=A_j Q_j+Q_{j-1}\ \ (j\ge J_0\ge 1),\quad Q_j\in\Z_{\ge1}\ (\forall j\ge0),
\quad A_j\in\Z_{\ge2}\ (\forall j\ge J_0),\quad \lim_{j\to\infty}A_j=\infty.
\end{equation}
Set $g:=\gcd(Q_{J_0-1},Q_{J_0})\ge1$ and write $\widetilde Q_j:=Q_j/g\in\Z_{\ge1}$
for $j\ge J_0-1$, so $\gcd(\widetilde Q_{J_0-1},\widetilde Q_{J_0})=1$. Define the
fundamental companion numerator solution $\widetilde P_j\in\Z$ by
$\widetilde P_{J_0-1}=0$, $\widetilde P_{J_0}=1$, and
$\widetilde P_{j+1}=A_j\widetilde P_j+\widetilde P_{j-1}$ for $j\ge J_0$. Then
$\widetilde P_{j+1}\widetilde Q_j-\widetilde P_j\widetilde Q_{j+1}=(-1)^{j-J_0}$,
so $\theta:=\lim_{j\to\infty}\widetilde P_j/\widetilde Q_j\in\R\setminus\Q$ exists
and satisfies
\begin{equation}\label{eq:growing-theta-error}
\frac{1}{2\,\widetilde Q_j\,\widetilde Q_{j+1}}
<\left|\theta-\frac{\widetilde P_j}{\widetilde Q_j}\right|
<\frac{1}{\widetilde Q_j\,\widetilde Q_{j+1}}
\qquad(\forall j\ge J_0).
\end{equation}
Let $L_G:=\frac{1}{g}(\Z+\Z\theta)$ denote the intrinsic rank-$2$ lattice of~$G$,
and let $L_{G,\mathrm{rat}}:=\frac{1}{g}\Z\subsetneq L_G$ denote its rank-$1$
rational sublattice. Since the natural critical weight for \eqref{eq:growing-rec}
is $A_j$ rather than $j^r$, we define the critical and $\tau$-weighted
approximation functionals ($\tau\in\R$) by
\[
\Lam_G(x):=\limsup_{j\to\infty} A_j\|Q_j x\|\in[0,\infty],\qquad
\Lam_{G,\tau}(x):=\limsup_{j\to\infty} A_j^\tau\|Q_j x\|\in[0,\infty]\quad(\Lam_G=\Lam_{G,1}),
\]
with sublevel sets $E_{G,\tau}(t):=\{x\in\R:\Lam_{G,\tau}(x)\le t\}$ ($E_G(t):=E_{G,1}(t)$),
unsigned exact limit sets $E_{G,\tau,\mathrm{lim}}(t):=\{x\in\R:\lim_{j\to\infty}A_j^\tau\|Q_j x\|=t\}$
($E_{G,\mathrm{lim}}(t):=E_{G,1,\mathrm{lim}}(t)$), and signed exact limit sets
$E_{G,\mathrm{lim}}^\pm(k):=\{x\in\R:\lim_{j\to\infty}A_j(Q_j x-\mathrm{round}(Q_j x))=\pm k\}$.

\begin{theorem}[Universal discrete spectrum, Kraaikamp--Liardet equivalence, and critical gauges for arbitrary growing-quotient recurrences]\label{thm:growing-rec-spectrum}
Let $G=(J_0,A_j,Q_j)$ satisfy~\eqref{eq:growing-rec}.
\begin{enumerate}[label=(\roman*),leftmargin=2.2em]
\item \textbf{Kraaikamp--Liardet ratio equivalence and zero-level TFAE:}
For every $x\in\R$, the Kraaikamp--Liardet normalized error functional satisfies
\[
\Lam_G^{\mathrm{KL}}(x)
:=\limsup_{j\to\infty}\frac{\|Q_j x\|}{\widetilde Q_j\,|\widetilde Q_j\theta-\widetilde P_j|}
=\Lam_G(x).
\]
Moreover, the following five conditions are equivalent:
\textup{(1)}~$\Lam_G(x)=0$; \textup{(2)}~$\Lam_G(x)<1$;
\textup{(3)}~$\Lam_G^{\mathrm{KL}}(x)<1$;
\textup{(4)}~$d_j(x):=\mathrm{round}(Q_{j+1}x)-A_j\mathrm{round}(Q_j x)-\mathrm{round}(Q_{j-1}x)=0$ for all sufficiently large~$j$; and
\textup{(5)}~$x\in L_G=\frac{1}{g}(\Z+\Z\theta)$.
\item \textbf{Full quantized spectrum and coarse metric collapse:}
We have $\Lam_G(\R)=\Z_{\ge0}\cup\{\infty\}$. Every open interval $(k,k+1)$
($k\in\Z_{\ge0}$) is an empty spectral gap: $E_G(t)=E_G(k)$ for all $t\in[k,k+1)$,
$E_G(t)=L_G$ for all $t\in[0,1)$, $\dimH E_G(t)=0$ for all finite $t<\infty$,
$\dimH\{\Lam_G<\infty\}=0$, $\operatorname{Leb}\{\Lam_G<\infty\}=0$, and
$\Lam_G(x)=\infty$ for Lebesgue-almost every $x\in\R$ (with
$\dimH(\{\Lam_G=\infty\}\cap[a,b])=1$ on every interval $a<b$).
\item \textbf{Three-way positive-level splitting and critical gauges on $\R$ and $\R/\Z$:}
For each integer $k\ge1$, let $h_k^G$ and $g_2^G$ be the piecewise-linear Hausdorff
gauges determined by $h_k^G(1/Q_n)=(2k+1)^{-n}$ and $g_2^G(1/Q_n)=2^{-n}$. Then:
\begin{enumerate}[label=(\alph*),leftmargin=1.6em]
\item On both $\R$ and $\R/\Z$ (\texttt{UnitAddCircle}), the exact level
$\{\Lam_G=k\}$ has $\sigma$-finite $\mathcal H^{h_k^G}$ measure that is infinite
on every nonempty open interval/arc, whereas $\{\Lam_G<k\}$ is
$\mathcal H^{h_k^G}$-null;
\item The unsigned exact limit set $E_{G,\mathrm{lim}}(k)$ has $\sigma$-finite
$\mathcal H^{g_2^G}$ measure that is infinite on every nonempty open interval/arc
(with gauge base~$2$ independent of $k\ge1$) and is $\mathcal H^{h_k^G}$-null;
\item Each signed exact limit set collapses to a single countable $L_G$-coset:
$E_{G,\mathrm{lim}}^\pm(k)=x_k^\pm+L_G$ for explicit telescoping witnesses
$x_k^\pm\in E_{G,\mathrm{lim}}^\pm(k)$; and
\item For every nondecreasing gauge $g$ on $[0,\infty)$ with $g(0)=0$, the four-way
$\liminf$ measure classification of Corollary~\ref{cor:liminf-gauges} holds verbatim
for $\{\Lam_G=k\}$ with respect to $\liminf_{n\to\infty}(2k+1)^n g(1/Q_n)$ and for
$E_{G,\mathrm{lim}}(k)$ with respect to $\liminf_{n\to\infty}2^n g(1/Q_n)$, including
the power-gauge threshold dichotomy for $g_\beta^G(1/Q_n)=\beta^{-n}$ at $\beta=2k+1$
and $\beta=2$.
\end{enumerate}
\end{enumerate}
\end{theorem}

\begin{proof}
Write $M_j=\mathrm{round}(Q_j x)\in\Z$ and $u_j=Q_j x-M_j\in[-1/2,1/2]$. From
$Q_{j+1}=A_jQ_j+Q_{j-1}$, the integer defect $d_j=M_{j+1}-A_jM_j-M_{j-1}\in\Z$
satisfies $d_j=A_ju_j+u_{j-1}-u_{j+1}$, hence $|A_j\|Q_j x\|-|d_j||\le 1$ and,
since $A_j\to\infty$ and $|u_{j\pm1}|\le 1/2$, only steps with $A_j|u_j|$ small or
large can occur once $A_j>2$: specifically, $|d_j-A_j u_j|\le|u_{j-1}|+|u_{j+1}|$.
Whenever $\Lam_G(x)<\infty$, $u_j=O(1/A_j)\to0$, so $A_j u_j-d_j\to0$ and
$A_j\|Q_j x\|-|d_j|\to0$. Because $|d_j|\in\Z_{\ge0}$, $\Lam_G(x)=\limsup_{j\to\infty}|d_j|\in\Z_{\ge0}\cup\{\infty\}$.
Furthermore, from the continued-fraction identity
$\widetilde Q_j|\widetilde Q_j\theta-\widetilde P_j|=(\widetilde Q_{j+1}/\widetilde Q_j+\theta_j)^{-1}=(A_j+O(1))^{-1}$
with $A_j\to\infty$, the ratio $\widetilde Q_j|\widetilde Q_j\theta-\widetilde P_j|\cdot A_j$
tends to~$1$, yielding $\Lam_G^{\mathrm{KL}}(x)=\Lam_G(x)$ for all $x\in\R$.
If $d_j=0$ for all $j\ge N$, unimodularity of
$\bigl(\begin{smallmatrix}\widetilde P_{N-1}&\widetilde Q_{N-1}\\\widetilde P_N&\widetilde Q_N\end{smallmatrix}\bigr)$
gives integers $a,b$ with $M_j=a\widetilde P_j+b\widetilde Q_j$ for all $j\ge N-1$,
whence $Q_j(x-(b+a\theta)/g)=u_j+a(\widetilde Q_j\theta-\widetilde P_j)\to0$ and
$x=(b+a\theta)/g\in L_G$; conversely, every $x\in L_G$ satisfies $A_j\|Q_j x\|=O(1/Q_j)\to0$.
For any bounded digit sequence $(d_j)_{j\ge J_0}$ ($(|d_j|\le B)$), the alternating
telescoping series around index $N$ produces a real $x$ whose nearest integers
satisfy $M_{j+1}-A_j M_j-M_{j-1}=d_j$ for all $j\ge N$, realizing every level
$k\in\Z_{\ge0}$ (constant digit $d_j=k$ gives $x_k^+\in E_{G,\mathrm{lim}}^+(k)$)
and $\infty$, and yielding $E_{G,\mathrm{lim}}^\pm(k)=x_k^\pm+L_G$ by subtracting
two solutions with eventually identical digits. Finally, because $A_j\to\infty$
implies $(2k+1)^j/Q_j\to0$, the Cantor-tree Frostman measure and covering gauge
machinery of Theorems~\ref{thm:gauge} and~\ref{thm:trichotomy} and Section~\ref{sec:poly-rec}
applies with branching $(2k+1)$ on $\{\Lam_G\le k\}$ and branching~$2$ on
$E_{G,\mathrm{lim}}(k)$, establishing~(ii) and~(iii)(a)--(d).
\end{proof}

To describe the weighted Hausdorff-dimension and supercritical structure of a
general recurrence~\eqref{eq:growing-rec}, define the \emph{upper logarithmic
quotient-to-denominator exponent}
\begin{equation}\label{eq:gamma-sup}
\gamma_{\sup}(G):=\limsup_{j\to\infty}\frac{\log A_j}{\log Q_j}\in[0,\infty].
\end{equation}
Whenever $\gamma_{\sup}(G)=0$ (which includes all polynomial recurrences
$A_j=cj^r+d$ as well as $A_j=2^{j+1}$, since $\log Q_j\ge\sum_{i\le j-1}\log A_i\gg j^2$),
$\log Q_{j+1}/\log Q_j=1+\frac{\log A_j}{\log Q_j}+o(1)\to1$. When $A_j$ grows
double-exponentially or faster, however, $\gamma_{\sup}(G)>0$ measures the
non-negligible ratio between the step-by-step window contraction $A_j^{-\tau}$ and
the grid spacing $Q_j^{-1}$, producing both a subcritical Hausdorff-dimension
distortion and a supercritical secondary phase transition inside $L_G$:

\begin{theorem}[Universal subcritical dimension and supercritical phase transition]\label{thm:growing-rec-dim-phase}
Let $G=(J_0,A_j,Q_j)$ satisfy~\eqref{eq:growing-rec} and let
$\gamma_{\sup}=\gamma_{\sup}(G)\in[0,\infty]$ be defined by~\eqref{eq:gamma-sup}.
\begin{enumerate}[label=(\roman*),leftmargin=2.2em]
\item \textbf{Subcritical Hausdorff dimension ($0\le\tau<1$):}
For every $0\le\tau<1$ and every finite threshold $0\le t<\infty$,
\begin{equation}\label{eq:growing-subcritical-dim}
\dimH E_{G,\tau}(t)
=\dimH E_{G,\tau,\mathrm{lim}}(0)
=\frac{1-\tau}{1+\tau\gamma_{\sup}},
\end{equation}
interpreted as $1$ when $\tau=0$ and as $0$ when $\tau>0$ and $\gamma_{\sup}=\infty$.
\item \textbf{Supercritical collapse to $L_G$ and secondary phase transition ($\tau>1$):}
For every $\tau>1$ and every finite $0\le t<\infty$, $E_{G,\tau}(t)\subseteq L_G$
(hence $E_{G,\tau}(t)$ is countable with $\dimH E_{G,\tau}(t)=0$), and the
rational sublattice satisfies $L_{G,\mathrm{rat}}=\frac{1}{g}\Z\subseteq E_{G,\tau}(t)$.
Moreover, for every irrational lattice element $x=(m+n\theta)/g\in L_G\setminus L_{G,\mathrm{rat}}$
($m,n\in\Z$, $n\ne0$) and all sufficiently large~$j$,
\begin{equation}\label{eq:growing-supercritical-sandwich}
\frac{|n|g}{6}\,\frac{A_j^{\tau-1}}{Q_j}
\le A_j^\tau\|Q_j x\|
\le |n|g\,\frac{A_j^{\tau-1}}{Q_j}.
\end{equation}
Consequently:
\begin{enumerate}[label=(\alph*),leftmargin=1.6em]
\item If $A_j^{\tau-1}/Q_j\to0$---in particular whenever
$(\tau-1)\gamma_{\sup}<1$ (and for \emph{all} $\tau>1$ when $\gamma_{\sup}=0$)---then
$E_{G,\tau}(t)=L_G$ for every finite $0\le t<\infty$;
\item If $\limsup_{j\to\infty}A_j^{\tau-1}/Q_j=\infty$---in particular whenever
$(\tau-1)\gamma_{\sup}>1$ (and for \emph{all} $\tau>1$ when $\gamma_{\sup}=\infty$)---then
$\Lam_{G,\tau}(x)=\infty$ for every $x\in L_G\setminus L_{G,\mathrm{rat}}$, so
$E_{G,\tau}(t)=L_{G,\mathrm{rat}}=\frac{1}{g}\Z\subsetneq L_G$ for every finite
$0\le t<\infty$.
\end{enumerate}
\end{enumerate}
\end{theorem}

\begin{proof}
For~(i), fix $0\le\tau<1$ and $t<\infty$. At level $j\ge J_0+1$, each step-$j$
surviving window has length $\asymp(Q_j A_j^\tau)^{-1}$, while the step-$j$ grid
of centers $m/Q_j$ inside a step-$(j-1)$ window of length $\asymp(Q_{j-1}A_{j-1}^\tau)^{-1}$
contains $\asymp\frac{Q_j}{Q_{j-1}A_{j-1}^\tau}\asymp A_{j-1}^{1-\tau}$ children.
Thus after $n$ steps from a base index $J$, the total number of surviving windows
in $[0,1]$ is $N_n\asymp\prod_{i=J}^{J+n-1}A_i^{1-\tau}\asymp Q_{J+n}^{1-\tau}$
(up to a factor $2^{\pm n}$, which is $Q_{J+n}^{o(1)}$ because $\log Q_j/j\to\infty$),
while the level-$(J+n)$ window diameter is $\delta_{J+n}\asymp(Q_{J+n}A_{J+n}^\tau)^{-1}$.
Taking logarithms gives
\[
\frac{\log N_n}{-\log\delta_{J+n}}
=\frac{(1-\tau)\log Q_{J+n}+O(n)}{\log Q_{J+n}+\tau\log A_{J+n}+O(1)}
=\frac{1-\tau+o(1)}{1+\tau\frac{\log A_{J+n}}{\log Q_{J+n}}+o(1)}.
\]
Choosing a subsequence $(j_k)$ along which $\frac{\log A_{j_k}}{\log Q_{j_k}}\to\gamma_{\sup}$
shows that for any $s>\frac{1-\tau}{1+\tau\gamma_{\sup}}$, the $s$-dimensional
covering sum $N_n\delta_{J+n}^s$ tends to $0$ along $J+n=j_k$, proving
$\dimH E_{G,\tau}(t)\le\frac{1-\tau}{1+\tau\gamma_{\sup}}$. Conversely, for any
$0<s<\frac{1-\tau}{1+\tau\gamma_{\sup}}$ (so $\gamma_{\sup}<\infty$), for all
sufficiently large $j$ we have $(Q_j A_j^\tau)^s\le Q_j^{1-\tau-\delta}$ for some
$\delta>0$; running our shifted-grid homogeneous Moran/Frostman tree construction
with decay schedule $\varepsilon_j\to0$ inside $E_{G,\tau,\mathrm{lim}}(0)$ shows
that the local mass of every ball $B(x,r)$ satisfies $\mu(B(x,r))\ll r^s$, giving
$\dimH E_{G,\tau,\mathrm{lim}}(0)\ge s$ and proving~\eqref{eq:growing-subcritical-dim}.

For~(ii), if $\tau>1$ and $x\in E_{G,\tau}(t)$ with $t<\infty$, then for large $j$,
$A_j\|Q_j x\|=A_j^{-(\tau-1)}(A_j^\tau\|Q_j x\|)\le(t+1)A_j^{-(\tau-1)}\to0$, so
$\Lam_G(x)=0$ and $x\in L_G$ by Theorem~\ref{thm:growing-rec-spectrum}(i). For
$x=m/g\in L_{G,\mathrm{rat}}$, $Q_j x=m\widetilde Q_j\in\Z$ for all $j\ge J_0-1$,
so $\|Q_j x\|=0$ and $x\in E_{G,\tau}(t)$. Now let $x=(m+n\theta)/g\in L_G$ with
$n\in\Z\setminus\{0\}$ and put $M_j:=m\widetilde Q_j+n\widetilde P_j\in\Z$. Then
$Q_j x-M_j=n\widetilde Q_j(\theta-\widetilde P_j/\widetilde Q_j)$, so by
\eqref{eq:growing-theta-error} and $Q_{j+1}=g\widetilde Q_{j+1}$,
\[
\frac{|n|g}{2 Q_{j+1}}
<|Q_j x-M_j|
<\frac{|n|g}{Q_{j+1}}.
\]
For all large $j$, $|n|g/Q_{j+1}<1/2$, so $\mathrm{round}(Q_j x)=M_j$ and
$\|Q_j x\|=|Q_j x-M_j|$. Using $A_j Q_j\le Q_{j+1}\le 2 A_j Q_j$ (for $j\ge J_0+1$)
and multiplying by $A_j^\tau$ gives~\eqref{eq:growing-supercritical-sandwich}.
Claims~(a) and~(b) follow immediately upon taking $\limsup_{j\to\infty}$ in~\eqref{eq:growing-supercritical-sandwich}
and relating $\frac{\log(A_j^{\tau-1}/Q_j)}{\log Q_j}=(\tau-1)\frac{\log A_j}{\log Q_j}-1$
to $(\tau-1)\gamma_{\sup}-1$.
\end{proof}

\begin{remark}[Comparison with Bugeaud--Liao--Rams]\label{rem:growing-BLR}
Part~(i) of Theorem~\ref{thm:growing-rec-dim-phase} overlaps with \cite{BLR} in two
boundary regimes. For $\tau=0$, since $Q_{j+1}/Q_j\ge A_j\to\infty$, Corollary~8 of
\cite{BLR} (stated there without proof) already gives
$\dimH E_{G,0,\mathrm{lim}}(0)=1$, hence also
$\dimH E_{G,0}(t)=1$. For $0<\tau<1$ and $\lim_{j\to\infty}\log A_j/\log Q_j=\infty$
(a limit, not only a $\limsup$), one has $\log Q_{j+1}/\log Q_j\to\infty$ because
$Q_{j+1}\ge A_jQ_j$; then, for $N$ large, the maps $x\mapsto Q_{j+N}x$ with
$\delta_j=(t+1)A_{j+N}^{-\tau}$ satisfy conditions~(1.1), (1.5) and~(1.6) of
\cite[Theorem~7]{BLR}, its formula gives local dimension $0$ at every point (hence
Hausdorff dimension $0$, covering $\R$ by countably many bounded intervals), and
$E_{G,\tau}(t)\subseteq\bigcup_N\{x:\|Q_jx\|\le(t+1)A_j^{-\tau}\ \forall j\ge N\}$
yields $\dimH E_{G,\tau}(t)=0$ by countable stability, hence also
$\dimH E_{G,\tau,\mathrm{lim}}(0)=0$ since $E_{G,\tau,\mathrm{lim}}(0)\subseteq E_{G,\tau}(t)$.
In all other cases---in
particular whenever $\gamma_{\sup}<\infty$, where $\log Q_{j+1}/\log Q_j$ stays
bounded---condition~(1.6) fails and \cite[Theorem~7]{BLR} does not apply. Its
right-hand side, evaluated formally at $f_j(x)=Q_jx$ and $\delta_j=\kappa A_j^{-\tau}$
($\kappa>0$, $j\ge J_0$), is
nevertheless $(1-\tau)/(1+\tau\gamma_{\sup})$ (using $\log Q_n=\sum_{J_0\le j<n}\log A_j+O(n)$
and $n=o(\log Q_n)$), so~\eqref{eq:growing-subcritical-dim} may be read as an extension
of that formula beyond condition~(1.6). (Hypotheses as stated in arXiv:1710.03596v1, whose proof sketch indicates that
condition~(1.6) may be removable by an interpolation device of Baker; this is not
stated there as a theorem.)
\end{remark}

\begin{corollary}[Three canonical growth regimes and the critical boundary $\tau_{\mathrm{crit}}=2$]\label{cor:growing-rec-three-regimes}
Consider the unit-initialized recurrences ($J_0=1$, $Q_0=Q_1=1$, $g=1$,
$L_{G,\mathrm{rat}}=\Z$, $L_G=\Z+\Z\theta_G$):
\begin{enumerate}[label=(\roman*),leftmargin=2.2em]
\item \textbf{Regime~I (exponential partial quotients, $A_j=2^{j+1}$):}
We have $\gamma_{\sup}=0$, $\dimH E_{G,\tau}(t)=\dimH E_{G,\tau,\mathrm{lim}}(0)=1-\tau$
for all $0\le\tau<1$ and $t<\infty$, and $E_{G,\tau}(t)=L_G$ for all $\tau>1$ and $t<\infty$
(no secondary supercritical collapse occurs).
\item \textbf{Regime~II (double-exponential partial quotients, $A_j=2^{2^j}$):}
For all $j\ge2$, $\frac14 A_j\le Q_j\le A_j-1$ (with $\lim_{j\to\infty}Q_j/A_j\approx0.316420$),
so $\gamma_{\sup}=1$. Consequently:
\begin{itemize}[leftmargin=1.5em]
\item for $0\le\tau<1$ and $t<\infty$, $\dimH E_{G,\tau}(t)=\dimH E_{G,\tau,\mathrm{lim}}(0)=\frac{1-\tau}{1+\tau}$;
\item for $1<\tau<2$ and $t<\infty$, $E_{G,\tau}(t)=L_G=\Z+\Z\theta_G$;
\item for $\tau>2$ and $t<\infty$, $E_{G,\tau}(t)=\Z\subsetneq L_G$; and
\item at the critical secondary-transition boundary $\tau=\tau_{\mathrm{crit}}=2$,
the sublevel sets stratify the rank-$2$ lattice $L_G$ by the coefficient~$|n|$:
$E_{G,2}(0)=\Z$, the generator $\theta_G\in L_G\setminus\Z$ satisfies
$\frac16\le\Lam_{G,2}(\theta_G)\le 4$ (with numerical limit
$\Lam_{G,2}(\theta_G)=(\lim_{j\to\infty}Q_j/A_j)^{-1}\approx3.15973\in[1/6,4]$, hence
$\theta_G\in E_{G,2}(4)\setminus\Z$), and $\bigcup_{K\in\Z_{\ge0}}E_{G,2}(K)=L_G$.
\end{itemize}
\item \textbf{Regime~III (super-double-exponential partial quotients, $A_j=2^{2^{j^2}}$):}
We have $\gamma_{\sup}=\infty$. Thus $\dimH E_{G,0}(t)=\dimH E_{G,0,\mathrm{lim}}(0)=1$,
whereas $\dimH E_{G,\tau}(t)=\dimH E_{G,\tau,\mathrm{lim}}(0)=0$ for all $0<\tau<1$
and $t<\infty$, and $E_{G,\tau}(t)=\Z\subsetneq L_G$ for all $\tau>1$ and $t<\infty$.
\end{enumerate}
\end{corollary}

\begin{proof}
For~(i), $\sum_{k=1}^{j-1}\log_2 A_k=\sum_{k=1}^{j-1}(k+1)=\frac{(j-1)(j+2)}{2}$,
so $\frac{\log A_j}{\log Q_j}\le\frac{2(j+1)}{j(j-1)}\to0$ and $\gamma_{\sup}=0$.
For~(ii), $A_{k+1}=A_k^2$ for $A_k=2^{2^k}$; induction from $Q_2=5\in[\frac14 A_2,A_2-1]$
($A_2=16$) gives $\frac14 A_j=\prod_{k=1}^{j-1}A_k\le Q_j$ and
$Q_{j+1}=A_jQ_j+Q_{j-1}\le A_j(A_j-1)+(A_j-1)=A_j^2-1=A_{j+1}-1$, so
$1\le\frac{\log A_j}{\log Q_j}\le 1+\frac{\log 4}{\log Q_j}\to1$ and $\gamma_{\sup}=1$.
At $\tau=2$, \eqref{eq:growing-supercritical-sandwich} with $1\le A_j/Q_j\le 4$
gives $\frac{|n|}{6}\le A_j^2\|Q_j x\|\le 4|n|$ eventually for $x=m+n\theta_G$ ($n\ne0$),
so $\Lam_{G,2}(x)\in[|n|/6,4|n|]$, proving $E_{G,2}(0)=\Z$, $\theta_G\in E_{G,2}(4)\setminus\Z$,
and $\bigcup_{K\ge0}E_{G,2}(K)=L_G$. For~(iii), $Q_j\le 2^{j\cdot 2^{(j-1)^2}}$ while
$A_j=2^{2^{j^2}}$, so $\frac{\log A_j}{\log Q_j}\ge\frac{2^{j^2-(j-1)^2}}{j}\ge\frac{2^j}{j+1}\to\infty$
and $\gamma_{\sup}=\infty$.
\end{proof}

\paragraph{Lean~4 formalization of Section~\ref{sec:growing-rec}.}
All results of Section~\ref{sec:growing-rec} are kernel-checked at Layer~A in two
self-contained modules:
\texttt{EpiFormal/\allowbreak General\allowbreak Growing\allowbreak Rec.lean}
(all $21$ audited declarations: \texttt{PolyRecData.\allowbreak toGrowingRecData},
\texttt{grow\_\allowbreak pow\_\allowbreak div\_\allowbreak Q\_\allowbreak tendsto\_\allowbreak zero},
\texttt{grow\_\allowbreak nat\_\allowbreak div\_\allowbreak log\_\allowbreak Q\_\allowbreak tendsto\_\allowbreak zero},
\texttt{grow\_\allowbreak alpha\_\allowbreak linear\_\allowbreak relation},
\texttt{growKLRatio\_\allowbreak eq\_\allowbreak growLambda},
\texttt{growLambda\_\allowbreak eq\_\allowbreak zero\_\allowbreak tfae},
\texttt{growLevelSet\_\allowbreak zero\_\allowbreak eq\_\allowbreak lattice},
\texttt{grow\_\allowbreak telescoping\_\allowbreak digit\_\allowbreak realize},
\texttt{grow\_\allowbreak spectral\_\allowbreak gap\_\allowbreak open\_\allowbreak interval},
\texttt{dimH\_\allowbreak growLevelSet\_\allowbreak nat\_\allowbreak eq\_\allowbreak zero},
\texttt{dimH\_\allowbreak growFiniteSet\_\allowbreak eq\_\allowbreak zero},
\texttt{volume\_\allowbreak growFiniteSet\_\allowbreak eq\_\allowbreak zero},
\texttt{ae\_\allowbreak growLambda\_\allowbreak eq\_\allowbreak top},
\texttt{dimH\_\allowbreak growInfiniteSet\_\allowbreak eq\_\allowbreak one},
\texttt{grow\_\allowbreak range\_\allowbreak lambda},
\texttt{growSignedExactLimitSet\_\allowbreak coset\_\allowbreak collapse},
\texttt{grow\_\allowbreak three\_\allowbreak way\_\allowbreak gauge\_\allowbreak splitting},
\texttt{grow\_\allowbreak power\_\allowbreak gauge\_\allowbreak dichotomy},
\texttt{growExactLevel\_\allowbreak power\_\allowbreak gauge\_\allowbreak dichotomy},
\texttt{growExactLevel\_\allowbreak liminf\_\allowbreak four\_\allowbreak way}, and
\texttt{growUnsignedExactLimit\_\allowbreak liminf\_\allowbreak four\_\allowbreak way}) and
\texttt{EpiFormal/\allowbreak General\allowbreak Growing\allowbreak Dim.lean}
(all $29$ audited declarations:
\texttt{grow\_\allowbreak Q\_\allowbreak le\_\allowbreak two\_\allowbreak pow\_\allowbreak mul\_\allowbreak prod\_\allowbreak A},
\texttt{grow\_\allowbreak prod\_\allowbreak A\_\allowbreak le\_\allowbreak Q\_\allowbreak simple},
\texttt{dimH\_\allowbreak grow\allowbreak Weighted\allowbreak LevelSet\_\allowbreak le},
\texttt{le\_\allowbreak dimH\_\allowbreak grow\allowbreak Weighted\allowbreak ExactLimitSet\_\allowbreak zero},
\texttt{dimH\_\allowbreak grow\allowbreak Weighted\allowbreak LevelSet\_\allowbreak eq},
\texttt{dimH\_\allowbreak grow\allowbreak Weighted\allowbreak ExactLimitSet\_\allowbreak zero\_\allowbreak eq},
\texttt{grow\allowbreak Weighted\allowbreak LevelSet\_\allowbreak subset\_\allowbreak lattice},
\texttt{dimH\_\allowbreak grow\allowbreak Weighted\allowbreak LevelSet\_\allowbreak of\_\allowbreak one\_\allowbreak lt},
\texttt{grow\allowbreak RatLattice\_\allowbreak subset\_\allowbreak weightedLevelSet},
\texttt{grow\allowbreak RatLattice\_\allowbreak ssubset\_\allowbreak growLatticeL},
\texttt{grow\allowbreak Weighted\allowbreak Seq\_\allowbreak bounds\_\allowbreak of\_\allowbreak lattice\_\allowbreak irrational},
\texttt{grow\allowbreak Weighted\allowbreak LevelSet\_\allowbreak eq\_\allowbreak lattice\_\allowbreak of\_\allowbreak tendsto\_\allowbreak zero},
\texttt{grow\allowbreak Weighted\allowbreak LevelSet\_\allowbreak eq\_\allowbreak lattice\_\allowbreak of\_\allowbreak mul\_\allowbreak gamma\_\allowbreak lt\_\allowbreak one},
\texttt{grow\allowbreak Weighted\allowbreak LevelSet\_\allowbreak eq\_\allowbreak lattice\_\allowbreak of\_\allowbreak gamma\_\allowbreak zero},
\texttt{grow\allowbreak Weighted\allowbreak LevelSet\_\allowbreak eq\_\allowbreak ratLattice\_\allowbreak of\_\allowbreak limsup\_\allowbreak top},
\texttt{grow\allowbreak Weighted\allowbreak LevelSet\_\allowbreak eq\_\allowbreak ratLattice\_\allowbreak of\_\allowbreak one\_\allowbreak lt\_\allowbreak mul\_\allowbreak gamma},
\texttt{grow\allowbreak Weighted\allowbreak LevelSet\_\allowbreak eq\_\allowbreak ratLattice\_\allowbreak of\_\allowbreak gamma\_\allowbreak top},
\texttt{growGammaSup\_\allowbreak exp\allowbreak Growing\allowbreak Rec},
\texttt{dimH\_\allowbreak exp\allowbreak Growing\allowbreak Rec\_\allowbreak sub\allowbreak critical},
\texttt{exp\allowbreak Growing\allowbreak Rec\_\allowbreak super\allowbreak critical},
\texttt{growGammaSup\_\allowbreak double\allowbreak Exp\allowbreak Growing\allowbreak Rec},
\texttt{dimH\_\allowbreak double\allowbreak Exp\allowbreak Growing\allowbreak Rec\_\allowbreak sub\allowbreak critical},
\texttt{double\allowbreak Exp\allowbreak Growing\allowbreak Rec\_\allowbreak super\allowbreak critical\_\allowbreak lt\_\allowbreak two},
\texttt{double\allowbreak Exp\allowbreak Growing\allowbreak Rec\_\allowbreak super\allowbreak critical\_\allowbreak gt\_\allowbreak two},
\texttt{double\allowbreak Exp\allowbreak Growing\allowbreak Rec\_\allowbreak critical\_\allowbreak two},
\texttt{growGammaSup\_\allowbreak super\allowbreak Double\allowbreak Exp\allowbreak Growing\allowbreak Rec},
\texttt{dimH\_\allowbreak super\allowbreak Double\allowbreak Exp\allowbreak Growing\allowbreak Rec\_\allowbreak zero},
\texttt{dimH\_\allowbreak super\allowbreak Double\allowbreak Exp\allowbreak Growing\allowbreak Rec\_\allowbreak pos\_\allowbreak sub\allowbreak critical},
and \texttt{super\allowbreak Double\allowbreak Exp\allowbreak Growing\allowbreak Rec\_\allowbreak super\allowbreak critical}), using only \texttt{propext},
\texttt{Classical.choice}, and \texttt{Quot.sound} ($0$~\texttt{sorryAx}).

\section{Continued-fraction window constraints and digit residues}\label{sec:windows-residues}

\subsection{Dimensions of continued-fraction window constraints}

Put $H_j=\lfloor Q_j\sqrt j\rfloor$. For irrational $x\in(0,1)$ let
$k(j)$ be its last convergent index with denominator at most $H_j$
(the integer part has index zero) and $r_j=k(j+1)-k(j)$.

\begin{theorem}[Hermite window constraints]\label{thm:windows}
For each prescribed sequence $\epsilon_j\in\{0,1\}$, the set with
$k(j)\equiv\epsilon_j\pmod2$ eventually has Hausdorff dimension one.
The set with eventually bounded $r_j$ has dimension $1/2$.
For every bounded sequence of positive integers $(m_j)$, the set
with $r_j=m_j$ eventually has dimension $1/2$. At a sufficiently
large starting index, either initial parity of $k(j)$ can be imposed.
\end{theorem}

\begin{proof}
Since $Q_{j+1}=(4j+2)Q_j+Q_{j-1}$ gives $Q_{j+1}/Q_j=4j+O(1)$, the window
thresholds $H_j=\lfloor Q_j\sqrt j\rfloor$ satisfy $H_{j+1}/H_j=4j+O(1)\asymp j$
and $\log H_j=j\log j+O(j)$. We work with fundamental cylinders of finite
regular continued fractions: for positive integers $a_1,\dots,a_r\ge1$, let
$I(a_1,\dots,a_r)$ be the closure of the set of irrationals in $(0,1)$ whose
first $r$ partial quotients are $a_1,\dots,a_r$. If $p_{r-1}/q_{r-1}$ and
$p_r/q_r$ are the last two convergents of $[0;a_1,\dots,a_r]$ and we write
$u:=q_r$, $v:=q_{r-1}$, then $I(a_1,\dots,a_r)$ is the closed interval with
endpoints $\frac{p_r}{u}$ and $\frac{p_r+p_{r-1}}{u+v}$, of length
\[
\bigl|I(a_1,\dots,a_r)\bigr|=\frac{1}{u(u+v)}\in\left[\frac{1}{2u^2},\frac{1}{u^2}\right].
\]
At stage $j$ we stop each constructed cylinder at the first convergent index
$\tau_j=k(j)+1$ whose denominator $u=q_{\tau_j}$ crosses $H_j$, maintaining the
two-sided invariant
\[
\frac43 H_j\le u\le \frac52 H_j.
\]
In the construction prescribing parity, the stopped cylinder will also have
last partial quotient $a_{\tau_j}=4$, so that $v=q_{\tau_j-1}$ satisfies
$u=4v+q_{\tau_j-2}>4v$, i.e.\ $v<u/4$. To initialize either construction at any
sufficiently large starting stage $j_0$, choose an integer $d$ with
$H_{j_0}/3\le d\le H_{j_0}/2$: the rational $1/d$ has the two terminal continued
fraction representations $[0;d]$ (length $1$, terminal denominator $d$) and
$[0;d-1,1]$ (length $2$, terminal denominator $(d-1)\cdot1+1=d$). Choosing the
representation whose terminal index has parity $\epsilon_{j_0}$ and appending
the digit $4$ gives a stopped cylinder at index $\tau_{j_0}=k(j_0)+1$ with
$k(j_0)\equiv\epsilon_{j_0}\pmod2$ and denominator
$u\in\{4d+1,4d+(d-1)\}\subset[4H_{j_0}/3,5H_{j_0}/2]$.

\emph{Full-dimensional parity construction.} Fix stage $j\ge j_0$ and a parent
cylinder stopped at $\tau_j=k(j)+1$ with last two denominators $u=q_{\tau_j}$
and $v=q_{\tau_j-1}\le u/4$. Write $T:=H_{j+1}$ and $R:=T/u\asymp j$. We choose
coprime positive integers $(m,n)$ in the rectangle
\[
\mathcal R_R:=\bigl[0.34R,0.36R\bigr]\times\bigl[0.20R,0.22R\bigr]\cap\Z^2.
\]
By M\"obius inversion $\sum_{d\mid\gcd(m,n)}\mu(d)=\mathbf 1_{\gcd(m,n)=1}$, the
number of coprime integer pairs in $\mathcal R_R$ is
\begin{align*}
\#\{(m,n)\in\mathcal R_R:\gcd(m,n)=1\}
&=\sum_{d=1}^{\lfloor0.36R\rfloor}\mu(d)
\left(\left\lfloor\frac{0.36R}{d}\right\rfloor-\left\lceil\frac{0.34R}{d}\right\rceil+1\right)\\
&\qquad\times\left(\left\lfloor\frac{0.22R}{d}\right\rfloor-\left\lceil\frac{0.20R}{d}\right\rceil+1\right)\\
&=\sum_{d=1}^{\lfloor0.36R\rfloor}\mu(d)
\left(\frac{0.02R}{d}+O(1)\right)\left(\frac{0.02R}{d}+O(1)\right)\\
&=0.0004R^2\sum_{d=1}^{\lfloor0.36R\rfloor}\frac{\mu(d)}{d^2}
+O\left(R\sum_{d=1}^{\lfloor0.36R\rfloor}\frac1d\right)\\
&=\frac{0.0004}{\zeta(2)}R^2+O(R\log R),
\end{align*}
since $\sum_{d>\lfloor0.36R\rfloor}d^{-2}=O(R^{-1})$. Thus there are $\gg R^2\asymp j^2$
coprime pairs $(m,n)\in\mathcal R_R$ uniformly for all large $j$.

For each coprime pair $(m,n)\in\mathcal R_R$, we have
$1<\frac{0.34}{0.22}\le\frac{m}{n}\le\frac{0.36}{0.20}=1.8$, so $m/n>1$ has a
standard finite regular continued-fraction expansion $[b_1;b_2,\dots,b_s]$ with
$b_1=\lfloor m/n\rfloor=1$, $s\ge2$, and $b_s\ge2$, as well as the paired
expansion $[b_1;b_2,\dots,b_s-1,1]$ of length $s+1$. By the standard matrix
identity for continuants, if $\begin{pmatrix}m & m'\\ n & n'\end{pmatrix}
=\begin{pmatrix}b_1 & 1\\ 1 & 0\end{pmatrix}\cdots\begin{pmatrix}b_s & 1\\ 1 & 0\end{pmatrix}$,
then replacing $\begin{pmatrix}b_s & 1\\ 1 & 0\end{pmatrix}$ by
$\begin{pmatrix}b_s-1 & 1\\ 1 & 0\end{pmatrix}\begin{pmatrix}1 & 1\\ 1 & 0\end{pmatrix}
=\begin{pmatrix}b_s & b_s-1\\ 1 & 1\end{pmatrix}$ leaves the first column
$\binom{m}{n}$ unchanged. Consequently, when we append either representation of
$m/n$ to the parent cylinder (whose denominator matrix has final row $(u,v)$),
the resulting terminal denominator is
\[
t=mu+nv,
\]
independently of which of the two terminal representations of $m/n$ is chosen
(and every intermediate denominator along the appended tail is strictly less
than $t$, while the penultimate denominator $t'=m'u+n'v$ satisfies $1\le t'<t$).
Using $0.34R\le m\le0.36R$, $0.20R\le n\le0.22R$, $u>0$, $0<v\le u/4$, and
$Ru=T$, we obtain the two-sided inequality
\[
\frac13 T<0.34T=0.34Ru<mu+nv\le\left(0.36+\frac{0.22}{4}\right)Ru=0.415T<\frac12 T.
\]
Thus $t\in(T/3,T/2)$, so the appended tail of $m/n$ has not yet crossed
$T=H_{j+1}$. Choose the representation of $m/n$ (of length $s$ or $s+1$) that
makes the terminal convergent index have parity $\epsilon_{j+1}\pmod2$, and
then append the digit $4$. The new terminal denominator is
\[
u_{\mathrm{new}}=4t+t',
\]
where $0<t'<t$ implies
\[
\frac43 T<4t<u_{\mathrm{new}}=4t+t'<5t<\frac52 T,
\]
and the new penultimate denominator $v_{\mathrm{new}}=t$ satisfies
$v_{\mathrm{new}}<u_{\mathrm{new}}/4$. Thus the first crossing of $T=H_{j+1}$
occurs precisely at the appended digit $4$ (index $\tau_{j+1}=k(j+1)+1$), so the
index immediately preceding the $4$ is indeed $k(j+1)$ and satisfies
$k(j+1)\equiv\epsilon_{j+1}\pmod2$, while $(u_{\mathrm{new}},v_{\mathrm{new}})$
restores the stage-$(j+1)$ invariant. Moreover, distinct coprime pairs
$(m,n)\in\mathcal R_R$ produce stopped child cylinders with disjoint interiors:
each child word is stopped at its first denominator crossing of $T$, so no child
word is a prefix of another, and deleting the parent prefix and the final digit
$4$ recovers the chosen representation of $m/n$ and hence the unique coprime
pair $(m,n)$.

\emph{Prescribed bounded increment construction.} Given a bounded sequence of
positive integers $(m_j)$ ($1\le m_j\le M_*$), start from a stage-$j$ cylinder
stopped at $\tau_j=k(j)+1$ with denominator $4H_j/3\le u\le5H_j/2$ and previous
denominator $0<v<u$. Append $m_j-1$ copies of the digit $1$ (appending no digit
if $m_j=1$), and let $u_*,v_*$ be the resulting last two denominators. By the
Fibonacci continuant relation, $u\le u_*\le F_{m_j+1}u\le F_{M_*+1}u$, so for all
sufficiently large $j$ we have $u_*<T/2$ ($T=H_{j+1}$), meaning none of the
$m_j-1$ ones crosses $T$. Next choose a final digit $a\in\Z_{\ge1}$ such that
the next denominator $u_{\mathrm{new}}=au_*+v_*$ satisfies
\[
\frac43 T\le au_*+v_*\le\frac52 T.
\]
Since $0<v_*<u_*$, this condition holds for all integers $a$ in an interval of
length $(7/6)T/u_*\asymp j$, giving $\gg T/u_*\gg j$ eligible choices of $a$.
For each such $a$, the first crossing of $T=H_{j+1}$ occurs at this final digit
$a$, so $\tau_{j+1}=\tau_j+m_j$, i.e.\ $r_j=k(j+1)-k(j)=\tau_{j+1}-\tau_j=m_j$,
and the new stopped denominator $u_{\mathrm{new}}\in[4T/3,5T/2]$ restores the
inductive hypothesis. Note that $k(j+1)=k(j)+m_j$ is determined thereafter by
the initial parity $k(j_0)\pmod2$ and the cumulative sums of $m_i$; in
particular, taking $m_j\equiv3\pmod2$ alternates the parity of $k(j)$ and
preserves the condition that $k(j)+j$ is odd.

\emph{Mass distribution lower bounds.} In both constructions, assign equal
conditional probability to the retained children of each level-$j$ stopped
cylinder. Let $b=2$ for the full-dimensional parity construction (where each
parent has $\gg R^2\gg j^2$ children) and $b=1$ for the prescribed-count
construction (where each parent has $\gg j$ children). By stage $j$, the
measure $\mu(I_j)$ of any retained stopped cylinder $I_j$ satisfies
\[
\mu(I_j)\le\prod_{i=j_0}^{j-1}\bigl(C_1 i^{-b}\bigr)=\exp\bigl(-bj\log j+O(j)\bigr).
\]
Every retained level-$j$ cylinder $I_j$ has last denominator
$u\in[4H_j/3,5H_j/2]$, so its length $|I_j|=1/(u(u+v))$ satisfies
$(25/2)^{-1}H_j^{-2}\le|I_j|\le(9/16)H_j^{-2}$, and the interiors of distinct
level-$j$ cylinders are disjoint. Consequently, any interval $U\subset(0,1)$ of
length $h\in[H_{j+1}^{-2},H_j^{-2}]$ can intersect the interiors of at most
$\lceil h/((25/2)^{-1}H_j^{-2})\rceil+2\le 15$ level-$j$ cylinders (and its
endpoints meet at most $2$ more). Since
$\log(h^{-1})\le2\log H_{j+1}=2j\log j+O(j)$, for every $0<s<b/2$ we have
\[
\mu(U)\le 17\exp\bigl(-bj\log j+O(j)\bigr)
=O\bigl(\exp(-2s j\log j- (b-2s)j\log j+O(j))\bigr)=O(h^s).
\]
By the mass distribution principle, the Hausdorff dimension is at least $b/2$
 for every $s<b/2$, yielding $\dimH\ge1$ in the first construction (hence $=1$)
and $\dimH\ge1/2$ in the second.

\emph{Upper bound $1/2$ for eventually bounded $r_j$.} If $r_j(x)$ is
eventually bounded, then there exist integers $C\ge1$ and $N_0\ge1$ such that
$k(j)\le Cj$ for all $j\ge N_0$. For any $n\ge2CN_0$, setting
$j=\lfloor n/(2C)\rfloor\ge N_0$ gives $k(j)\le Cj\le n/2<n$, so the $n$th
convergent denominator $d_n(x)=q_n(x)$ has already passed $H_j$:
\[
\frac{\log d_n(x)}{n}\ge\frac{\log H_{\lfloor n/(2C)\rfloor}}{n}
=\frac{\frac{n}{2C}\log\frac{n}{2C}+O(n)}{n}\longrightarrow\infty
\qquad(n\to\infty).
\]
Thus for every $M>0$, every point $x$ with eventually bounded $r_j(x)$ satisfies
$d_n(x)\ge\ee^{Mn}$ for all sufficiently large $n$. Fix any $s>1/2$ and choose
$t\in(1/2,s)$. Since $d_n(a_1,\dots,a_n)\ge\prod_{i=1}^n a_i$, summing
$d_n^{-2t}$ over all $(a_1,\dots,a_n)\in\Z_{\ge1}^n$ gives
\[
\sum_{a_1,\dots,a_n\ge1} d_n(a_1,\dots,a_n)^{-2t}
\le\sum_{a_1,\dots,a_n\ge1}\prod_{i=1}^n a_i^{-2t}=\zeta(2t)^n<\infty.
\]
For the subfamily of level-$n$ cylinders $I(a_1,\dots,a_n)$ with
$d_n\ge\ee^{Mn}$, each cylinder has diameter
$|I(a_1,\dots,a_n)|\le d_n^{-2}\le\ee^{-2Mn}$, and their $s$-dimensional
covering cost is bounded by
\[
\sum_{d_n\ge\ee^{Mn}}\bigl|I(a_1,\dots,a_n)\bigr|^s
\le\sum_{d_n\ge\ee^{Mn}} d_n^{-2s}
\le\ee^{-2M(s-t)n}\sum_{a_1,\dots,a_n\ge1} d_n^{-2t}
\le\bigl(\zeta(2t)\ee^{-2M(s-t)}\bigr)^n.
\]
Choosing $M>\frac{\log\zeta(2t)}{2(s-t)}$ makes the base
$\zeta(2t)\ee^{-2M(s-t)}<1$, so the $s$-cost tends to $0$ as $n\to\infty$.
Taking the countable union over starting indices $N_1$ of the tail sets
$\bigcap_{n\ge N_1}\{d_n\ge\ee^{Mn}\}$ shows that the set of points with
eventually bounded $r_j$ has $\mathcal H^s$-measure zero for all $s>1/2$, hence
Hausdorff dimension at most $1/2$. Together with the lower bound $1/2$, both
sets have exact Hausdorff dimension $1/2$.
\end{proof}

For related spectra of continued-fraction growth see Fan, Liao, Wang
and Wu~\cite{FLLW}. The lower constructions here prescribe the
Hermite windows themselves. Removing the countable set $\Q+\Q\ee$
leaves both dimensions unchanged. These statements concern sets of
variable $x$ and leave the window pattern of the single point $\pi$
undetermined (translation by its integer part preserves denominators).

\subsection{Short continued-fraction windows and divergent window counts}

For $x\in(0,1)\setminus\Q$, retain $H_j,k(j),r_j$ from the window
constraints above.

\begin{theorem}[Short windows at almost every point]\label{thm:short-ae}
There is a constant $c>0$ such that for Lebesgue-almost every $x\in(0,1)$,
\[
 \#\{1\le j\le J:r_j(x)\le1\}\ge c\log J
\]
for all sufficiently large $J$ (with a starting index depending on $x$).
\end{theorem}

\begin{proof}
For irrational $x=[0;a_1,a_2,\dots]\in(0,1)\setminus\Q$, let
$\tau_j(x):=k(j)+1$ be the first convergent index whose denominator $q_{\tau_j}$
strictly exceeds $H_j$. Stop each continued-fraction cylinder at index $\tau_j$
and denote by $\mathcal F_j$ the $\sigma$-field generated by the countable
partition of $(0,1)\setminus\Q$ into stopped cylinders
$I(a_1,\dots,a_{\tau_j})$. Because the stopping condition is determined by the
first crossing of $H_j$, the stopping words at level $j$ are prefix-free and cover
$(0,1)\setminus\Q$. Since $H_j<H_{j+1}$, we have $\tau_j\le\tau_{j+1}$, so
$\mathcal F_j\subseteq\mathcal F_{j+1}$, and the short-window event
\[
E_j:=\bigl\{x\in(0,1)\setminus\Q:r_j(x)=\tau_{j+1}(x)-\tau_j(x)\le1\bigr\}
\]
belongs to $\mathcal F_{j+1}$.

Fix a stopped cylinder $I=I(a_1,\dots,a_r)$ in $\mathcal F_j$ (so $r=\tau_j$),
and write $p/q=p_r/q_r$ and $p'/q'=p_{r-1}/q_{r-1}$ for its last two
convergents, with $q>H_j$ and $0\le q'\le q$. Any irrational $x\in I$ is
uniquely parameterized by its remaining tail $t=[0;a_{r+1},a_{r+2},\dots]\in(0,1)\setminus\Q$
via the M\"obius transformation
\[
x=\Phi_I(t):=\frac{p+p't}{q+q't},\qquad t\in(0,1).
\]
Since $|pq'-p'q|=1$, the Jacobian derivative of $\Phi_I$ is
$|\Phi'_I(t)|=(q+q't)^{-2}$, while the total Lebesgue measure of $I$ is
$\operatorname{Leb}(I)=\int_0^1(q+q't)^{-2}\,dt=\frac{1}{q(q+q')}$. Dividing
$|\Phi'_I(t)|$ by $\operatorname{Leb}(I)$ shows that, conditionally on $I$, the
tail $t\in(0,1)$ has the explicit probability density
\[
p_I(t)=\frac{|\Phi'_I(t)|}{\operatorname{Leb}(I)}=\frac{q(q+q')}{(q+q't)^2},
\qquad 0<t<1.
\]
We now bound $\mathbb P(E_j\mid I)$. There are two cases:
\begin{itemize}[leftmargin=1.8em]
\item If $q>H_{j+1}$, then the denominator $q_{\tau_j}=q$ has already crossed
$H_{j+1}$, so $\tau_{j+1}=\tau_j$ and $r_j=0\le1$ identically on $I$; hence
$\mathbb P(E_j\mid I)=1$.
\item If $H_j<q\le H_{j+1}$, put $A:=\lfloor H_{j+1}/q\rfloor+1\ge2$, so that
$A-1\le H_{j+1}/q<A$. Whenever the next partial quotient $a=a_{r+1}$ satisfies
$a\ge A$, the next denominator is $q_{r+1}=aq+q'\ge Aq>H_{j+1}$, which forces
$\tau_{j+1}=r+1=\tau_j+1$ and hence $r_j=1$ ($E_j$ occurs). In terms of the tail
$t=[0;a,a_{r+2},\dots]\in(0,1)$, the condition $a=\lfloor 1/t\rfloor\ge A$ is
equivalent to $0<t<1/A$ (up to the rational endpoint $1/A$). Integrating the
conditional density $p_I(t)$ over $(0,1/A)$ gives
\[
\mathbb P(E_j\mid I)\ge\mathbb P(a_{r+1}\ge A\mid I)
=\int_0^{1/A}\frac{q(q+q')}{(q+q't)^2}\,dt
=\left[-\frac{q(q+q')}{q'(q+q't)}\right]_{t=0}^{t=1/A}
=\frac{q+q'}{Aq+q'}.
\]
Because $A\ge1$ and $q'\ge0$, we have $A(q+q')=Aq+Aq'\ge Aq+q'$, so
$\frac{q+q'}{Aq+q'}\ge\frac1A$. Moreover, since $q>H_j$ and $H_{j+1}/q\ge1$, we
have $A=\lfloor H_{j+1}/q\rfloor+1\le\frac{H_{j+1}}{q}+1\le\frac{2H_{j+1}}{q}<\frac{2H_{j+1}}{H_j}$,
which yields
\[
\mathbb P(E_j\mid I)\ge\frac{q+q'}{Aq+q'}\ge\frac1A\ge\frac{H_j}{2H_{j+1}}.
\]
\end{itemize}
Since $H_j/(2H_{j+1})<1$, the bound $\mathbb P(E_j\mid I)\ge H_j/(2H_{j+1})$
holds in the first case $q>H_{j+1}$ as well, and therefore holds uniformly on
every atom $I$ of $\mathcal F_j$:
\[
\mathbb P(E_j\mid\mathcal F_j)\ge\frac{H_j}{2H_{j+1}}.
\]
By \eqref{eq:rec} and $H_j=\lfloor Q_j\sqrt j\rfloor$, we have
$H_{j+1}/H_j=4j+O(1)$, so there exist an index $j_0\ge1$ and a constant
$c_0>0$ such that
\[
\mathbb P(E_j\mid\mathcal F_j)\ge\frac{c_0}{j}\qquad\text{almost surely for all }j\ge j_0.
\]

Write $Z_J:=\sum_{j=1}^J\mathbf 1_{E_j}$. Since $E_1,\dots,E_{j-1}$ are
$\mathcal F_j$-measurable, $Z_{j-1}$ is $\mathcal F_j$-measurable, while on binary
indicators $\mathbf 1_{E_j}\in\{0,1\}$ we have the exact identity
$\ee^{-\mathbf 1_{E_j}}=1-(1-\ee^{-1})\mathbf 1_{E_j}$. Using $1-u\le\ee^{-u}$
for all $u\in\R$, the one-step conditional expectation for $j\ge j_0$ satisfies
\[
\mathbb E\bigl[\ee^{-\mathbf 1_{E_j}}\bigm|\mathcal F_j\bigr]
=1-(1-\ee^{-1})\mathbb P(E_j\mid\mathcal F_j)
\le 1-\frac{(1-\ee^{-1})c_0}{j}
\le\exp\left(-\frac{b}{j}\right),\qquad b:=(1-\ee^{-1})c_0>0.
\]
Iterating the tower property $\mathbb E[\ee^{-Z_j}]=\mathbb E\bigl[\ee^{-Z_{j-1}}\mathbb E[\ee^{-\mathbf 1_{E_j}}\mid\mathcal F_j]\bigr]$
from $j=j_0$ to $J$ and using $\sum_{j=j_0}^J\frac1j\ge\log J-\log j_0$ gives
\[
\mathbb E\bigl[\ee^{-Z_J}\bigr]
\le\exp\left(-b\sum_{j=j_0}^J\frac1j\right)
\le C J^{-b},\qquad C:=\ee^{b\log j_0}.
\]
By Markov's inequality applied to the nonnegative random variable $\ee^{-Z_J}$,
\[
\mathbb P\left(Z_J<\frac{b}{2}\log J\right)
=\mathbb P\left(\ee^{-Z_J}>J^{-b/2}\right)
\le J^{b/2}\mathbb E\bigl[\ee^{-Z_J}\bigr]
\le C J^{-b/2}.
\]
Evaluating along the dyadic subsequence $J_m:=2^m$ ($m\ge1$), the probabilities
$\mathbb P(Z_{2^m}<(b/2)\log(2^m))\le C\cdot(2^{-b/2})^m$ form a convergent
geometric series $\sum_{m=1}^\infty C(2^{-b/2})^m<\infty$ since $b>0$. By the
first Borel--Cantelli lemma, for Lebesgue-almost every $x\in(0,1)$ there exists
$m_0(x)\ge2$ such that $Z_{2^m}(x)\ge\frac{b}{2}\log(2^m)$ for all
$m\ge m_0(x)$. Finally, for any integer $J\ge 2^{m_0(x)}$, choose $m\ge m_0(x)$
with $2^m\le J<2^{m+1}$; since $J\mapsto Z_J(x)$ is nondecreasing and
$m\ge\frac{m+1}{2}>\frac{\log_2 J}{2}$ for $m\ge1$, we obtain
\[
Z_J(x)\ge Z_{2^m}(x)\ge\frac{b}{2}\log(2^m)>\frac{b}{4}\log(2^{m+1})>\frac{b}{4}\log J,
\]
proving the theorem with $c:=b/4>0$.
\end{proof}

The argument applies to any increasing integer thresholds tending to
infinity with successive ratios $O(j)$. The estimate concerns
Lebesgue measure on the variable $x$ and supplies a frequency bound
for the short-window criterion of the Phase-1 project.

\begin{theorem}[Full dimension for divergent window counts]\label{thm:no-short-dim}
Let $\varphi=(1+\sqrt5)/2$ and prescribe arbitrary
$\epsilon_j\in\{0,1\}$. For any $0<\eta<1/\log\varphi$, the set of
$x\in(0,1)$ satisfying both $r_j(x)\ge\eta\log j$ and
$k(j)\equiv\epsilon_j\pmod2$ eventually has Hausdorff dimension at
least $1-\eta\log\varphi$. The set satisfying the prescribed
eventual parity and $r_j(x)\to\infty$ has Hausdorff dimension one
and Lebesgue measure zero.
\end{theorem}

\begin{proof}
We modify the stopped-cylinder construction of Theorem~\ref{thm:windows},
initialized at a large stage $j_0$ with the prescribed parity
$k(j_0)\equiv\epsilon_{j_0}\pmod2$. At stage $j\ge j_0$, a retained cylinder ends
at index $\tau_j=k(j)+1$ with last partial quotient $a_{\tau_j}=4$ and last two
denominators $u=q_{\tau_j}$ and $v=q_{\tau_j-1}$, where
\[
\frac43 H_j\le u\le\frac52 H_j,\qquad 0<v\le\frac{u}{4}.
\]
First append $\ell_j:=\lceil\eta\log j\rceil$ partial quotients equal to $1$.
Write $u':=q_{\tau_j+\ell_j}$ and $v':=q_{\tau_j+\ell_j-1}$ for the resulting
last two denominators. Since multiplication by $\begin{pmatrix}1 & 1\\ 1 & 0\end{pmatrix}^{\ell_j}
=\begin{pmatrix}F_{\ell_j+1} & F_{\ell_j}\\ F_{\ell_j} & F_{\ell_j-1}\end{pmatrix}$
governs the denominator recurrence during a block of $\ell_j$ ones (where
$F_0=0,F_1=1,F_{i+1}=F_i+F_{i-1}$ are the Fibonacci numbers), we have
\[
\begin{pmatrix}u'\\ v'\end{pmatrix}
=\begin{pmatrix}F_{\ell_j+1} & F_{\ell_j}\\ F_{\ell_j} & F_{\ell_j-1}\end{pmatrix}
\begin{pmatrix}u\\ v\end{pmatrix},
\]
so $0<v'\le u'$ and, using $v\le u$ and $F_{\ell_j+1}+F_{\ell_j}=F_{\ell_j+2}\le\varphi^{\ell_j+1}$,
\[
u\le u'=F_{\ell_j+1}u+F_{\ell_j}v\le F_{\ell_j+2}u\le\varphi^{\ell_j+1}u
\le\varphi^2 j^{\eta\log\varphi}u.
\]
With $T:=H_{j+1}$ and $H_{j+1}/H_j=4j+O(1)$, the remaining scale ratio
$R:=T/u'$ satisfies
\[
R=\frac{T}{u'}\ge\frac{H_{j+1}}{\varphi^2 j^{\eta\log\varphi}\cdot(5H_j/2)}
\ge c_1 j^{1-\eta\log\varphi}
\]
for a constant $c_1>0$. Since $0<\eta<1/\log\varphi$, the exponent
$\gamma:=1-\eta\log\varphi>0$ is strictly positive, so $R\to\infty$ uniformly as
$j\to\infty$; in particular, $u'=T/R<T$ for all large $j$, so the prelude of
$\ell_j$ ones stays strictly below $T$.

Next, choose coprime positive integers $(m,n)$ in the rectangle
\[
\mathcal R'_R:=\bigl[0.34R,0.35R\bigr]\times\bigl[0.12R,0.13R\bigr]\cap\Z^2.
\]
By the same M\"obius inversion calculation as in Theorem~\ref{thm:windows}, the
number of coprime pairs $(m,n)\in\mathcal R'_R$ is
$\frac{0.01\times0.01}{\zeta(2)}R^2+O(R\log R)\gg R^2\ge c_2 j^{2\gamma}$. For
every $(m,n)\in\mathcal R'_R$, the quotient $m/n$ satisfies
\[
2<\frac{34}{13}=\frac{0.34R}{0.13R}\le\frac{m}{n}\le\frac{0.35R}{0.12R}=\frac{35}{12}<3.
\]
Thus $m/n\in(2,3)$ has first partial quotient $\lfloor m/n\rfloor=2$ and two
finite regular continued-fraction representations of opposite lengths modulo
two (both of length at least $2$). Append the digits of $m/n$ using the
representation whose terminal convergent index has parity
$\epsilon_{j+1}\pmod2$. As shown in Theorem~\ref{thm:windows}, both
representations produce the same terminal denominator
\[
t=mu'+nv'
\]
(and all intermediate denominators along the tail of $m/n$ are at most $t$).
Using $0.34R\le m\le0.35R$, $0.12R\le n\le0.13R$, $0<v'\le u'$, and $Ru'=T$, we
obtain
\[
\frac13 T<0.34T=0.34Ru'<t=mu'+nv'\le(0.35+0.13)Ru'=0.48T<\frac12 T.
\]
Thus $t\in(T/3,T/2)$ stays below $T=H_{j+1}$. Now append the final digit $4$.
The new denominator is $u_{\mathrm{new}}=4t+t'$ with $0<t'<t$, so
\[
\frac43 T<4t<u_{\mathrm{new}}=4t+t'<5t<\frac52 T
\]
and $v_{\mathrm{new}}=t<u_{\mathrm{new}}/4$. Hence the first crossing of
$T=H_{j+1}$ occurs at this final digit $4$ (index $\tau_{j+1}=k(j+1)+1$), so the
index immediately before the $4$ is $k(j+1)$ with
$k(j+1)\equiv\epsilon_{j+1}\pmod2$, and the window increment satisfies
\[
r_j=\tau_{j+1}-\tau_j=\ell_j+\operatorname{len}(m/n)+1\ge\ell_j+3>\eta\log j.
\]
Distinct coprime pairs $(m,n)\in\mathcal R'_R$ produce stopped child cylinders
with disjoint interiors because each child stops at its first crossing of $T$,
and deleting the parent prefix, the fixed prelude of $\ell_j$ ones, and the
terminal digit $4$ uniquely recovers the chosen representation of $m/n$.

Since each stage-$j$ parent has at least $c_2 j^{2\gamma}$ children, assigning
equal conditional masses gives each level-$j$ stopped cylinder a measure at most
$\prod_{i=j_0}^{j-1}(c_2^{-1}i^{-2\gamma})=\exp(-2\gamma j\log j+O(j))$, while
its length is comparable to $H_j^{-2}$. As in Theorem~\ref{thm:windows}, any
interval of length $h\in[H_{j+1}^{-2},H_j^{-2}]$ meets at most $17$ level-$j$
cylinders, and since $\log H_{j+1}=j\log j+O(j)$, its measure is $O(h^s)$ for
every $0<s<\gamma$. The mass distribution principle yields Hausdorff dimension
at least $\gamma=1-\eta\log\varphi$. Finally, every point in the constructed
set for a given $\eta>0$ satisfies $r_j(x)\ge\eta\log j\to\infty$ and
$k(j)\equiv\epsilon_j\pmod2$ eventually, so the set with $r_j(x)\to\infty$ and
prescribed eventual parity contains these sets for all $\eta\in(0,1/\log\varphi)$;
letting $\eta\downarrow0$ shows that it has Hausdorff dimension~$1$, while
Theorem~\ref{thm:short-ae} implies that it has Lebesgue measure zero (since
almost every $x$ has $r_j(x)\le1$ infinitely often).
\end{proof}

\begin{remark}[Machine-checked counting core, Frostman assembly, CF tree realizations, regular CF bridge, bounded-window upper bound, and almost-everywhere short windows for Theorems~\ref{thm:windows}--\ref{thm:no-short-dim}]\label{rem:windows-lean}
The M\"obius counting cores, the stopped-tree Frostman dimension assembly, all three concrete continued-fraction stopped-cylinder tree realizations, the pointwise regular continued-fraction bridge into mathlib's \texttt{GenContFract.of}, the $\dimH\le1/2$ upper bound for eventually bounded window counts, and the measure-theoretic instantiation of Theorem~\ref{thm:short-ae} together with the Lebesgue-nullity deduction of Theorem~\ref{thm:no-short-dim} (completing all of Theorems~\ref{thm:windows}, \ref{thm:short-ae}, and~\ref{thm:no-short-dim} at Layer~A) have been verified by the Lean~4 kernel (Lean~4.35.0 with mathlib~v4.35.0):
(i)~the M\"obius inversion identity \texttt{WindowMobiusCount.\allowbreak coprime\_\allowbreak count\_\allowbreak eq}, the partial-sum lower bound $359/900$ (\texttt{moebius\_\allowbreak partial}), and the rectangle-counting lower bounds $R^2/10000$ (\texttt{WindowMobiusCount.\allowbreak coprime\_\allowbreak count\_\allowbreak ge} on $[0.34R,0.36R]\times[0.20R,0.22R]$) and $R^2/100000$ (\texttt{NoShortDim.\allowbreak coprimeS\_\allowbreak count\_\allowbreak ge} on $[0.34R,0.35R]\times[0.12R,0.13R]$) for $R\ge10^6$;
(ii)~the scale growth and inverse-square decay of $H_j=\lfloor Q_j\sqrt{j}\rfloor$ (\texttt{WindowPaperScale.\allowbreak Hpaper\_\allowbreak inv\_\allowbreak sq\_\allowbreak tendsto}, \texttt{Hpaper\_\allowbreak log\_\allowbreak le}), the uniform-branching pruning and Cantor double-counting lemmas (\texttt{WindowTreePrune.*}), the real-exponent asymptotic Frostman threshold (\texttt{WindowAsympThresh.\allowbreak asymp\_\allowbreak thresh\_\allowbreak exists\_\allowbreak rpow}), and the full per-window Frostman discharge (\texttt{WindowFrostmanDischarge.\allowbreak frostman\_\allowbreak of\_\allowbreak uniform\_\allowbreak tree}, \texttt{dimH\_\allowbreak ge\_\allowbreak of\_\allowbreak stopped\_\allowbreak tree\_\allowbreak rpow}) instantiated at the paper scale~$H_j$ in \texttt{WindowPaperScale.\allowbreak dimH\_\allowbreak ge\_\allowbreak one\_\allowbreak windows\_\allowbreak paper\_\allowbreak scale} ($b=2$, $\dimH\ge 1$), \texttt{dimH\_\allowbreak ge\_\allowbreak half\_\allowbreak prescribed\_\allowbreak paper\_\allowbreak scale} ($b=1$, $\dimH\ge 1/2$), and \texttt{dimH\_\allowbreak ge\_\allowbreak gamma\_\allowbreak no\_\allowbreak short\_\allowbreak paper\_\allowbreak scale} / \texttt{dimH\_\allowbreak ge\_\allowbreak one\_\allowbreak of\_\allowbreak gamma\_\allowbreak family} ($\dimH\ge 1-\eta\log\varphi\uparrow 1$), together with concrete tree- and measure-based non-vacuity witnesses (\texttt{WindowPaperScale.\allowbreak nonvacuous\_\allowbreak factorial\_\allowbreak tree\_\allowbreak dimH}, \texttt{nonvacuous\_\allowbreak factorial\_\allowbreak measure\_\allowbreak dimH}) and the local-scale / cascade measure bounds (\texttt{WindowFrostmanDischarge.\allowbreak mass\_\allowbreak distrib\_\allowbreak of\_\allowbreak local\_\allowbreak diam\_\allowbreak bound}, \texttt{mu\_\allowbreak window\_\allowbreak cascade\_\allowbreak full}, \texttt{dimH\_\allowbreak ge\_\allowbreak of\_\allowbreak mu\_\allowbreak window\_\allowbreak cascade});
(iii)~the concrete continued-fraction stopped-cylinder tree \texttt{WindowCFRealize.\allowbreak cfWindowsTree} for Theorem~\ref{thm:windows} (constructed via \texttt{euclidCF}, parity twin \texttt{flipLast}, and cylinder geometry \texttt{cylPair}), discharging every tree-interface hypothesis unconditionally (\texttt{WindowCFRealize.\allowbreak dimH\_\allowbreak ge\_\allowbreak one\_\allowbreak cfWindowsTree});
(iv)~the pointwise regular continued-fraction bridge in \texttt{WindowCFBridge.lean} ($20$ audited declarations), proving strict cylinder nesting (\texttt{cylPair\_\allowbreak nest\_\allowbreak strict}), open-cylinder M\"obius inversion (\texttt{exists\_\allowbreak mobius\_\allowbreak of\_\allowbreak mem\_\allowbreak Ioo\_\allowbreak cylPair}), exact partial-quotient and convergent-denominator identification with mathlib's \texttt{GenContFract.of} (\texttt{stream\_\allowbreak mobius\_\allowbreak spec}, \texttt{contsAux\_\allowbreak mobius\_\allowbreak eq}, \texttt{dens\_\allowbreak mobius\_\allowbreak eq}, \texttt{dens\_\allowbreak of\_\allowbreak mono}), irrationality and regular CF stopping index \texttt{cfStopIdx} (defined directly against \texttt{NoShortDim.H j} $=\lfloor Q_j\sqrt{j}\rfloor$) matching for any stopped CF word tree (\texttt{stopped\_\allowbreak word\_\allowbreak tree\_\allowbreak cf\_\allowbreak bridge}), and the end-to-end Theorem~\ref{thm:windows} exact parity dimension $\dimH(\texttt{prescribedParitySet}\,\epsilon)=1$ (\texttt{WindowCFBridge.\allowbreak dimH\_\allowbreak prescribedParitySet\_\allowbreak eq\_\allowbreak one});
(v)~the concrete prelude-of-ones CF tree \texttt{WindowNoShortCFRealize.\allowbreak cfNoShortTree} for Theorem~\ref{thm:no-short-dim} ($18$ audited declarations, discharging \texttt{dimH\_\allowbreak ge\_\allowbreak gamma\_\allowbreak no\_\allowbreak short\_\allowbreak paper\_\allowbreak scale} and proving \texttt{dimH\_\allowbreak ge\_\allowbreak gamma\_\allowbreak noShortSet} and exact $\dimH(\texttt{divWindowSet}\,\epsilon)=1$ in \texttt{dimH\_\allowbreak divWindowSet\_\allowbreak eq\_\allowbreak one} unconditionally via $\eta\downarrow0$) and the concrete $b=1$ prescribed-count CF tree \texttt{WindowPrescribedCFRealize.\allowbreak cfPrescribedTree} for Theorem~\ref{thm:windows} ($16$ audited declarations, discharging \texttt{dimH\_\allowbreak ge\_\allowbreak half\_\allowbreak prescribed\_\allowbreak paper\_\allowbreak scale} and proving \texttt{prescribedIncSetAt\_\allowbreak disjoint}, \texttt{dimH\_\allowbreak ge\_\allowbreak half\_\allowbreak prescribedIncSetAt} / \texttt{exists\_\allowbreak start\_\allowbreak dimH\_\allowbreak ge\_\allowbreak half\_\allowbreak prescribedIncSetAt} for both disjoint initial parities at a fixed starting index, \texttt{dimH\_\allowbreak ge\_\allowbreak half\_\allowbreak prescribedIncSet}, and \texttt{dimH\_\allowbreak ge\_\allowbreak half\_\allowbreak boundedIncSet} unconditionally);
(vi)~the bounded-window upper bound and exact $1/2$-dimension closure in \texttt{WindowBoundedUpper.lean} ($21$ audited declarations), proving finite-word denominator product and length lower bounds (\texttt{cfProd\_\allowbreak den\_\allowbreak ge\_\allowbreak length}, \texttt{cfProd\_\allowbreak den\_\allowbreak ge\_\allowbreak prod}), the factorial growth bound on~$H_j$ (\texttt{Hpaper\_\allowbreak ge\_\allowbreak fact}), super-exponential denominator growth on \texttt{boundedIncSet} (\texttt{boundedIncSet\_\allowbreak dens\_\allowbreak exp\_\allowbreak growth}), level-$n$ cylinder covering and $\zeta(2t)^n$ \texttt{ENNReal} \texttt{tsum} bounds (\texttt{exists\_\allowbreak word\_\allowbreak cover\_\allowbreak of\_\allowbreak irrational}, \texttt{tsum\_\allowbreak cfProd\_\allowbreak den\_\allowbreak rpow\_\allowbreak le}), $\mathcal H^s$-nullity for every $s>1/2$ (\texttt{hausdorffMeasure\_\allowbreak boundedIncSet\_\allowbreak zero}), and the exact identities \texttt{dimH\_\allowbreak boundedIncSet\_\allowbreak le\_\allowbreak half}, \texttt{dimH\_\allowbreak boundedIncSet\_\allowbreak eq\_\allowbreak half}, \texttt{dimH\_\allowbreak prescribedIncSet\_\allowbreak eq\_\allowbreak half}, and \texttt{dimH\_\allowbreak prescribedIncSetAt\_\allowbreak eq\_\allowbreak half}, completing Layer~A verification of all of Theorem~\ref{thm:windows};
(vii)~the full measure-theoretic realization of Theorem~\ref{thm:short-ae} and the Lebesgue-nullity deduction in Theorem~\ref{thm:no-short-dim} in \texttt{WindowShortAeRealize.lean} ($14$ audited declarations): prefix-free stopped-word cylinder partitions of $(0,1)\setminus\Q$, constancy of \texttt{cfStopIdx} on stopped cylinders, measurability of the short-window events $E_j$, the explicit atom conditional volume lower bound $1/(40j)$ for $j\ge3$, the exponential step inequality yielding the explicit constant $c=(1-\ee^{-1})/160>0$ in \texttt{thm\_\allowbreak short\_\allowbreak ae\_\allowbreak explicit} and \texttt{thm\_\allowbreak short\_\allowbreak ae}, and the complete conjunction $\dimH(\texttt{divWindowSet}\,\epsilon)=1\wedge\operatorname{Leb}(\texttt{divWindowSet}\,\epsilon)=0$ in \texttt{volume\_\allowbreak divWindowSet\_\allowbreak eq\_\allowbreak zero} and \texttt{thm\_\allowbreak no\_\allowbreak short\_\allowbreak dim\_\allowbreak complete}.
All declarations in these modules depend only on the Lean base axioms \texttt{propext}, \texttt{Classical.choice}, and \texttt{Quot.sound}, as audited with \texttt{\#print axioms} in \texttt{Audit.lean} (see \texttt{verification/\allowbreak audit\_\allowbreak final\_\allowbreak 20261122.log}, $872$ audited declarations across $131$ Lean files with zero \texttt{sorryAx} and zero project-added axioms).
\end{remark}

In particular, the set with eventually $r_j\ge3$ has Lebesgue
measure zero and full Hausdorff dimension~$1$. The same is true if one
also imposes any prescribed parity of the terminal convergent
indices.

\subsection{Eventual periodicity of digit residues}

For irrational $x=[0;a_1,a_2,\ldots]$ in $(0,1)$ let $E$ consist of
the points for which $(a_i\bmod4)$ is eventually periodic. Let $F$
consist of points with every digit in $\mathcal A=\{1,5,9,\ldots\}$.
Continued fractions with a fixed arithmetic-progression alphabet are
studied by Chousionis, Leykekhman and Urba\'nski~\cite{CLU}. Here
we compare different sequences of allowed residues, including sequences
that change from one position to the next.

\begin{theorem}[Digit residues modulo four]\label{thm:residues}
$\dimH E=\dimH F$, with $357/512<\dimH F<179/256$; the closed bracket
$357/512\le\dimH F\le179/256$ is machine-checked in Lean~4 (Remark~\ref{rem:residue-kernel}).
\end{theorem}

\begin{proof}
\emph{Step 1: Uniform Lipschitz equivalence and $\dimH E=\dimH F$.}
Fix any sequence of residue shifts $\mathbf c=(c_1,c_2,\dots)\in\{0,1,2,3\}^{\N}$,
and define the digit-shift map $\Psi_{\mathbf c}:F\to(0,1)\setminus\Q$ by
\[
\Psi_{\mathbf c}\bigl([0;a_1,a_2,\dots]\bigr):=[0;a_1+c_1,a_2+c_2,\dots]
\qquad(a_i\in\mathcal A=\{1,5,9,\dots\}).
\]
We prove that $\Psi_{\mathbf c}$ is Lipschitz on $F$ with Lipschitz constant
$40/3$, \emph{uniformly} over all shift sequences $\mathbf c$. Let
$x=[0;a_1,a_2,\dots]$ and $y=[0;b_1,b_2,\dots]$ be two distinct points of $F$,
and let $n\ge0$ be the length of their longest common prefix (so $a_i=b_i$ for
$1\le i\le n$ and $a_{n+1}\ne b_{n+1}$). By symmetry we may assume
$a:=a_{n+1}<b:=b_{n+1}$. Since $a,b\in\mathcal A=\{1,5,9,\dots\}$ are distinct
positive integers congruent to $1\bmod4$, we have $a\ge1$ and $b-a\ge4$ (in
particular $b\ge5$). Write the source tails starting at position $n+1$ as
\[
x_t:=[0;a,x_1]=\frac{1}{a+x_1},\qquad
y_t:=[0;b,y_1]=\frac{1}{b+y_1}\qquad\bigl(x_1,y_1\in(0,1)\bigr),
\]
and the shifted image tails as
\begin{align*}
x'_t&:=[0;a+c,x'_1]=\frac{1}{a+c+x'_1},\\
y'_t&:=[0;b+c,y'_1]=\frac{1}{b+c+y'_1}\qquad\bigl(c:=c_{n+1}\in\{0,1,2,3\},\ x'_1,y'_1\in(0,1)\bigr).
\end{align*}
Since $0<x_1,y_1<1$ and $a+1<b$, the source tail distance is bounded below by
\[
|x_t-y_t|=\frac{1}{a+x_1}-\frac{1}{b+y_1}
>\frac{1}{a+1}-\frac{1}{b}=\frac{b-a-1}{(a+1)b}>0.
\]
Similarly, since $0<x'_1,y'_1<1$ and $c\ge0$, the shifted tail distance is
bounded above by
\[
|x'_t-y'_t|=\frac{1}{a+c+x'_1}-\frac{1}{b+c+y'_1}
<\frac{1}{a+c}-\frac{1}{b+c+1}=\frac{b-a+1}{(a+c)(b+c+1)}
\le\frac{b-a+1}{a(b+1)}.
\]
Combining these two inequalities bounds the tail distance ratio by
\[
\frac{|x'_t-y'_t|}{|x_t-y_t|}
\le\frac{b-a+1}{b-a-1}\cdot\frac{(a+1)b}{a(b+1)}.
\]
Put $d:=b-a\ge4$. Then $\frac{d+1}{d-1}=1+\frac{2}{d-1}$ is strictly decreasing
in $d\ge4$, so $\frac{b-a+1}{b-a-1}\le\frac{4+1}{4-1}=\frac53$. Moreover,
since $b\ge a+4$ and $a\ge1$, we have
\[
\frac{(a+1)b}{a(b+1)}
=\frac{a+1}{a}\cdot\frac{b}{b+1}
<2\cdot1=2
\quad\left(\text{and in fact }\frac{b-a+1}{b-a-1}\frac{(a+1)b}{a(b+1)}\le\frac53\cdot\frac{2\cdot5}{1\cdot6}=\frac{50}{18}<\frac{10}{3}\right),
\]
so in all cases
\[
|x'_t-y'_t|\le\frac{10}{3}|x_t-y_t|.
\]
Now restore the common prefix of length $n$. For the source points,
$x=M(x_t)$ and $y=M(y_t)$, where $M(t)=\frac{p_n+p_{n-1}t}{q_n+q_{n-1}t}$ is the
M\"obius map of the common prefix $(a_1,\dots,a_n)$, with last denominator
$q:=q_n\ge1$ and previous denominator $0\le q_{n-1}\le q$. For all $t\in[0,1]$,
$q\le q+q_{n-1}t\le 2q$, so the derivative $|M'(t)|=(q+q_{n-1}t)^{-2}$ satisfies
\[
\frac{1}{4q^2}\le |M'(t)|\le\frac{1}{q^2}\qquad\text{for all }t\in[0,1].
\]
By the mean value theorem, $|x-y|=|M(x_t)-M(y_t)|\ge\frac{1}{4q^2}|x_t-y_t|$.
For the shifted points, $\Psi_{\mathbf c}(x)=M'(x'_t)$ and
$\Psi_{\mathbf c}(y)=M'(y'_t)$, where $M'$ is the M\"obius map of the shifted
prefix $(a_1+c_1,\dots,a_n+c_n)$ with last denominator $q_{\mathrm{new}}$.
Because the continuant $q_n(d_1,\dots,d_n)$ is strictly increasing in each
partial quotient $d_i\ge1$ and $a_i+c_i\ge a_i$, we have
$q_{\mathrm{new}}\ge q$, and the mean value theorem gives
\[
\bigl|\Psi_{\mathbf c}(x)-\Psi_{\mathbf c}(y)\bigr|
=\bigl|M'(x'_t)-M'(y'_t)\bigr|
\le\frac{1}{q_{\mathrm{new}}^2}|x'_t-y'_t|
\le\frac{1}{q^2}|x'_t-y'_t|.
\]
Combining the prefix and tail bounds yields the uniform Lipschitz inequality
\[
\bigl|\Psi_{\mathbf c}(x)-\Psi_{\mathbf c}(y)\bigr|
\le\frac{1}{q^2}\cdot\frac{10}{3}|x_t-y_t|
\le 4\cdot\frac{10}{3}|x-y|=\frac{40}{3}|x-y|\qquad\text{for all }x,y\in F.
\]
Hence $\dimH\Psi_{\mathbf c}(F)\le\dimH F$ for every shift sequence
$\mathbf c\in\{0,1,2,3\}^{\N}$. Conversely, every point
$x=[0;d_1,d_2,\ldots]\in E$ has $(d_i\bmod4)$ eventually periodic; setting
$c_i:=(d_i-1)\bmod4\in\{0,1,2,3\}$ gives an eventually periodic sequence
$\mathbf c=(c_1,c_2,\dots)\in\{0,1,2,3\}^{\N}$ such that
$a_i:=d_i-c_i=4\lfloor(d_i-1)/4\rfloor+1\in\mathcal A$ for every $i\ge1$, so
$x=\Psi_{\mathbf c}([0;a_1,a_2,\dots])\in\Psi_{\mathbf c}(F)$. Thus $E$ is the
union of $\Psi_{\mathbf c}(F)$ over the countable set of eventually periodic
sequences $\mathbf c\in\{0,1,2,3\}^{\N}$ (each determined by a preperiod length,
a period length, and finitely many residues in $\{0,1,2,3\}$). Countable
stability gives $\dimH E\le\dimH F$, while $F\subset E$ (corresponding to
$c_i\equiv0$) gives $\dimH F\le\dimH E$. Therefore $\dimH E=\dimH F$.

\emph{Step 2: Transfer operator, Bowen's formula, and integral tail bounds.}
For the numerical bounds on $\dimH F$, let $f_a(x)=1/(a+x)$ for $0\le x\le1$.
For an alphabet $\mathcal D\subseteq\mathcal A$ and $s>1/2$ define the
Perron--Frobenius transfer operator on bounded positive functions $v:[0,1]\to(0,\infty)$
by
\[
 (\mathcal L_{s,\mathcal D}v)(x)
   =\sum_{a\in\mathcal D}(a+x)^{-2s}v(f_a(x)).
\]
The topological pressure $P(s,\mathcal D)$ is the limit of
$n^{-1}\log(\mathcal L_{s,\mathcal D}^n1)(0)$ as $n\to\infty$.
Indeed, the derivative of an $n$-letter cylinder map
$f_{a_1}\circ\cdots\circ f_{a_n}(x)$ has magnitude $(q_n+xq_{n-1})^{-2}$, which
is maximal at $x=0$ and varies by at most a factor of $4$ on $[0,1]$. The
pressure $s\mapsto P(s,\mathcal D)$ is continuous and strictly decreasing on
$(1/2,\infty)$, and Bowen's formula identifies its unique zero with the
Hausdorff dimension of the limit set coded by $\mathcal D$
\cite[\S2, Theorem~2.3]{CLU}. To satisfy the strict uniform contraction
convention in that reference (since $f_1'(0)=-1$), one passes to the system of
all two-letter branches $f_a\circ f_b$ ($a,b\in\mathcal D$): their derivatives
satisfy $|(f_a\circ f_b)'(x)|=(ab+ax+1)^{-2}\le(1\cdot1+0+1)^{-2}=1/4<1$ on
$[0,1]$, their cylinder interiors are disjoint, their transfer operator is
$\mathcal L_{s,\mathcal D}^2$ (so its pressure is $2P(s,\mathcal D)$, with the
same zero), and the limit set coded by paired words in $\mathcal D^2$ is
identical to the limit set coded by $\mathcal D$.

For any bounded positive function $v:[0,1]\to(0,\infty)$ with
$\inf_{[0,1]}v>0$, a uniform pointwise upper inequality
$\mathcal L_{s,\mathcal D}v\le(1-\eta)v$ on $[0,1]$ ($0<\eta<1$) forces
$P(s,\mathcal D)\le\log(1-\eta)<0$: by positivity of $\mathcal L_{s,\mathcal D}$,
iteration gives $\mathcal L_{s,\mathcal D}^n v\le(1-\eta)^n v$, and comparing $1$
with $v/\inf v$ yields
$(\mathcal L_{s,\mathcal D}^n1)(0)\le(1-\eta)^n v(0)/\inf v$.
Likewise, a uniform pointwise lower inequality
$\mathcal L_{s,\mathcal D}v\ge(1+\eta)v$ on $[0,1]$ ($\eta>0$) forces
$(\mathcal L_{s,\mathcal D}^n1)(0)\ge(1+\eta)^n v(0)/\sup v$ and hence
$P(s,\mathcal D)\ge\log(1+\eta)>0$. We apply the lower test at $s=357/512$ to
the finite subalphabet $\mathcal D_N:=\{4k+1:0\le k<N\}$ with $N=2^{80}$ (whose
limit set is contained in $F$, so $\dimH F\ge\dimH F_{\mathcal D_N}>357/512$),
and the upper test at $s=179/256$ to the full alphabet $\mathcal D=\mathcal A$
(certifying $\dimH F<179/256$).

In both tests, $v$ is chosen as a positive step function on the $512$ uniform
grid cells $C_j=[l_j,u_j)=[j/512,(j+1)/512)$ ($0\le j\le510$) and
$C_{511}=[511/512,1]$. On a cell $[l,u]$, for each explicit branch $a=4k+1$
($0\le k<130$) and $x\in[l,u]$, the weight $(a+x)^{-2s}$ is decreasing in $x$
and thus lies in $[(a+u)^{-2s},(a+l)^{-2s}]$, while $f_a(x)=1/(a+x)$ lies in the
closed image interval $[f_a(u),f_a(l)]=[1/(a+u),1/(a+l)]$; we bound $v(f_a(x))$
from below and above by the minimum and maximum of $v$ over all grid cells
intersected by $[f_a(u),f_a(l)]$. For all tail indices $k\ge130$, the digit is
$a=4k+1\ge521>512$, so for every $x\in[0,1]$ the image point satisfies
$0<f_a(x)=1/(4k+1+x)\le 1/521<1/512$ and therefore falls strictly inside the
first cell $C_0=[0,1/512)$, where $v$ takes the constant value $v_0$. Putting
$p:=2s>1$, we bound the tail sums over $k\ge130$ by monotone integral comparison:
for every $x\in[0,1]$ and every $t\in[k,k+1]$, we have
$4k+1+x\le 4k+2\le 4t+2$, so
\[
(4k+1+x)^{-p}\ge(4k+2)^{-p}\ge\int_k^{k+1}(4t+2)^{-p}\,dt,
\]
and summing from $k=130$ to $N-1$ (where $4\cdot130+2=522$) gives the lower tail
bound
\[
 \sum_{k=130}^{N-1}(4k+1+x)^{-p}v_0
 \ge v_0\int_{130}^N(4t+2)^{-p}\,dt
 =v_0\left[\frac{(4t+2)^{1-p}}{4(1-p)}\right]_{t=130}^{t=N}
 =v_0\,\frac{522^{1-p}-(4N+2)^{1-p}}{4(p-1)}.
\]
In the opposite direction, for every $x\in[0,1]$ and every $t\in[k-1,k]$, we
have $4k+1+x\ge 4k+1\ge 4t+1$, so
\[
(4k+1+x)^{-p}\le(4k+1)^{-p}\le\int_{k-1}^k(4t+1)^{-p}\,dt,
\]
and summing from $k=130$ to $\infty$ (where $4\cdot129+1=517$) gives the upper
tail bound
\[
 \sum_{k=130}^{\infty}(4k+1+x)^{-p}v_0
 \le v_0\int_{129}^\infty(4t+1)^{-p}\,dt
 =v_0\left[\frac{(4t+1)^{1-p}}{4(1-p)}\right]_{t=129}^{t=\infty}
 =v_0\,\frac{517^{1-p}}{4(p-1)} .
\]

\emph{Step 3: Exact integer-arithmetic verification.}
The companion script \texttt{residue\_\allowbreak certificate\_\allowbreak general\_\allowbreak N512.py} constructs a
positive integer vector of cell values for each test. Every branch
weight and integral power is enclosed by rational endpoints. In
particular, to enclose $S(y/z)^{-m/d}$, where $S=10^{13}$ and
$y,z,m,d$ are positive integers, the script checks the exact
integer inequalities $L^d y^m\le S^d z^m\le U^d y^m$ for the proposed
integer endpoints $L,U$. Summation and comparison against the
corresponding cell value use exact integer arithmetic. At $s=357/512$, the
smallest relative lower margin $\min_j((\mathcal L_{s,\mathcal D_N}v-v)/v)|_{C_j}$
equals
$47991263886531665573/9862965570000000000000>48/10000>0$; at $s=179/256$, the
smallest relative upper margin $\min_j((v-\mathcal L_{s,\mathcal A}v)/v)|_{C_j}$
equals
$1350497164217210109/4931093915000000000000>27/100000>0$. The
strict signs certify the stated open interval $(357/512,179/256)$, of width
$1/512$. The same integer inequalities, re-encoded as Boolean checks, are
evaluated by the Lean kernel (Remark~\ref{rem:residue-kernel}).
\end{proof}

As a comparison, Table~1 of \cite{CLU} bounds the dimension for
$\{1,4,7,\ldots\}$ by $0.743586$. Its Proposition~2.7 transfers this upper bound
to $\{1,5,9,\ldots\}$ by the increasing bijection $1+3k\mapsto 1+4k$.

\begin{theorem}[General-modulus residue sandwich, transfer-operator soundness, and certified $m\in\{2,3,4,5,6\}$ brackets]\label{thm:residues-general}
Let $m\ge2$ be an integer modulus. For every sequence $\mathbf r=(r_i)_{i\ge1}$
with $1\le r_i\le m$, write
\[
F_m(\mathbf r)
:=\bigl\{x=[0;a_1,a_2,\dots]\in(0,1)\setminus\Q:\ a_i\ge r_i\text{ and }a_i\equiv r_i\pmod m\ \forall i\ge1\bigr\},
\]
put $F_{m,r}:=F_m(r,r,\dots)$ for $1\le r\le m$, and let $E_m\subset(0,1)\setminus\Q$
be the set of irrationals whose partial quotients $(a_i\bmod m)_{i\ge1}$ are
eventually periodic.
\begin{enumerate}[label=\textup{(\roman*)},leftmargin=2.2em]
\item \emph{Uniform Lipschitz shift invariance and sandwich.}
For every integer shift sequence $\mathbf c=(c_i)_{i\ge1}\in\Z_{\ge0}^{\N}$ and every
$\mathbf r\in\Z_{\ge1}^{\N}$, the shift map $\Psi_{\mathbf c}([0;a_1,a_2,\dots]):=[0;a_1+c_1,a_2+c_2,\dots]$
is uniformly $24$-Lipschitz on $F_m(\mathbf r)$ (in particular, with $K=24$
uniformly in $m\ge2$). Consequently, for every sequence $\mathbf r$ with
$1\le r_i\le m$,
\[
\dimH F_{m,m}\le \dimH F_m(\mathbf r)\le \dimH F_{m,1},
\]
and every eventually periodic modulus-$m$ residue set satisfies $\dimH E_m=\dimH F_{m,1}$.
\item \emph{Unconditional transfer-operator criterion and Lean soundness.}
Without appealing to Bowen's pressure formula or external \cite{CLU} citations,
a step-function upper inequality $\mathcal L_{s,\mathcal A_{m,r}}v\le(1-\eta)v$
or finite-$K$ subalphabet lower inequality $\mathcal L_{s,\mathcal D_K}v\ge(1+\eta)v$
yields $\dimH F_{m,r}\le s$ or $s\le\dimH F_{m,r}$ directly from
Hausdorff-measure covering and Frostman mass distribution. In particular,
$\frac12<\dimH F_{m,r}$ for all $m\ge1,r\ge1$, and $\dimH F_{4,1}\le\frac45$.
\item \emph{Exact integer certificates for $m\in\{2,3,4,5,6\}$.}
Running the verified step-function inequalities in exact integer arithmetic
(\texttt{verification/\allowbreak residue\_\allowbreak certificate\_\allowbreak general\_\allowbreak N512.py},
grid $N=512$, scale $S=10^{13}$, $2^{80}$-precision enclosures) certifies the
narrow rational brackets in Table~\ref{tab:residues-general}, with strict
separation $\mathrm{upper}(F_{m,m})<\mathrm{lower}(F_{m,1})$ (and hence
$\dimH F_{m,m}<\dimH F_{m,1}$) for every $m\in\{2,3,4,5,6\}$.
\end{enumerate}
\end{theorem}

\begin{proof}
For~(i), on any two distinct points $x,y\in F_m(\mathbf r)$ with first differing
partial quotient $a_{n+1}<b_{n+1}$, the congruence $a_{n+1}\equiv b_{n+1}\pmod m$
with $m\ge2$ forces $b_{n+1}-a_{n+1}\ge m\ge2$, giving a uniform tail expansion
ratio bound~$6$ and common-prefix Möbius distortion bound~$4$, hence global
$24$-Lipschitz constant across all~$m\ge2$; covering $E_m$ by countably many
Möbius images of shifted sets $F_m(\mathbf r)$ then yields $\dimH E_m=\dimH F_{m,1}$.
For~(ii), submultiplicativity of cylinder diameters yields the upper
Hausdorff-measure estimate directly from $\mathcal L_{s,\mathcal A}v\le v$, while
Frostman measure pushforward on compact subtrees gives the lower bound from
$\mathcal L_{s,\mathcal D_K}v\ge v$; both elementary tail bounds (monotone
integral upper bound vs.\ midpoint-shifted integral lower bound) reduce step-function
verification to finitely many grid cells. For~(iii), exact integer execution
checks every cell inequality for each $(m,r)$ in Table~\ref{tab:residues-general}.
All algebraic and measure-theoretic deductions in~(i)--(ii), together with the
step-function soundness theorems for~(ii)--(iii), are machine-checked in Lean~4;
the concrete grid arithmetic for $m\in\{2,3,4,5,6\}$ is performed both in exact
Python integer arithmetic and, for the closed brackets, by the Lean kernel
(Remark~\ref{rem:residue-kernel}).
\end{proof}

\begin{table}[ht]
\centering\small
\begin{tabular}{c|c|c|c|c|c}
$m$ & $r$ & $K_{\mathrm{tail}}$ & Lower bound & Upper bound & Width \\\hline
$2$ & $1$ & $258$ & $105/128\,(0.820312)$ & $421/512\,(0.822266)$ & $1/512$ \\
$2$ & $2$ & $258$ & $23/32\,(0.718750)$ & $369/512\,(0.720703)$ & $1/512$ \\\hline
$3$ & $1$ & $172$ & $95/128\,(0.742188)$ & $381/512\,(0.744141)$ & $1/512$ \\
$3$ & $3$ & $172$ & $41/64\,(0.640625)$ & $329/512\,(0.642578)$ & $1/512$ \\\hline
$4$ & $1$ & $130$ & $357/512\,(0.697266)$ & $179/256\,(0.699219)$ & $1/512$ \\
$4$ & $4$ & $130$ & $309/512\,(0.603516)$ & $155/256\,(0.605469)$ & $1/512$ \\\hline
$5$ & $1$ & $120$ & $171/256\,(0.667969)$ & $343/512\,(0.669922)$ & $1/512$ \\
$5$ & $5$ & $120$ & $149/256\,(0.582031)$ & $299/512\,(0.583984)$ & $1/512$ \\\hline
$6$ & $1$ & $120$ & $331/512\,(0.646484)$ & $83/128\,(0.648438)$ & $1/512$ \\
$6$ & $6$ & $120$ & $291/512\,(0.568359)$ & $73/128\,(0.570312)$ & $1/512$
\end{tabular}
\caption{Certified open intervals $(\mathrm{lower},\mathrm{upper})$ for $\dimH F_{m,r}$ ($m\in\{2,3,4,5,6\}$, $r\in\{1,m\}$) from \texttt{residue\_\allowbreak certificate\_\allowbreak general\_\allowbreak N512.py} ($N=512$, $S=10^{13}$, $2^{80}$-precision enclosures), all of width $1/512$; for every $m$, $\mathrm{upper}(F_{m,m})<\mathrm{lower}(F_{m,1})$. The closed brackets are kernel-checked in Lean~4 (Remark~\ref{rem:residue-kernel}).}\label{tab:residues-general}
\end{table}

\begin{remark}[Kernel-checked brackets]\label{rem:residue-kernel}
Independently of Bowen's formula, every closed bracket
$\mathrm{lower}\le\dimH F_{m,r}\le\mathrm{upper}$ in Table~\ref{tab:residues-general}
(twenty one-sided bounds) is a theorem of Lean~4 with Mathlib, proved from
Theorem~\ref{thm:residues-general}(ii) and the integer certificate data. The proof is
organised so that no floating-point computation is trusted: a generic soundness theorem
(\texttt{EpiFormal/\allowbreak ResidueCertGenSound.lean}) shows that, if Boolean checkers on
natural numbers (\texttt{EpiFormal/\allowbreak ResidueCertGenBase.lean}) return \texttt{true}
on a vector $v\in\Z_{\ge1}^{512}$ and a coefficient table, then $\dimH F_{m,r}\le s$
(respectively $\ge s$). The checkers verify that $v$ is non-increasing (so the maximum and
minimum of the step function over each image cell are its values at the cell endpoints,
whose indices are the closed forms $\min(N-1,\lfloor N^2/(aN+e)\rfloor)$), that each
coefficient $c$ satisfies the exponent-cleared enclosure
$S^qN^p\le c^q(aN+j)^p$ (upper) or $c^q(aN+j+1)^p\le S^qN^p$ (lower) for $2s=p/q$, the
row inequalities, and the integral tail bounds. These checks are evaluated by the Lean
kernel (\texttt{decide +kernel}, GMP-accelerated natural-number arithmetic; no
\texttt{native\_decide}) on data emitted by \texttt{verification/\allowbreak
emit\_cert\_gen.py} in $64$-row chunks; the resulting bounds and the strict separations
$\dimH F_{m,m}<\dimH F_{m,1}$ ($2\le m\le6$) are collected in
\texttt{EpiFormal/\allowbreak ResidueCertTable.lean}. Strictness of the brackets
(open intervals) uses the positive margins and the pressure argument of
Theorem~\ref{thm:residues}, which is not formalised.
\end{remark}


\subsection{Resolvent identity, global strict separation for all \texorpdfstring{$m\ge2$}{m >= 2}, and reciprocal-Fibonacci asymptotics}\label{sec:residues-asymptotic}

Table~\ref{tab:residues-general} certifies $\dimH F_{m,m}<\dimH F_{m,1}=\dimH E_m$
for $m\in\{2,3,4,5,6\}$ via $512$-cell transfer-operator grid certificates. To
prove this strict separation unconditionally for \emph{every} modulus $m\ge2$ and
to determine the exact asymptotic law of $\dimH F_{m,r}$ as $m\to\infty$, we
introduce a family of closed-form rational test functions derived from the
metallic-mean continuants of~$r$.

For any integer $r\ge1$, define the metallic-mean continuant sequence
$(q_n(r))_{n\ge0}$ (\texttt{metalQ r n}, with predecessor $q_{-1}(r):=0$,
\texttt{metalQPrev r n}) by
\begin{equation}\label{eq:metallic-continuant}
q_0(r)=1,\qquad q_1(r)=r,\qquad q_{n+2}(r)=r\,q_{n+1}(r)+q_n(r)\quad(n\ge0).
\end{equation}
In particular, for $r=1$, $q_n(1)=F_{n+1}$ is the $(n+1)$-st Fibonacci number
($F_1=F_2=1,F_3=2,F_4=3,F_5=5,\dots$, \texttt{metalQ\_\allowbreak one\_\allowbreak eq\_\allowbreak fib}).
For $s\ge1/2$, define the metallic-continuant series (\texttt{metalC r s}) and
the associated resolvent test profile on $[0,1]$ (\texttt{resolventV r s x}) by
\begin{equation}\label{eq:Cr-Dr}
C_r(s):=\sum_{n=0}^\infty q_n(r)^{-2s},\quad
C_r:=C_r(1/2)=\sum_{n=0}^\infty\frac{1}{q_n(r)},\quad
C_1=\psi_{\mathrm{Fib}}:=\sum_{k=1}^\infty\frac{1}{F_k}\approx 3.3598856662,
\end{equation}
\[
V_{r,s}(x):=\sum_{n=0}^\infty\bigl(q_n(r)+q_{n-1}(r)\,x\bigr)^{-2s}\qquad(x\in[0,1]),
\]
so that $V_{r,s}(0)=C_r(s)$ and the $n=0$ term of $V_{r,s}(x)$ is $(1+0\cdot x)^{-2s}=1$.

\begin{theorem}[Resolvent identity and reciprocal-Fibonacci asymptotics]\label{thm:residues-asymptotic}
\begin{enumerate}[label=(\roman*),leftmargin=2.2em]
\item \textbf{Exact metallic-continuant resolvent identity:}
For every integer $r\ge1$, every $s\ge1/2$, and every $x\in[0,1]$, the first
branch $a=r$ ($k=0$) of the transfer operator $\mathcal L_{s,F_{m,r}}$ acts on
$V_{r,s}$ by the exact algebraic identity
\begin{equation}\label{eq:resolvent-branch}
(r+x)^{-2s}\,V_{r,s}\!\left(\frac{1}{r+x}\right)
=V_{r,s}(x)-1.
\end{equation}
\item \textbf{Global strict separation $\dimH F_{m,m}<\dimH F_{m,1}=\dimH E_m$ for all $m\ge2$:}
For every modulus $m\ge2$,
\begin{equation}\label{eq:global-strict-sep}
\dimH F_{m,m}
<\dimH F_{m,1}
=\dimH E_m.
\end{equation}
\item \textbf{Two-term asymptotic expansions ($m\to\infty$):}
For every fixed integer $r\ge1$, there exists an explicit constant $B_r>0$ such
that for all sufficiently large $m$,
\begin{equation}\label{eq:dimH-Fmr-two-term}
\left|\dimH F_{m,r}-\left(\frac12+\frac{C_r}{2m}-\frac{C_r^2\log m}{2m^2}\right)\right|
\le\frac{B_r}{m^2},
\end{equation}
where $C_r=\sum_{n=0}^\infty 1/q_n(r)$ and $C_1=\psi_{\mathrm{Fib}}=\sum_{k=1}^\infty 1/F_k$.
In contrast, for the $0\bmod m$ residue alphabet $F_{m,m}$ there exists $B_0>0$
such that for all sufficiently large $m$,
\begin{equation}\label{eq:dimH-Fmm-two-term}
\left|\dimH F_{m,m}-\left(\frac12+\frac{1}{2m}-\frac{\log m}{2m^2}\right)\right|
\le\frac{B_0}{m^2}.
\end{equation}
\item \textbf{Asymptotic dimension-excess ratios and eventual strict monotonicity in~$r$:}
For every fixed $r\ge1$, $\lim_{m\to\infty}m(\dimH F_{m,r}-1/2)=C_r/2$, and
\[
\lim_{m\to\infty}\frac{\dimH F_{m,r}-1/2}{\dimH F_{m,m}-1/2}=C_r,\qquad
\lim_{m\to\infty}\frac{\dimH E_m-1/2}{\dimH F_{m,m}-1/2}=\psi_{\mathrm{Fib}},
\]
and since $r\mapsto C_r$ is strictly decreasing on $\Z_{\ge1}$, whenever
$1\le r_1<r_2$ we have $\dimH F_{m,r_2}<\dimH F_{m,r_1}$ for all sufficiently
large~$m$.
\end{enumerate}
\end{theorem}

\begin{proof}
For~(i), since $q_{n+2}(r)\ge 2\,q_n(r)$, $q_n(r)\ge 2^{\lfloor n/2\rfloor}$ grows
exponentially and $V_{r,s}(x)$ converges absolutely for all $s\ge1/2$ and $x\in[0,1]$.
For each $n\ge0$, the recurrence $q_{n+1}(r)=r\,q_n(r)+q_{n-1}(r)$ gives
\[
(r+x)\cdot\left(q_n(r)+q_{n-1}(r)\,\frac{1}{r+x}\right)
=r\,q_n(r)+q_{n-1}(r)+q_n(r)x
=q_{n+1}(r)+q_n(r)x.
\]
Raising both sides to the power $-2s$ and summing over $n\ge0$ shifts the index
$n\mapsto n+1$ and separates the $n=0$ term $(q_0(r)+q_{-1}(r)x)^{-2s}=1^{-2s}=1$,
yielding~\eqref{eq:resolvent-branch} (\texttt{resolventV\_\allowbreak branch\_\allowbreak eq}).

For~(ii), by \eqref{eq:resolvent-branch} the first branch $k=0$ ($a=r$) of the
transfer operator $\mathcal L_{s,F_{m,r}}$ cancels $V_{r,s}(x)-1$ identically:
\[
(\mathcal L_{s,F_{m,r}}V_{r,s})(x)-V_{r,s}(x)
=\sum_{k=1}^\infty(km+r+x)^{-2s}\,V_{r,s}\!\left(\frac{1}{km+r+x}\right)-1.
\]
Since $V_{r,s}$ is decreasing on $[0,1]$ with $C_r(s)\ge V_{r,s}(1/(km+r+x))\ge V_{r,s}(1/(m+r))$,
combining Theorem~\ref{thm:residues-general}(ii) with standard integral comparisons for
$\sum_{k=1}^\infty(km+r+x)^{-2s}$ gives the non-asymptotic resolvent criteria
(\texttt{dimH\_\allowbreak residueConstSet\_\allowbreak le\_\allowbreak of\_\allowbreak resolvent}
and \texttt{le\_\allowbreak dimH\_\allowbreak residueConstSet\_\allowbreak of\_\allowbreak resolvent}):
for $s>1/2$, $C_r(s)\,m^{-2s}\bigl(1+\frac{1}{2s-1}\bigr)\le 1$ implies $\dimH F_{m,r}\le s$,
and $V_{r,s}\bigl(\frac{1}{m+r}\bigr)\frac{(m+r+1)^{1-2s}}{m(2s-1)}>1$ implies $\dimH F_{m,r}\ge s$,
while applying the same comparison to the constant test function $v\equiv 1$ on
$F_{m,m}$ ($a=km$, $k\ge1$) yields \texttt{dimH\_\allowbreak residueConstSet\_\allowbreak mm\_\allowbreak le}
and \texttt{le\_\allowbreak dimH\_\allowbreak residueConstSet\_\allowbreak mm}.
Evaluating these two criteria at explicit rational test exponents separates
$\dimH F_{m,m}$ and $\dimH F_{m,1}$ for every $m\ge2$
(\texttt{dimH\_\allowbreak residueConstSet\_\allowbreak mm\_\allowbreak lt\_\allowbreak m1}),
and Theorem~\ref{thm:residues-general}(i) identifies $\dimH E_m=\dimH F_{m,1}$.

For~(iii) and~(iv), for fixed $r\ge1$ and large $m$, write $s=\frac12+\frac{\alpha}{2m}$
with $\alpha=C_r-\frac{C_r^2\log m}{m}\pm\frac{2B_r}{m}$. Because $2s-1=\alpha/m\to0$,
we have $C_r(s)=C_r+O(2s-1)=C_r+O(m^{-1})$,
$V_{r,s}(1/(m+r))=C_r(s)+O(m^{-1})=C_r+O(m^{-1})$, and
$m^{-(2s-1)}=\exp(-\frac{\alpha\log m}{m})=1-\frac{C_r\log m}{m}+O\bigl((\frac{\log m}{m})^2\bigr)$.
Substituting $\alpha_u=C_r-\frac{C_r^2\log m}{m}+\frac{2B_r}{m}$ into the upper
resolvent criterion and $\alpha_l=C_r-\frac{C_r^2\log m}{m}-\frac{2B_r}{m}$ into the
lower resolvent criterion establishes~\eqref{eq:dimH-Fmr-two-term}
(\texttt{dimH\_\allowbreak residueConstSet\_\allowbreak asymp\_\allowbreak two\_\allowbreak term}),
while the constant-$1$ criteria establish~\eqref{eq:dimH-Fmm-two-term}
(\texttt{dimH\_\allowbreak residueConstSet\_\allowbreak mm\_\allowbreak asymp\_\allowbreak two\_\allowbreak term}).
Dividing $m(\dimH F_{m,r}-1/2)\to C_r/2$ by $m(\dimH F_{m,m}-1/2)\to 1/2$ gives the
exact dimension-excess ratio $C_r$ (and $\psi_{\mathrm{Fib}}$ for $E_m$ via
$\dimH E_m=\dimH F_{m,1}$), and strict monotonicity $q_n(r_1)<q_n(r_2)$ for $n\ge1$
implies $C_{r_2}<C_{r_1}$ and hence $\dimH F_{m,r_2}<\dimH F_{m,r_1}$ for all large~$m$.
\end{proof}

\begin{table}[ht]
\centering\small\setlength{\tabcolsep}{4.5pt}
\begin{tabular}{c|c|c|c|c}
$r$ & Continuants $(q_n(r))_{n\ge0}$ & $C_r=\sum_{n=0}^\infty 1/q_n(r)$ & $s_{\mathrm{2term}}(1000,r)$ & Ratio ($m=5000$) \\\hline
$1$ & $1,1,2,3,5,8,13,\dots$ ($F_{n+1}$) & $\psi_{\mathrm{Fib}}\approx 3.359885666243$ & $0.5016409526$ & $3.346356$ \\
$2$ & $1,2,5,12,29,70,169,\dots$ (Pell) & $1.842203049828$ & $0.5009093801$ & $1.839556$ \\
$3$ & $1,3,10,33,109,360,\dots$ & $1.476794726319$ & $0.5007308647$ & $1.475593$ \\
$4$ & $1,4,17,72,305,1292,\dots$ & $1.327004277866$ & $0.5006574201$ & $1.326264$ \\
$5$ & $1,5,26,135,701,3640,\dots$ & $1.247635731173$ & $0.5006184416$ & $1.247109$ \\\hline
$m$ & $F_{m,m}$ ($0\bmod m$ alphabet) & $C_\infty=1.000000000000$ & $0.5004965461$ & $1.000000$
\end{tabular}
\caption{High-precision ($50$-digit \texttt{mpmath}, \texttt{verify\_\allowbreak paper2\_\allowbreak extensions.py}) values of the reciprocal metallic-continuant constants $C_r=\sum_{k=1}^\infty 1/q_k(r)$ ($C_1=\psi_{\mathrm{Fib}}$), the two-term asymptotic approximation $\frac12+\frac{C_r}{2m}-\frac{C_r^2\log m}{2m^2}$ at $m=1000$, and the two-term dimension-excess ratio $(\dimH F_{m,r}-1/2)/(\dimH F_{m,m}-1/2)$ at $m=5000$ converging to $C_r$.}\label{tab:residues-asymptotic}
\end{table}

\subsection{Multi-residue phase transition, Euler--Mascheroni three-term asymptotics, and supercritical power-law rate}\label{sec:multi-residue}

Theorem~\ref{thm:residues-asymptotic} treats single congruence classes $r\pmod m$.
More generally, fix two finite digit sets $S\subset\Z_{\ge1}$ and
$B\subset\Z_{\ge0}$ with $d:=|S|+|B|\ge1$. For every integer modulus $m\ge1$,
define the \emph{multi-residue digit alphabet}
(\texttt{multiResidueAlphabet m S B}) and its continued-fraction Cantor set
(\texttt{multiResidueSet m S B}) by
\[
\mathcal A_{m,S,B}:=\bigl\{a\in\Z_{\ge1} : (\exists r\in S,\ a\ge r\wedge a\equiv r\!\!\pmod m)\ \text{or}\ (\exists b\in B,\ \exists k\ge1,\ a+b=km)\bigr\},
\]
\[
F_{m,S,B}:=\bigl\{[0;a_1,a_2,\dots]\in(0,1)\setminus\Q : a_n\in\mathcal A_{m,S,B}\ \forall n\ge1\bigr\}.
\]
In particular, $F_{m,\{r\},\emptyset}=F_{m,r}$ for $r\ge1$ and
$F_{m,\emptyset,\{0\}}=F_{m,m}$.
For finite $S\subset\Z_{\ge1}$, denote by $W_S(p,n):=\sum_{w\in S^n}q(w)^{-p}$
the $n$-th continuant word sum (\texttt{wordSum S p n}, with $W_S(p,0)=1$, where
$q(w)$ is the continuant denominator of the finite word $w$), by
\[
s_0(S):=\inf\Bigl\{s>0:\sum_{n=0}^\infty W_S(2s,n)<\infty\Bigr\}
\]
the Poincar\'e abscissa (\texttt{poincareAbscissa S}), and, whenever convergent,
by $C_S(p):=\sum_{n=0}^\infty W_S(p,n)$ (\texttt{multiC S p}, with $C_S:=C_S(1)$),
$C_S':=\sum_{n=0}^\infty\sum_{w\in S^n}\frac{\log q(w)}{q(w)}$ (\texttt{multiCLog S}),
and $V_{S,s}(x):=\sum_{n=0}^\infty\sum_{w\in S^n}(q(w)+p(w)x)^{-2s}$
(\texttt{multiResolventV S s x}, where $p(w)/q(w)=[0;w]$).

\begin{theorem}[Multi-residue phase transition and three-term asymptotics]\label{thm:multi-residue-asymptotic}
\begin{enumerate}[label=\textup{(\roman*)},leftmargin=2.2em]
\item \textbf{Bowen-free Poincar\'e-abscissa dimension for finite alphabets:}
For every nonempty finite alphabet $S\subset\Z_{\ge1}$, writing
$E_S:=\{[0;a_1,a_2,\dots]:a_n\in S\}$, one has $\dimH E_S=s_0(S)$, established
directly from continuant multiplicativity and the multi-branch resolvent
identity
\begin{equation}\label{eq:multi-resolvent-branch}
\sum_{a\in S}(a+x)^{-2s}\,V_{S,s}\!\left(\frac{1}{a+x}\right)=V_{S,s}(x)-1
\qquad\bigl(s>s_0(S),\ x\in[0,1]\bigr)
\end{equation}
via the unconditional transfer-operator criterion of Theorem~\ref{thm:residues-general}\textup{(ii)},
without invoking Bowen's pressure formula.
\item \textbf{Universal multi-residue phase transition at $1/2$:}
For all finite sets $S\subset\Z_{\ge1}$ and $B\subset\Z_{\ge0}$ with
$d=|S|+|B|\ge1$,
\[
\lim_{m\to\infty}\dimH F_{m,S,B}=\max\!\Bigl(\frac12,\,s_0(S)\Bigr).
\]
Both regimes occur already for two-digit alphabets: $S=\{2,3\}$ is subcritical
with $\sum_{n=0}^\infty W_{\{2,3\}}(1,n)<\infty$ and $s_0(\{2,3\})<1/2$, whereas
$S=\{1,2\}$ is supercritical with $s_0(\{1,2\})\ge 51/100>1/2$.
\item \textbf{Elementary Euler--Mascheroni expansion of $\zeta(1+\varepsilon)$:}
With $C_\zeta:=4\zeta(3/2)>0$ and $\gamma=\lim_{N\to\infty}(H_N-\log N)$ the
Euler--Mascheroni constant, every $\varepsilon\in(0,1]$ satisfies
\[
\left|\sum_{k=1}^\infty k^{-(1+\varepsilon)}-\frac{1}{\varepsilon}-\gamma\right|
\le C_\zeta\,\varepsilon.
\]
\item \textbf{Master subcritical three-term asymptotics to $O((\log m)^2/m^3)$:}
Whenever $d=|S|+|B|\ge1$ and $\sum_{n=0}^\infty W_S(1,n)<\infty$
(so $S=\emptyset$ or $s_0(S)<1/2$), there exists a constant $K>0$ such that for
all sufficiently large~$m$,
\begin{equation}\label{eq:multi-dimH-three-term}
\left|\dimH F_{m,S,B}-\left(\frac12+\frac{d\,C_S}{2m}-\frac{(d\,C_S)^2\log m}{2m^2}+\frac{d^2\,C_S(\gamma C_S-C_S')}{2m^2}\right)\right|
\le K\,\frac{(\log m)^2}{m^3}.
\end{equation}
Consequently,
\[
\lim_{m\to\infty}m^2\!\left(\dimH F_{m,S,B}-\frac12-\frac{d\,C_S}{2m}+\frac{(d\,C_S)^2\log m}{2m^2}\right)
=\frac{d^2\,C_S(\gamma C_S-C_S')}{2},
\]
\[
\lim_{m\to\infty}\frac{\dimH F_{m,S,B}-1/2}{\dimH F_{m,m}-1/2}=d\,C_S.
\]
In particular:
\begin{itemize}[leftmargin=1.5em]
\item for every fixed $r\ge1$,
$\dimH F_{m,r}=\frac12+\frac{C_r}{2m}-\frac{C_r^2\log m}{2m^2}+\frac{C_r(\gamma C_r-C_r')}{2m^2}+O\!\left(\frac{(\log m)^2}{m^3}\right)$,
where $C_r':=\sum_{n=0}^\infty\frac{\log q_n(r)}{q_n(r)}$;
\item for $F_{m,m}$ ($S=\emptyset, B=\{0\}$),
$\dimH F_{m,m}=\frac12+\frac{1}{2m}-\frac{\log m}{2m^2}+\frac{\gamma}{2m^2}+O\!\left(\frac{(\log m)^2}{m^3}\right)$;
\item for $S=\{2,3\}, B=\emptyset$, \eqref{eq:multi-dimH-three-term} holds with $d=2$.
\end{itemize}
\item \textbf{Supercritical two-sided power-law convergence rate $\Theta(m^{-2s_0})$:}
Whenever $s_0(S)>1/2$, there exist constants $0<K_1\le K_2$ (depending on
$S,B$) such that for all sufficiently large~$m$,
\begin{equation}\label{eq:multi-supercritical-rate}
K_1\,m^{-2s_0(S)}\le \dimH F_{m,S,B}-s_0(S)\le K_2\,m^{-2s_0(S)}.
\end{equation}
In particular, \eqref{eq:multi-supercritical-rate} holds for $S=\{1,2\}$ and
every finite $B\subset\Z_{\ge0}$.
\end{enumerate}
\end{theorem}

\begin{proof}
For~(i), continuant concatenation $q(uv)\in[q(u)q(v),\,2q(u)q(v)]$ yields
the sub- and super-multiplicative word-sum bounds
$2^{-p}\,W_S(p,k)W_S(p,n)\le W_S(p,k+n)\le W_S(p,k)W_S(p,n)$
(\texttt{wordSum\_\allowbreak add\_\allowbreak le},
\texttt{wordSum\_\allowbreak add\_\allowbreak ge}), which force
$1\le W_S(2s_0(S),n)\le 2^{2s_0(S)}$ (\texttt{one\_\allowbreak le\_\allowbreak wordSum\_\allowbreak poincareAbscissa},
\texttt{wordSum\_\allowbreak poincareAbscissa\_\allowbreak le}), absolute
convergence of $W_S(2s,n)$ for $s>s_0(S)$
(\texttt{summable\_\allowbreak wordSum\_\allowbreak of\_\allowbreak poincareAbscissa\_\allowbreak lt}),
and divergence for $0<s<s_0(S)$
(\texttt{not\_\allowbreak summable\_\allowbreak wordSum\_\allowbreak of\_\allowbreak lt\_\allowbreak poincareAbscissa}).
Shifting words $w\in S^n$ to $aw$ gives \eqref{eq:multi-resolvent-branch}
(\texttt{multiResolventV\_\allowbreak branch\_\allowbreak eq}); hence
$\mathcal L_{s,S}V_{S,s}=V_{S,s}-1\le V_{S,s}$ for $s>s_0(S)$ gives
$\dimH E_S\le s_0(S)$, while truncated partial sums
$\sum_{j=0}^{N-1}\mathcal L_{s,S}^j 1$ for $0<s<s_0(S)$ exceed their own image
for large~$N$, giving $\dimH E_S\ge s_0(S)$ via
Theorem~\ref{thm:residues-general}(ii) without Bowen's formula
(\texttt{dimH\_\allowbreak cfAlphabetSet\_\allowbreak finset\_\allowbreak eq\_\allowbreak poincareAbscissa}).

For~(ii), once $m>\max S+\max B$, every tail digit
$a\in\mathcal A_{m,S,B}\setminus S$ is at least $m-\max B\asymp m$, so
$\mathcal L_{s,\mathcal A_{m,S,B}}V_{S,s}(x)-V_{S,s}(x)
=\sum_{a\in\mathcal A_{m,S,B}\setminus S}(a+x)^{-2s}V_{S,s}(1/(a+x))-1
\le O(m^{-2s}/(2s-1))-1$ forces $\limsup_{m\to\infty}\dimH F_{m,S,B}\le \max(1/2,s_0(S))$;
conversely, $F_{m,S,B}\supseteq E_S$ gives $\dimH F_{m,S,B}\ge s_0(S)$ and the
tail branch sum gives $\dimH F_{m,S,B}>1/2$, proving
$\lim_{m\to\infty}\dimH F_{m,S,B}=\max(1/2,s_0(S))$
(\texttt{tendsto\_\allowbreak dimH\_\allowbreak multiResidueSet}).
For $S=\{2,3\}$, $W_{\{2,3\}}(1,1)=\frac12+\frac13=\frac56<1$ gives geometric
summability at $s=1/2$ and, by continuity, below $1/2$
(\texttt{subcritical\_\allowbreak pair\_\allowbreak two\_\allowbreak three}); for $S=\{1,2\}$, evaluating
$\mathcal L_{51/100,\{1,2\}}v$ on $v(x)=2/(2+x)$ gives
$\frac{2}{2x+3}(1+x)^{-1/50}+\frac{2}{2x+5}(2+x)^{-1/50}\ge\frac{35}{36}\bigl(\frac{2}{2+x}+\frac{2}{(2x+3)(2x+5)(2+x)}\bigr)\ge v(x)$
on $[0,1]$ since $(2x+3)(2x+5)\le35$ and $(35/36)^{50}\le1/3$, whence
$s_0(\{1,2\})\ge 51/100>1/2$
(\texttt{supercritical\_\allowbreak pair\_\allowbreak one\_\allowbreak two}).

For~(iii), writing $H_N-\log(N+1)=\sum_{k=0}^{N-1}(a_k^{\mathrm{hom}}-b_k)$
and $\sum_{k=1}^N k^{-(1+\varepsilon)}-(1-(N+1)^{-\varepsilon})/\varepsilon$
via telescoping differences of $(k+1)^{-\varepsilon}/\varepsilon$ and
$\log(k+1)$, applying the mean value theorem twice yields $|a_k-b_k|\le 4\varepsilon(k+1)^{-3/2}$,
and passing $N\to\infty$ proves~\textup{(iii)}
(\texttt{abs\_\allowbreak tsum\_\allowbreak rpow\_\allowbreak sub\_\allowbreak inv\_\allowbreak sub\_\allowbreak eulerMascheroni\_\allowbreak le}).

For~(iv), Taylor-bounding word powers on $[1,1+\varepsilon]$ gives
$C_S(1+\varepsilon)=C_S-C_S'\varepsilon+O(\varepsilon^2)$
(\texttt{multiC\_\allowbreak one\_\allowbreak add\_\allowbreak bounds}), and
multiplying $V_{S,s}(1/(km+O(1)))=C_S(2s)+O((km)^{-1})$ against the Hurwitz
tail expanded via~\textup{(iii)} turns both resolvent inequalities into
\[
d\,C_S(1+\varepsilon)\left(\frac{1}{\varepsilon}+\gamma+O(\varepsilon)+O(m^{-1})\right)m^{-\varepsilon}=1+O(m^{-1}),\qquad \varepsilon:=2s-1=\frac{\alpha}{m}.
\]
Expanding $m^{-\varepsilon}=1-\varepsilon\log m+\frac12(\varepsilon\log m)^2+O((\log m/m)^3)$
yields $\alpha=dC_S-\frac{(dC_S)^2\log m}{m}+\frac{d^2C_S(\gamma C_S-C_S')}{m}+O\!\left(\frac{(\log m)^2}{m^2}\right)$,
which gives~\eqref{eq:multi-dimH-three-term} and its limits and specializations
(\texttt{dimH\_\allowbreak multiResidueSet\_\allowbreak asymp\_\allowbreak three\_\allowbreak term},
\texttt{tendsto\_\allowbreak m2\_\allowbreak residual\_\allowbreak multiResidue},
\texttt{tendsto\_\allowbreak dimH\_\allowbreak multiResidue\_\allowbreak ratio},
\texttt{dimH\_\allowbreak residueConstSet\_\allowbreak asymp\_\allowbreak three\_\allowbreak term},
\texttt{dimH\_\allowbreak residueConstSet\_\allowbreak mm\_\allowbreak asymp\_\allowbreak three\_\allowbreak term},
\texttt{dimH\_\allowbreak multiResidueSet\_\allowbreak two\_\allowbreak three\_\allowbreak asymp}).

For~(v), submultiplicativity together with $W_S(2s_0,n)\le 2^{2s_0}\le 4$ gives
the quantitative pole bound $C_S(2(s_0+\delta))\le 32/\delta$ for
$0<\delta\le1$ (\texttt{multiC\_\allowbreak supercritical\_\allowbreak le}).
Taking $\delta=K_2 m^{-2s_0}$ in the upper resolvent criterion gives
$C_S(2s)\cdot O(m^{-2s_0})\le 1$, proving
$\dimH F_{m,S,B}-s_0(S)\le K_2 m^{-2s_0(S)}$
(\texttt{dimH\_\allowbreak multiResidueSet\_\allowbreak supercritical\_\allowbreak upper}).
Conversely, pick any $r_0\in S$, form the finite subalphabet $A_1:=S\cup\{m+r_0\}\subset\mathcal A_{m,S,B}$
with $m+r_0\notin S$, set $\delta:=K_1 m^{-2s_0}$ ($K_1:=1/(4096\log(\max S+1))$)
and $s:=s_0+\delta$, and define the finite-$n$ partial-sum test function
$v(y):=\sum_{j=0}^{n-1}(\mathcal L_{s,S}^j 1)(y)$ with $n:=\lfloor 1/(2\delta\log(\max S+1))\rfloor$.
On $[0,1]$, $1/48\le (\mathcal L_{s,S}^j 1)(y)\le 4$ holds for all $j<n$, while
the extra branch $a_1=m+r_0$ contributes
$(m+r_0+x)^{-2s}v(1/(m+r_0+x))\ge (2m)^{-2s}\cdot\frac{n}{48}\ge 1$, making
$\mathcal L_{s,A_1}v=v-1+\mathcal L_{s,S}^n 1+(m+r_0+\cdot)^{-2s}v\circ f_{m+r_0}\ge v$
and forcing $\dimH F_{m,S,B}\ge\dimH E_{A_1}\ge s_0(S)+K_1 m^{-2s_0(S)}$
(\texttt{dimH\_\allowbreak multiResidueSet\_\allowbreak supercritical\_\allowbreak lower},
\texttt{dimH\_\allowbreak multiResidueSet\_\allowbreak supercritical\_\allowbreak rate},
\texttt{dimH\_\allowbreak multiResidueSet\_\allowbreak one\_\allowbreak two\_\allowbreak supercritical\_\allowbreak rate}).
\end{proof}

\begin{table}[ht]
\centering\small\setlength{\tabcolsep}{4.5pt}
\begin{tabular}{c|c|c|c|c|c}
$(S,B)$ & $d$ & $C_S$ & $C_S'$ & $c_2=\frac12 d^2 C_S(\gamma C_S-C_S')$ & $m^2\cdot\mathrm{resid}$ ($16384$) \\\hline
$S=\{1\},\,B=\emptyset$ ($r=1$) & $1$ & $3.359885666$ & $1.968398116$ & $-0.048751$ & $+0.113703$ \\
$S=\{2\},\,B=\emptyset$ ($r=2$) & $1$ & $1.842203050$ & $1.110456868$ & $-0.043392$ & $-0.016699$ \\
$S=\{3\},\,B=\emptyset$ ($r=3$) & $1$ & $1.476794726$ & $0.770974876$ & $+0.060146$ & $+0.073533$ \\
$S=\emptyset,\,B=\{0\}$ ($r=m$) & $1$ & $1.000000000$ & $0.000000000$ & $+0.288608$ ($\gamma/2$) & $+0.292414$ \\\hline
$S=\{2,3\},\,B=\emptyset$ & $2$ & $3.921382875$ & $10.062207836$ & $-61.163567$ & $-58.118127$ \\
$S=\{1\},\,B=\{0\}$ & $2$ & $3.359885666$ & $1.968398116$ & $-0.195004$ & $+1.105827$ \\
$S=\emptyset,\,B=\{0,1\}$ & $2$ & $1.000000000$ & $0.000000000$ & $+1.154431$ ($2\gamma$) & $+1.184953$ \\
$S=\{2,4,5\},\,B=\emptyset$ & $3$ & $7.292391500$ & $49.948931200$ & $-1500.981411$ & $-1375.047467$
\end{tabular}
\caption{High-precision ($32$-digit \texttt{mpmath} Chebyshev collocation, \texttt{verify\_\allowbreak residue\_\allowbreak multi\_\allowbreak asymptotic.py}, non-rigorous numerical corroboration) values of $C_S$, $C_S'$, the theoretical third-order coefficient $c_2=\frac12 d^2 C_S(\gamma C_S-C_S')$, and the observed second-order residual $m^2\bigl(\dimH F_{m,S,B}-\frac12-\frac{dC_S}{2m}+\frac{(dC_S)^2\log m}{2m^2}\bigr)$ at $m=16384$. In the supercritical regime, the script also checks $S=\{1,2\}$ ($s_0\approx0.531280506$, $m^{2s_0}(\dimH F_{m,\{1,2\},\emptyset}-s_0)\approx 28.78$ at $m=4096$) and $S=\{1,2,3\}$ ($s_0\approx0.705660908$, $m^{2s_0}(\dimH-s_0)\approx 8.81$ at $m=4096$).}\label{tab:multi-residue-asymptotic}
\end{table}

\begin{remark}[A heuristic fourth-order term; not proved]\label{rem:fourth-term}
Put $\Gamma:=d\,C_S$, $\Gamma':=d\,C_S'$ and $\varepsilon:=2\dimH F_{m,S,B}-1$. Replacing the
leading eigenvalue of the transfer operator by its scalar part and $\zeta(1+\varepsilon)$ by
$\varepsilon^{-1}+\gamma$ (part~(iii)) leads to the scalar model
\[
\varepsilon\,m=(\Gamma-\varepsilon\Gamma')(1+\gamma\varepsilon)\,m^{-\varepsilon},
\]
whose small root we denote by $\varepsilon_m^{\mathrm{mod}}$. Expanding it (symbolically, with
\texttt{sympy}) gives
\[
\frac{1+\varepsilon_m^{\mathrm{mod}}}{2}=\frac12+\frac{\Gamma}{2m}-\frac{\Gamma^2\log m}{2m^2}
+\frac{\Gamma(\gamma\Gamma-\Gamma')}{2m^2}+\frac{3\Gamma^3(\log m)^2}{4m^3}
+O\Bigl(\frac{\log m}{m^3}\Bigr),
\]
whose first four terms are exactly those proved in~\eqref{eq:multi-dimH-three-term}. This
suggests, \emph{as a conjecture}, that in the subcritical regime
\[
\dimH F_{m,S,B}=\frac12+\frac{\Gamma}{2m}-\frac{\Gamma^2\log m}{2m^2}
+\frac{\Gamma(\gamma\Gamma-\Gamma')}{2m^2}+\frac{3\Gamma^3(\log m)^2}{4m^3}
+O\Bigl(\frac{\log m}{m^3}\Bigr).
\]
Table~\ref{tab:fourth-term} compares the model directly with collocation values of the
dimension; this comparison does not depend on fitting the $\log m/m^3$ and $1/m^3$ terms, which a
free least-squares fit cannot separate from the $(\log m)^2/m^3$ term for $\log m\le 10$. If the
conjecture holds, the tabulated quantity tends to $0$, whereas if the coefficient
$\frac34\Gamma^3$ were wrong it would tend to a nonzero constant. In 7 of the ten
cases its absolute value decreases over $m=1024,4096,16384$; in 2 more it
decreases from $m=4096$ to $m=16384$; for $S=\{1,3\}$ ($\Gamma\approx40.9$) it still increases,
i.e.\ $m\le16384$ is not yet in the asymptotic range for this case. The values for $m=1024$
and $r\in\{1,2,3,m\}$ come from \texttt{verification/\allowbreak asymptotic\_\allowbreak residues\_\allowbreak check\_\allowbreak 20261123.py}
(collocation sizes $16$ and $24$, which differ there by less than $10^{-15}$); the others from
\texttt{verify\_\allowbreak residue\_\allowbreak multi\_\allowbreak asymptotic.py}. For $m\in\{4096,16384\}$ and
$r\in\{1,2,3,m\}$ the first script reproduces the second to within $10^{-15}$
(changing the tabulated quantity by less than $10^{-5}$), with
collocation sizes differing by less than $10^{-15}$. All of this is non-rigorous numerical
evidence (\texttt{verification/\allowbreak fourth\_\allowbreak term\_\allowbreak fit\_\allowbreak 20261124.py}); neither the
fourth term nor the size of the remainder is proved.
\end{remark}

\begin{table}[ht]
\centering\small\setlength{\tabcolsep}{4pt}
\begin{tabular}{c|c|c|ccc|c}
$S$ & $B$ & $\Gamma=dC_S$ & $m=1024$ & $m=4096$ & $m=16384$ & $\frac34\Gamma^3$ \\\hline
$\{1\}$ & $\emptyset$ & $3.36$ & $-0.124$ & $-0.088$ & $-0.065$ & $+28.4$ \\
$\{2\}$ & $\emptyset$ & $1.842$ & $-0.093$ & $-0.066$ & $-0.049$ & $+4.689$ \\
$\{3\}$ & $\emptyset$ & $1.477$ & $-0.095$ & $-0.068$ & $-0.050$ & $+2.416$ \\
$\emptyset$ & $\{0\}$ & $1$ & $+0.001$ & $+0.001$ & $+0.000$ & $+0.750$ \\
$\emptyset$ & $\{0,1\}$ & $2$ & $+0.039$ & $+0.028$ & $+0.020$ & $+6.000$ \\
$\{1\}$ & $\{0\}$ & $6.72$ & $+0.766$ & $+0.626$ & $+0.484$ & $+228$ \\
$\{2,3\}$ & $\emptyset$ & $7.843$ & $+28.4$ & $+23.7$ & $+18.4$ & $+362$ \\
$\{1,4\}$ & $\emptyset$ & $20.33$ & $+1163$ & $+1332$ & $+1160$ & $+6301$ \\
$\{2,4,5\}$ & $\emptyset$ & $21.88$ & $+2304$ & $+2897$ & $+2616$ & $+7853$ \\
$\{1,3\}$ & $\emptyset$ & $40.91$ & $+15538$ & $+30653$ & $+34429$ & $+51350$ \\
\end{tabular}
\caption{The quantity $m^3\bigl(\dimH F_{m,S,B}-\frac{1+\varepsilon_m^{\mathrm{mod}}}{2}\bigr)/(\log m)^2$
of Remark~\ref{rem:fourth-term} (non-rigorous collocation values). The last column is the
coefficient of $(\log m)^2/m^3$ in the model, for scale: the tabulated quantity would tend to a
nonzero constant if this coefficient were wrong.}\label{tab:fourth-term}
\end{table}

\paragraph{Lean status of Section~\ref{sec:windows-residues}.}
All of Theorems~\ref{thm:windows}, \ref{thm:short-ae}, and~\ref{thm:no-short-dim}---including the exact dimension-$1$ identity for \texttt{prescribed\allowbreak ParitySet}, the exact dimension-$1/2$ identities for \texttt{bounded\allowbreak IncSet}, \texttt{prescribed\allowbreak IncSet}, and \texttt{prescribed\allowbreak IncSetAt}, the almost-everywhere short-window lower bound $Z_J(x)\ge\frac{1-\ee^{-1}}{160}\log J$ (\texttt{thm\_\allowbreak short\_\allowbreak ae\_\allowbreak explicit}), and the full conjunction $\dimH(\texttt{div\allowbreak WindowSet}\,\epsilon)=1\wedge\operatorname{Leb}(\texttt{div\allowbreak WindowSet}\,\epsilon)=0$ (\texttt{thm\_\allowbreak no\_\allowbreak short\_\allowbreak dim\_\allowbreak complete})---are now kernel-checked end-to-end at Layer~A for mathlib's \texttt{GenContFract.of} across thirteen Lean~4 modules: the counting and Frostman-scale pipeline (\texttt{Window\allowbreak Mobius\allowbreak Count.lean}, \texttt{No\allowbreak Short\allowbreak Dim\allowbreak Count.lean}, \texttt{Window\allowbreak Tree\allowbreak Prune.lean}, \texttt{Window\allowbreak Asymp\allowbreak Thresh.lean}, \texttt{Window\allowbreak Frostman\allowbreak Discharge.lean}, and \texttt{Window\allowbreak Paper\allowbreak Scale.lean}), the CF realization and bounded-increment pipeline (\texttt{Window\allowbreak CF\allowbreak Realize.lean}, \texttt{Window\allowbreak CF\allowbreak Bridge.lean}, \texttt{Window\allowbreak No\allowbreak Short\allowbreak CF\allowbreak Realize.lean}, \texttt{Window\allowbreak Prescribed\allowbreak CF\allowbreak Realize.lean}, and \texttt{Window\allowbreak Bounded\allowbreak Upper.lean}), and the short-window measure pipeline (\texttt{Window\allowbreak Short\allowbreak Ae.lean} and \texttt{Window\allowbreak Short\allowbreak Ae\allowbreak Realize.lean}, $14$ declarations; base axioms only, zero \texttt{sorryAx}; see Remark~\ref{rem:windows-lean}).
For Theorems~\ref{thm:residues}, \ref{thm:residues-general}, and~\ref{thm:residues-asymptotic}, the entire general-modulus $m\ge2$ Lipschitz shift equivalence, sandwich, and eventual-periodicity theorem is kernel-checked in \texttt{EpiFormal/\allowbreak Residue\allowbreak ModGeneral.lean} ($7$ declarations: \texttt{lipschitz\allowbreak OnWith\_\allowbreak cfShift}, \texttt{dimH\_\allowbreak residueSeqSet\_\allowbreak sandwich}, \texttt{dimH\_\allowbreak eventually\allowbreak PeriodicResidueSet\_\allowbreak eq}), the unconditional transfer-operator Hausdorff-dimension framework, $\frac12<\dimH F_{m,r}$, $\dimH F_{4,1}\le 4/5$, and step-function certificate soundness (\texttt{step\_\allowbreak cert\_\allowbreak upper\_\allowbreak sound}, \texttt{step\_\allowbreak cert\_\allowbreak lower\_\allowbreak sound}) are kernel-checked in \texttt{EpiFormal/\allowbreak Transfer\allowbreak Operator\allowbreak Dim.lean} ($11$ declarations, base axioms only), and the complete resolvent identity, global strict separation $\dimH F_{m,m}<\dimH F_{m,1}=\dimH E_m$ for all $m\ge2$, two-term asymptotic expansions with $C_r$ and $C_1=\psi_{\mathrm{Fib}}$, exact dimension-excess ratios, and eventual strict monotonicity in~$r$ of Theorem~\ref{thm:residues-asymptotic} are kernel-checked in \texttt{EpiFormal/\allowbreak Residue\allowbreak Asymptotic.lean} (all $15$ declarations: \texttt{metalQ\_\allowbreak one\_\allowbreak eq\_\allowbreak fib}, \texttt{metalC\_\allowbreak one\_\allowbreak eq\_\allowbreak psiFib}, \texttt{resolventV\_\allowbreak branch\_\allowbreak eq}, \texttt{dimH\_\allowbreak residueConstSet\_\allowbreak le\_\allowbreak of\_\allowbreak resolvent}, \texttt{le\_\allowbreak dimH\_\allowbreak residueConstSet\_\allowbreak of\_\allowbreak resolvent}, \texttt{dimH\_\allowbreak residueConstSet\_\allowbreak mm\_\allowbreak le}, \texttt{le\_\allowbreak dimH\_\allowbreak residueConstSet\_\allowbreak mm}, \texttt{dimH\_\allowbreak residueConstSet\_\allowbreak mm\_\allowbreak lt\_\allowbreak m1}, \texttt{dimH\_\allowbreak residueConstSet\_\allowbreak mm\_\allowbreak lt\_\allowbreak eventuallyPeriodic}, \texttt{dimH\_\allowbreak residueConstSet\_\allowbreak asymp\_\allowbreak two\_\allowbreak term}, \texttt{dimH\_\allowbreak residueConstSet\_\allowbreak mm\_\allowbreak asymp\_\allowbreak two\_\allowbreak term}, \texttt{tendsto\_\allowbreak mul\_\allowbreak dimH\_\allowbreak residueConstSet\_\allowbreak sub\_\allowbreak half}, \texttt{tendsto\_\allowbreak dimH\_\allowbreak residue\_\allowbreak ratio\_\allowbreak eq\_\allowbreak metalC}, \texttt{tendsto\_\allowbreak dimH\_\allowbreak eventuallyPeriodic\_\allowbreak ratio\_\allowbreak eq\_\allowbreak psiFib}, and \texttt{eventually\_\allowbreak dimH\_\allowbreak residueConstSet\_\allowbreak strictAnti}; base axioms only, zero \texttt{sorryAx}), alongside the $14$ helper lemmas in \texttt{Window\allowbreak And\allowbreak Residue\allowbreak Core.lean}; the integer grid certificates for the narrow brackets in Table~\ref{tab:residues-general} are run in exact Python arithmetic (\texttt{residue\_\allowbreak certificate\_\allowbreak general\_\allowbreak N512.py}) and, as closed brackets, re-checked by the Lean kernel in \texttt{Residue\allowbreak Cert\allowbreak Table.lean} (Remark~\ref{rem:residue-kernel}). Moreover, all parts~(i)--(v) of Theorem~\ref{thm:multi-residue-asymptotic} (Bowen-free proof of the Poincar\'e-abscissa dimension formula, multi-residue phase transition at $1/2$ with two-digit witnesses $\{2,3\}$ and $\{1,2\}$, elementary Euler--Mascheroni expansion of $\zeta(1+\varepsilon)$, three-term subcritical asymptotics to $O((\log m)^2/m^3)$, and supercritical $\Theta(m^{-2s_0})$ convergence rate) are kernel-checked at Layer~A in \texttt{EpiFormal/\allowbreak ResidueMultiAsymptotic.lean} ($24$ public theorems audited separately in \texttt{AuditResidueMultiAsymptotic.lean}, strictly on the three Lean base axioms \texttt{propext}, \texttt{Classical.choice}, and \texttt{Quot.sound}, with zero \texttt{sorryAx}, recorded in \texttt{verification/\allowbreak audit\_\allowbreak residue\_\allowbreak multi\_\allowbreak asymptotic\_\allowbreak 20261123.log}).

\section{Numerical evidence}\label{sec:num}

All high-precision computations use \texttt{mpmath} with dynamic precision
$\mathrm{digits}(Q_j)+300$ (or 100 decimal digits for fixed-horizon checks up
to $j=35$) to survive the catastrophic cancellation in
$Q_jx-\mathrm{round}(Q_jx)$. These are corroborations, not proofs.

\paragraph{Spectral gap.} For $x=B\pi$ ($B=1,\dots,10$), $\sqrt2$, the golden
ratio $\varphi$, and random reals, $\max_{j\le60}j\|Q_jx\|$ lies in $22$--$29$
and grows $\sim j/4$ in this finite sample; the data do not establish an
infinite limsup. For $x\in L$ (e.g.\ $(\ee+1)/2$, $1+(\ee+1)/2$),
$j\|Q_jx\|<10^{-12}$ already for $j\ge10$ and $\sim10^{-117}$ at $j=60$,
confirming $\Lam=0$.

\paragraph{Dimension (box-counting).} Exact rational covering trees give the
following box-counting estimates for $E_c\cap[0,1]$:
\begin{center}
\begin{tabular}{c|cccc}
$c$ & $N=4$ & $N=12$ & path $N=400$ & bound $n=1500$\\\hline
$1/4$ & 0.371 & 0.323 & 0.171 & 0.180\\
$1/2$ & 0.630 & --- & 0.251 & 0.270\\
$1$   & 0.863 & --- & 0.343 & 0.360\\
$2$   & 1.000 & --- & 0.442 & 0.446
\end{tabular}
\end{center}
The finite counts suggest the branching scale $\lfloor8c\rfloor+1$ ($3,5,9,17$), and
the rigorous upper bound decays as $\log B/\log n\to0$. \textbf{Caveat:} the
decay is only $O(1/\log n)$; at accessible $N$ the estimates are far from $0$
(e.g.\ $s=0.01$ needs $n\approx10^{60}$). The numerics are consistent with
$\dimH=0$ but cannot resolve it directly; the proof is the covering argument of
Theorem~\ref{thm:finite}. Below the gap ($c=0.24$) the covering upper bound
still grows exponentially while the true branching tends to $0$: the covering
method is not sharp below the gap, which is why the exact classification
(Theorem~\ref{thm:class}) is needed there.

\paragraph{Discrete spectral ladder and three-way splitting.} The 100-digit
\texttt{mpmath} script \texttt{verify\_\allowbreak full\_\allowbreak spectrum\_\allowbreak discrete.py} verifies
the quantized levels $\Lam(x_k)=k/4$ for $k=1,2,3,4$ (with $|4ju_j(x_k)-k|\approx7.14\times10^{-3}\cdot k$ at $j=35$), confirms that two points in
$E_{\mathrm{lim}}^+(k/4)$ with distinct initial digits have difference in $L$ to
within $5.95\times10^{-97}$, and checks that the alternating binary digit choice
$d_j=(-1)^j k\in\{-k,+k\}$ in $E_{\mathrm{lim}}(k/4)$ outside $E_{\mathrm{lim}}^+(k/4)\cup E_{\mathrm{lim}}^-(k/4)$ satisfies $4j\|Q_jx\|\to k$ while $4ju_j(x)$ oscillates
between $+k$ and $-k$. The companion script \texttt{verify\_\allowbreak weighted\_\allowbreak lower.py}
confirms the Falconer tree-density ratio convergence toward $1-\alpha$ for
subcritical weights $\alpha\in\{0.25,0.50,0.75\}$.

\paragraph{General polynomial-growth recurrences.} The script
\texttt{verify\_\allowbreak poly\_\allowbreak recurrence\_\allowbreak 2b.py}
(at least $400$-digit effective \texttt{mpmath} precision for Experiments~1--3
and $50$-digit precision up to $N=10^5$ for Experiments~4--5) passes all
assertions across $r\in\{1,2,3\}$ and $c\in\{1,2,4\}$: it verifies the
non-asymptotic factorial lower bound of Lemma~\ref{lem:poly-growth}, the
$O(j^{-r})$ convergence of $j^r\|Q_jx\|\to k/c$ for $k\in\{1,3\}$, the
intrinsic lattice $L^{(r,c,d)}=\frac{1}{g}(\Z+\Z\theta)$ (including the
non-coprime case $(2,4,0)$ with $Q_0=3,Q_1=6,g=3$, where
$q_0\alpha_U+q_1\alpha_V=1/3$ to within $5.14\times10^{-1079}$ and $x_2-x_1\in
L^{(2,4,0)}$ to within $6.0\times10^{-1072}$) alongside alternating-digit
oscillation in $E_{\mathrm{lim}}^{(r,c,d)}(k/c)$, the subcritical Falconer
dimension trajectories approaching $1-\alpha/r$ (for both fixed $t>0$ and the
slowly decaying $t=0$ radius) versus decay toward $0$ at critical $\alpha=r$
($s_{10^5}\approx0.033$ for $(2,1,-1)$ and $0.056$ for $(3,2,-3)$), and the
exact gauge ratio $\log(2k+1)/\log2$ at $n=10^5$ (with the $17.6\%$--$39.5\%$
finite-$n$ excess over the asymptotic scale factor $1/r$ matching the
$O(\log\log n/\log n)$ subleading term $\frac{\log n+\log\log n+\log r}{\log n+(\log c-r)/r}$
to within $0.009$).

\paragraph{Reciprocal-Fibonacci residue asymptotics and growing-quotient phase transitions.}
The $50$-digit \texttt{mpmath} script
\texttt{verify\_\allowbreak paper2\_\allowbreak extensions.py} (logged in
\texttt{verification/\allowbreak verify\_\allowbreak paper2\_\allowbreak extensions.log})
computes $C_r$ ($1\le r\le 5$, Table~\ref{tab:residues-asymptotic}) to $50$ decimal
places, verifies the global separation sandwich
$\frac12+\frac{1}{2m}<\frac12+\frac{11}{16m}$ alongside the certified bounds of
Table~\ref{tab:residues-general}, checks the convergence of the two-term
excess ratio $(\dimH F_{m,r}-1/2)/(\dimH F_{m,m}-1/2)\to C_r$ up to $m=20000$
(reaching $3.355958$ vs.\ $\psi_{\mathrm{Fib}}\approx3.359886$ at $r=1$), and
verifies all three growing-quotient regimes of
Corollary~\ref{cor:growing-rec-three-regimes} ($\log A_j/\log Q_j\to 0$, $1$,
and $\infty$, respectively), including the supercritical behavior of
$A_j^\tau\|Q_j\theta_G\|$ on the double-exponential recurrence $A_j=2^{2^j}$ at
$\tau=1.5$ ($1.71\times10^{-19}$ at $j=7$), $\tau=2.0$ (stabilizing to
$(\lim Q_j/A_j)^{-1}\approx 3.15973\in[1/6,4]$), and $\tau=2.5$
($5.83\times10^{19}$ at $j=7$).

\paragraph{Multi-residue phase transition and Euler--Mascheroni three-term asymptotics.}
The $32$-digit \texttt{mpmath} Chebyshev-collocation script
\texttt{verification/\allowbreak verify\_\allowbreak residue\_\allowbreak multi\_\allowbreak asymptotic.py}
(\texttt{ALL CHECKS PASSED}) independently corroborates Theorem~\ref{thm:multi-residue-asymptotic}:
all ten $N=512$ kernel brackets of Table~\ref{tab:residues-general} are
satisfied by $n=16$ vs.\ $n=22$ Chebyshev estimates; across ten subcritical
single- and multi-residue alphabets evaluated up to $m=16384$
(Table~\ref{tab:multi-residue-asymptotic}), the second-order residual
$m^2\bigl(\dimH F_{m,S,B}-\frac12-\frac{dC_S}{2m}+\frac{(dC_S)^2\log m}{2m^2}\bigr)$
converges monotonically toward the theoretical constant
$c_2=\frac12 d^2 C_S(\gamma C_S-C_S')$ with bounded third-order ratio
$m^3(\dimH-\text{3-term})/(\log m)^2$; the algebraic supercritical test on
$S=\{1,2\}$, $s=51/100$, $v(x)=2/(2+x)$ achieves
$\min_{[0,1]}(\mathcal L_{51/100,\{1,2\}}v/v)\approx1.01099707>1$; and the
supercritical pairs $S=\{1,2\}$ ($s_0\approx0.531280506$) and $S=\{1,2,3\}$
($s_0\approx0.705660908$) exhibit the predicted $\Theta(m^{-2s_0})$ excess
scaling up to $m=4096$.

\section{Honesty statement, scope, and formal verification}\label{sec:verification}

\begin{itemize}[leftmargin=1.5em]
\item Theorems~\ref{thm:class} and~\ref{thm:gap} have short self-contained paper
proofs in Section~\ref{sec:class} and are kernel-checked at Layer~A in
\texttt{EpiFormal/\allowbreak Classification.lean}
(\texttt{theorem2\_classification} and \texttt{corollary3\_gap}), using only the
three Lean base axioms \texttt{propext}, \texttt{Classical.choice} and
\texttt{Quot.sound}; the former $\delta_j$ axioms are now proved theorems
(\texttt{DeltaSign.lean}).
\item Corollaries~\ref{thm:spectral} and~\ref{thm:sub} (for $c<1/4$), and the
countability and dimension-$0$ statements for $c<1/4$, are kernel-checked at
Layer~A in \texttt{EpiFormal/\allowbreak SpectralGap.lean} (five declarations,
no \texttt{sorryAx}). The complete spectrum discretization
$\Lam(\R)=\frac14\Z_{\ge0}\cup\{\infty\}$ and open-interval spectral-gap
exclusion (Theorem~\ref{thm:full-spectrum}), the gap constancy $E_c=E_{k/4}$
(Corollary~\ref{thm:sub}), the three-way digit rigidity and signed-limit
$L$-difference, countability, and dimension-$0$ lemmas
(Theorem~\ref{thm:trichotomy}), and the $r=1$ linear $(cj+d)$ defect
approximation lemmas are kernel-checked at Layer~A in
\texttt{EpiFormal/\allowbreak FullSpectrumDiscrete.lean} (15 audited
declarations, base axioms only, no \texttt{sorryAx}). The remaining signed-coset
collapse $E_{\mathrm{lim}}^\pm(k/4)=x_k^\pm+L$ in
Theorem~\ref{thm:trichotomy}(iii) is now likewise kernel-checked at Layer~A in
\texttt{EpiFormal/\allowbreak TrichotomyAndPolyExtension.lean} (8 declarations,
ending in \texttt{trichotomy\_\allowbreak part\_\allowbreak iii\_\allowbreak complete},
base axioms only, no \texttt{sorryAx}), completing Layer~A verification of all
three parts~(i)--(iii) of Theorem~\ref{thm:trichotomy}.
\item For the general polynomial-growth denominator recurrences ($r\ge1$) of
Section~\ref{sec:poly-rec}, the entire theory---Lemmas~\ref{lem:poly-growth}(i)--(iv)
and~\ref{lem:poly-lattice}, Theorem~\ref{thm:general-rec} (full spectrum
$\frac{1}{c}\Z_{\ge0}\cup\{\infty\}$ for all $r\ge1$),
Theorem~\ref{thm:poly-gauge-trichotomy}(i)--(iii) (critical-gauge and
binary-gauge metric measure trichotomy on $\R$ and $\R/\Z$ plus signed-coset
collapse), Corollary~\ref{cor:poly-exact-gauge},
Remark~\ref{rem:poly-gauge-scale}, and
Theorem~\ref{thm:poly-phase-transition}(i)--(iii) (critical $\alpha=r$,
subcritical $0\le\alpha<r$ exact dimension $1-\alpha/r$, and supercritical
$\alpha>r$ collapse)---is kernel-checked at Layer~A across
\texttt{EpiFormal/\allowbreak TrichotomyAndPolyExtension.lean} ($37$
declarations), \texttt{EpiFormal/\allowbreak PolyGaugeAndPhaseTransition.lean}
($22$ declarations), \texttt{EpiFormal/\allowbreak PolyWeightedUpper.lean} ($1$
declaration), and \texttt{EpiFormal/\allowbreak PolyWeightedLower.lean} ($8$
declarations; all $68$ declarations on base axioms only, zero \texttt{sorryAx}),
and corroborated by high-precision \texttt{mpmath} checks. Furthermore, the
universal extension to arbitrary growing-quotient recurrences
$Q_{j+1}=A_jQ_j+Q_{j-1}$ ($A_j\ge2$, $A_j\to\infty$) in
Section~\ref{sec:growing-rec}---Theorem~\ref{thm:growing-rec-spectrum}(i)--(iii)
(Kraaikamp--Liardet ratio identity, zero-level TFAE, full discrete spectrum
$\Z_{\ge0}\cup\{\infty\}$, and three-way positive-level splitting and critical
gauges on $\R$ and $\R/\Z$), Theorem~\ref{thm:growing-rec-dim-phase}(i)--(ii)
(subcritical dimension $(1-\tau)/(1+\tau\gamma_{\sup})$ and supercritical
secondary phase transition between $L_G$ and $\frac{1}{g}\Z$ at
$\tau_{\mathrm{crit}}=1+1/\gamma_{\sup}$), and
Corollary~\ref{cor:growing-rec-three-regimes}(i)--(iii) (Regimes~I--III including
the critical boundary $\tau=2$ for $A_j=2^{2^j}$)---is kernel-checked at Layer~A
across \texttt{EpiFormal/\allowbreak General\allowbreak Growing\allowbreak Rec.lean}
($21$ declarations) and \texttt{EpiFormal/\allowbreak General\allowbreak Growing\allowbreak Dim.lean}
($29$ declarations; all $50$ declarations on base axioms only, zero \texttt{sorryAx}).
\item Theorem~\ref{thm:finite} ($\dimH E_c=0$ for all $c\in\R$) is \emph{fully
machine-checked at Layer~A} (\texttt{DimHCore.lean} + \texttt{DimHFinite.lean},
base axioms only: \texttt{propext}, \texttt{Classical.choice},
\texttt{Quot.sound}; no $\delta$ axioms, no \texttt{sorryAx}). The assembly uses
a surviving-center \texttt{Finset} recursion (\texttt{survSet}) rather than path
encoding; the growth estimate is explicit
(\texttt{tendsto\_pow\_div\_factorial\_atTop}), not asymptotic. Moreover, the
complementary infinite-level conclusions of
Corollary~\ref{cor:infinite-level}---Lebesgue-nullity of $\{\Lam<\infty\}$,
full Lebesgue measure on every interval $[a,b]$ and a.e.\ $\Lam=\infty$, and
$\dimH\{\Lam=\infty\}=1$---are likewise kernel-checked at Layer~A in
\texttt{EpiFormal/\allowbreak{}GaugeExact.lean}
(\texttt{volume\_finiteLevel\_zero}, \texttt{ae\_apprExponent\_eq\_top},
\texttt{volume\_infiniteLevel\_Icc}, and
\texttt{dimH\_infiniteLevel\_eq\_one}).
\item Theorems~\ref{thm:upper} ($0\le\alpha\le1$, all $c\in\R$)
and~\ref{conj:weight} ($0\le\alpha<1$, prescribed limits, sublevel sets for all
$c\ge0$, and the zero-level set) are \emph{fully machine-checked at Layer~A}
(both the upper and lower bounds). The upper bound is proved in
\texttt{Weighted\allowbreak Fiber.lean}, \texttt{Weighted\allowbreak DimHCore.lean},
\texttt{Weighted\allowbreak Upper.lean}, and \texttt{Weighted\allowbreak Final.lean}
(as \texttt{dimH\_\allowbreak levelSet\ensuremath{\alpha}\_le}); the
matching lower bound and exact dimension identities are proved in
\texttt{Weighted\allowbreak ExactDim.lean}
(as \texttt{one\_\allowbreak le\_\allowbreak hausdorff\allowbreak Measure\_\allowbreak cantorSet},
\texttt{dimH\_\allowbreak exact\allowbreak LimitSet\ensuremath{\alpha}\_ge},
\texttt{dimH\_\allowbreak exact\allowbreak LimitSet\ensuremath{\alpha}\_eq},
\texttt{dimH\_\allowbreak levelSet\ensuremath{\alpha}\_eq}, and
\texttt{dimH\_\allowbreak appr\allowbreak Exponent\ensuremath{\alpha}\_eq\_\allowbreak zero}; base
axioms only, no \texttt{sorryAx}). The existence of $x_0$ with $\Lam(x_0)=1/4$
(as \texttt{lambda\_\allowbreak x0\_\allowbreak eq\_\allowbreak quarter} in \texttt{Quarter\allowbreak Point\allowbreak Final.lean}) and the
unweighted-zero strictness witness
(\texttt{Unweighted\allowbreak Correction.\allowbreak witness} in Remark~\ref{rem:unweighted})
are likewise kernel-checked at Layer~A.
\item Theorem~\ref{thm:gauge}, Corollary~\ref{cor:exact-gauge},
Remark~\ref{rem:gauge-scale}, and Theorem~\ref{thm:trichotomy}(i)--(iii) are now
\emph{fully machine-checked at Layer~A} both on $\R$ and (for the metric gauge
statements of parts~(i)--(ii)) natively on the
quotient metric space $\R/\Z$ (\texttt{UnitAddCircle}) across
\texttt{EpiFormal/\allowbreak GaugeNull.lean} (3 declarations),
\texttt{EpiFormal/\allowbreak GaugeExact.lean} (22 declarations),
\texttt{EpiFormal/\allowbreak GaugeCircle.lean} (53 declarations),
\texttt{EpiFormal/\allowbreak FullSpectrumDiscrete.lean} (15 declarations), and
\texttt{EpiFormal/\allowbreak TrichotomyAndPolyExtension.lean} (8 part-(iii)
declarations), using only
the three Lean base axioms \texttt{propext}, \texttt{Classical.choice}, and
\texttt{Quot.sound}, with zero \texttt{sorryAx}. In particular,
\texttt{GaugeCircle.lean} explicitly constructs the piecewise-linear gauges
$h_k$ (\texttt{critGauge}), $g_2$ (\texttt{binaryGauge}), and $g_\beta$
(\texttt{interpGauge}), discharges all external continuity hypotheses for
Theorem~\ref{thm:gauge} and Remark~\ref{rem:gauge-scale}, constructively proves
open-interval and open-arc $\mathcal H^{h_k}$-infinitude
(Corollary~\ref{cor:exact-gauge}(i) on $\R$ and \texttt{UnitAddCircle}), and
verifies the binary Cantor-tree positivity, $\mathcal H^{g_2}$-infinitude on
every nonempty open subset of $\R/\Z$, $\sigma$-finiteness,
$\mathcal H^{h_k}$-nullity, and $g_\beta$ threshold dichotomy for
$E_{\mathrm{lim}}(k/4)$ in Theorem~\ref{thm:trichotomy}(ii), while
\texttt{TrichotomyAndPolyExtension.lean} proves the signed-limit coset collapse
$E_{\mathrm{lim}}^\pm(k/4)=x_k^\pm+L$ in Theorem~\ref{thm:trichotomy}(iii).
\item Theorem~\ref{thm:finite-state-tail} (the marked-component critical gauge
theorem, both on $\R$ via \texttt{thm\_\allowbreak finite\_\allowbreak
state\_\allowbreak tail\_\allowbreak complete} and on $\R/\Z$ via
\texttt{thm\_\allowbreak finite\_\allowbreak state\_\allowbreak tail\_\allowbreak
circle\_\allowbreak complete}) is now \emph{fully machine-checked at Layer~A}
in \texttt{EpiFormal/\allowbreak Finite\allowbreak State\allowbreak Complete.lean}
(12 audited declarations, base axioms only, zero \texttt{sorryAx}).
\item For the continued-fraction window and residue results
(Theorems~\ref{thm:windows}--\ref{thm:residues} and~\ref{thm:residues-general}),
the entirety of Theorem~\ref{thm:windows} (both the exact $\dimH=1$
prescribed-parity statement and the exact $\dimH=1/2$ upper and lower bounds
for bounded/prescribed increments), the Lebesgue-measure zero-tail law of
Theorem~\ref{thm:short-ae} (explicit positive constant $c = (1-\ee^{-1})/160$ in
\texttt{thm\_\allowbreak short\_\allowbreak ae\_\allowbreak explicit}), and
both parts of Theorem~\ref{thm:no-short-dim} ($\dimH=1$ and
$\operatorname{Leb}=0$, via
\texttt{volume\_\allowbreak divWindowSet\_\allowbreak eq\_\allowbreak zero}
and \texttt{thm\_\allowbreak no\_\allowbreak short\_\allowbreak dim\_\allowbreak
complete}, together with M\"obius counting cores, paper-scale Frostman
discharge, concrete stopped-cylinder CF trees, pointwise regular CF
identification with mathlib's \texttt{GenContFract.of}, and level-$n$ cylinder
$\zeta(2t)^n$ covering) are kernel-checked at Layer~A across
\texttt{Window\allowbreak Mobius\allowbreak Count.lean},
\texttt{No\allowbreak Short\allowbreak Dim\allowbreak Count.lean},
\texttt{Window\allowbreak Frostman\allowbreak Discharge.lean},
\texttt{Window\allowbreak Paper\allowbreak Scale.lean},
\texttt{Window\allowbreak CF\allowbreak Realize.lean},
\texttt{Window\allowbreak CF\allowbreak Bridge.lean},
\texttt{Window\allowbreak No\allowbreak Short\allowbreak CF\allowbreak Realize.lean},
\texttt{Window\allowbreak Prescribed\allowbreak CF\allowbreak Realize.lean},
\texttt{Window\allowbreak Bounded\allowbreak Upper.lean},
\texttt{Window\allowbreak Short\allowbreak Ae.lean},
\texttt{Window\allowbreak Short\allowbreak Ae\allowbreak Realize.lean}, and the fourteen
core lemmas in \texttt{Window\allowbreak And\allowbreak Residue\allowbreak Core.lean}
(base axioms only, zero \texttt{sorryAx}; see Remark~\ref{rem:windows-lean}).
For Theorems~\ref{thm:residues} and~\ref{thm:residues-general}, the
transfer-operator Hausdorff-dimension bridge, the qualitative gap
$\dimH > 1/2$ for $E_{\equiv 1\pmod 4}$, the proved upper bound
$\dimH \le 4/5$, the finite-step matrix-certificate soundness theorems
(\texttt{step\_\allowbreak cert\_\allowbreak upper\_\allowbreak sound} and
\texttt{step\_\allowbreak cert\_\allowbreak lower\_\allowbreak sound}), and
the M\"obius/eventual/periodic residue invariance identities are kernel-checked
in \texttt{Residue\allowbreak ModGeneral.lean} and
\texttt{Transfer\allowbreak Operator\allowbreak Dim.lean}, while the sharp
$512$-point dyadic numerical intervals in Theorem~\ref{thm:residues} and
Table~\ref{tab:residues-general} are certified by the reproducible exact
integer-arithmetic script \texttt{residue\_\allowbreak certificate\_\allowbreak general\_\allowbreak N512.py},
and their closed versions are kernel-checked in
\texttt{Residue\allowbreak Cert\allowbreak Table.lean} (Remark~\ref{rem:residue-kernel}). Finally,
Theorem~\ref{thm:residues-asymptotic}(i)--(iv) (exact metallic-continuant
resolvent identity, global strict separation $\dimH F_{m,m}<\dimH F_{m,1}=\dimH E_m$
for all $m\ge2$, two-term asymptotic expansions with $C_r$ and $C_1=\psi_{\mathrm{Fib}}$,
exact dimension-excess ratios, and eventual strict monotonicity in~$r$) is
kernel-checked at Layer~A in
\texttt{EpiFormal/\allowbreak Residue\allowbreak Asymptotic.lean} ($15$ audited
declarations, base axioms only, zero \texttt{sorryAx}) and corroborated by
\texttt{verify\_\allowbreak paper2\_\allowbreak extensions.py}. Furthermore,
all five parts~(i)--(v) of Theorem~\ref{thm:multi-residue-asymptotic}
(Bowen-free proof of the classical Poincar\'e-abscissa dimension formula for
finite alphabets, universal
multi-residue phase transition at $1/2$ with two-digit witnesses $\{2,3\}$
and $\{1,2\}$, elementary Euler--Mascheroni expansion of $\zeta(1+\varepsilon)$,
subcritical three-term asymptotics to $O((\log m)^2/m^3)$ with coefficient
$\frac12 d^2 C_S(\gamma C_S-C_S')$, and supercritical $\Theta(m^{-2s_0})$
convergence rate) are kernel-checked at Layer~A in
\texttt{EpiFormal/\allowbreak Residue\allowbreak Multi\allowbreak Asymptotic.lean}
($24$ public theorems audited separately in
\texttt{AuditResidueMultiAsymptotic.lean}, base axioms only, zero \texttt{sorryAx},
log \texttt{verification/\allowbreak audit\_\allowbreak residue\_\allowbreak multi\_\allowbreak asymptotic\_\allowbreak 20261123.log})
and numerically corroborated by
\texttt{verify\_\allowbreak residue\_\allowbreak multi\_\allowbreak asymptotic.py}.
\item This article does \emph{not} prove the irrationality of $e+\pi$ (which
remains an open problem). It is a metric-theory spin-off; the connection is
that the zero-level classification machinery arose from the $e+\pi$ escape
program \cite{Phase1}.
\item The zero-level classification and initial gap (Theorems~\ref{thm:class}
and~\ref{thm:gap}) are a subsequence variant of the Kraaikamp--Liardet
criterion \cite{KL}, quoted as Lemma~4.1 of \cite{BDS}; see
Section~\ref{sec:prior}. The independent contributions concern the complete
spectrum discretization and three-way positive-level splitting
(Theorems~\ref{thm:full-spectrum}, \ref{thm:trichotomy},
\ref{thm:general-rec}--\ref{thm:poly-phase-transition},
and~\ref{thm:growing-rec-spectrum}--\ref{thm:growing-rec-dim-phase}, all kernel-checked
in Lean~4), the critical gauge theorems
(Theorem~\ref{thm:gauge} and Corollaries~\ref{cor:exact-gauge}
and~\ref{cor:liminf-gauges}), the kernel-checked marked-component gauge
(Theorem~\ref{thm:finite-state-tail}, mechanized in
\texttt{Finite\allowbreak State\allowbreak Complete.lean}), the machine-checked
dimension-zero and subcritical dimension theorems
(Theorems~\ref{thm:finite}, \ref{thm:upper}, and~\ref{conj:weight}), the
window and residue theorems (Theorems~\ref{thm:windows}--\ref{thm:residues-general}
and Theorem~\ref{thm:residues-asymptotic}(ii)), and explicit residue-class
instances of known asymptotic methods
(Theorems~\ref{thm:residues-asymptotic}(iii)--(iv)
and~\ref{thm:multi-residue-asymptotic}; see Section~\ref{sec:prior}).
\end{itemize}

\section*{Use of AI and formal verification}
% TODO before submission: adapt this statement to the target journal's AI-use policy.
This work was carried out with substantial assistance from generative AI
systems (including ChatGPT, Codex, and Claude), both for developing and drafting
mathematical arguments and text, preparing the numerical certificates, and
writing the Lean~4 formalization. The results marked as machine-checked (Layer~A)
were verified by the Lean~4 kernel, and their axiom dependencies were audited
with \texttt{\#print axioms} in \texttt{Audit.lean} ($937$ audited declarations
across $134$ Lean files as of 2026-11-22, $9179$ build jobs, of which $23$ additionally use the three whitelisted classical
axioms \texttt{e\_transcendental}, \texttt{pi\_transcendental}, and
\texttt{mu\_e\_eq\_two} in the companion $e+\pi$ paper's modules; the other $914$
depend only on the Lean base axioms \texttt{propext}, \texttt{Classical.choice},
and \texttt{Quot.sound}---with zero
project-added axioms and zero \texttt{sorryAx}); the $27$ declarations of the
residue-certificate layer (Remark~\ref{rem:residue-kernel}) are audited separately in
\texttt{AuditResidueCert.lean}, and the $24$ public theorems of the multi-residue
phase-transition and Euler--Mascheroni three-term asymptotic layer
(\texttt{EpiFormal/\allowbreak ResidueMultiAsymptotic.lean},
Theorem~\ref{thm:multi-residue-asymptotic}) are audited separately in
\texttt{AuditResidueMultiAsymptotic.lean} (logged in
\texttt{verification/\allowbreak audit\_\allowbreak residue\_\allowbreak multi\_\allowbreak asymptotic\_\allowbreak 20261123.log}),
again on the three Lean base axioms only with zero \texttt{sorryAx} (total
$344$ Lean files across the repository with zero escape keywords). The build and main audit were independently reproduced on 2026-11-23 from a
clean clone of the public repository on a separate machine (Lean~4.35.0,
Mathlib v4.35.0): \texttt{lake build} succeeded, the $937$
\texttt{\#print axioms} outputs coincide entry by entry with the recorded log, and the $27$ residue-certificate entries depend only on the base axioms
(\texttt{verification/\allowbreak colab\_\allowbreak rebuild\_\allowbreak audit\_\allowbreak 20261123.md}), and on 2026-11-24 a second clean-clone rebuild on a separate machine (commit \texttt{d920a81}, which contains \texttt{ResidueMultiAsymptotic.lean}) reproduced the $937$ main entries, the $27$ residue-certificate entries and the $24$ entries of \texttt{AuditResidueMultiAsymptotic.lean} entry by entry, all $27+24$ on the three base axioms only (\texttt{verification/\allowbreak colab\_\allowbreak rebuild\_\allowbreak audit\_\allowbreak 20261124.md}). The portions not resolved here or lying outside
the Lean kernel are the finite-dimensional linear-algebraic Perron--Frobenius
eigenvector existence input on marked strongly connected components in
Theorem~\ref{thm:finite-state-tail} (whose positive right eigenvector is passed
explicitly to \texttt{thm\_\allowbreak finite\_\allowbreak state\_\allowbreak tail\_\allowbreak complete}
and \texttt{thm\_\allowbreak finite\_\allowbreak state\_\allowbreak tail\_\allowbreak circle\_\allowbreak complete}
in \texttt{FiniteStateComplete.lean}), the strictness of the open intervals in
Theorem~\ref{thm:residues} and Table~\ref{tab:residues-general} via the
cited pressure formalism (the closed brackets themselves are kernel-checked,
Remark~\ref{rem:residue-kernel}), the asymptotic rate on the critical boundary
$s_0(S)=1/2$ and the exact supercritical leading constant in
Theorem~\ref{thm:multi-residue-asymptotic}, and the irrationality of $e+\pi$,
which remains an open problem.
% TODO before submission: state who checked the paper proofs, and how.

\begin{thebibliography}{99}
\bibitem{BDS} A.~B\'ir\'o, J.-M.~Deshouillers, V.~T.~S\'os, \emph{Good
approximation and characterization of subgroups of $\R/\Z$}, Studia Sci.\ Math.\ Hungar.\ 38 (2001), 97--113, doi:10.1556/\allowbreak sscmath.38.2001.1-4.7.

\bibitem{BV} V.~Beresnevich and S.~Velani, \emph{A mass transference principle
and the Duffin--Schaeffer conjecture}, Ann.\ of Math.\ (2) 164 (2006), 971--992.

\bibitem{BLR} Y.~Bugeaud, L.~Liao, M.~Rams, \emph{Metrical results on the
distribution of fractional parts of powers of real numbers}, Proc.\ Edinb.\ Math.\ Soc.\ (2) 62 (2019), 505--521, doi:10.1017/\allowbreak S0013091518000585; arXiv:1710.03596. (Cor.~8, the hypothesis is $Q_{j+1}/Q_j\to\infty$.)

\bibitem{CZ} W.~Cheng and X.~Zhang, \emph{Exact approximation order of real numbers in
Cantor series expansions}, arXiv:2606.30435v2 (2026).

\bibitem{CMMS} C.~Cabrelli, F.~Mendivil, U.~M.~Molter, R.~Shonkwiler,
\emph{On the Hausdorff $h$-measure of Cantor sets}, Pacific J.\ Math.\ 217
(2004), 45--59, doi:10.2140/\allowbreak pjm.2004.217.45.

\bibitem{CLU} V.~Chousionis, D.~Leykekhman, M.~Urba\'nski, \emph{On the
dimension spectrum of infinite subsystems of continued fractions}, Trans.\ Amer.\ Math.\ Soc.\ 373 (2020), 1009--1042, arXiv:1805.11904.

\bibitem{DFSU} T.~Das, L.~Fishman, D.~Simmons, M.~Urba\'nski, \emph{Hausdorff
dimensions of perturbations of a conformal iterated function system via
thermodynamic formalism}, Selecta Math.\ (N.S.) 29 (2023), no.~2, Paper No.~19,
doi:10.1007/\allowbreak s00029-022-00820-z; arXiv:2007.10554.

\bibitem{DLMF} \emph{NIST Digital Library of Mathematical Functions}, \S25.2,
Release~1.2.8 of 2026-09-15, \url{https://dlmf.nist.gov/25.2}.

\bibitem{Falconer} K.~Falconer, \emph{Fractal Geometry: Mathematical
Foundations and Applications}, 2nd ed., Wiley, 2003. (Countable sets have
Hausdorff dimension~$0$.)

\bibitem{FLLW} A.-H.~Fan, L.-M.~Liao, B.-W.~Wang, J.~Wu, \emph{On Khintchine
exponents and Lyapunov exponents of continued fractions}, Ergodic Theory Dynam.\ Systems 29 (2009), 73--109.

\bibitem{FWW} D.-J.~Feng, Z.-Y.~Wen, J.~Wu, \emph{Some dimensional results
for homogeneous Moran sets}, Sci.\ China Ser.\ A 40 (1997), 475--482,
doi:10.1007/\allowbreak BF02896955.

\bibitem{Good41} I.~J.~Good, \emph{The fractional dimensional theory of continued
fractions}, Proc.\ Cambridge Philos.\ Soc.\ 37 (1941), 199--228,
doi:10.1017/\allowbreak S030500410002171X.

\bibitem{HU02} S.~Heinemann and M.~Urba\'nski, \emph{Hausdorff dimension estimates
for infinite conformal IFSs}, Nonlinearity 15 (2002), 727--734,
doi:10.1088/\allowbreak 0951-7715/\allowbreak 15/\allowbreak 3/\allowbreak 312.

\bibitem{Hensley92} D.~Hensley, \emph{Continued fraction Cantor sets, Hausdorff
dimension, and functional analysis}, J.\ Number Theory 40 (1992), 336--358,
doi:10.1016/\allowbreak 0022-314X(92)90006-B.

\bibitem{KL} C.~Kraaikamp and P.~Liardet, \emph{Good approximations and
continued fractions}, Proc.\ Amer.\ Math.\ Soc.\ 112 (1991), 303--309, doi:10.1090/\allowbreak s0002-9939-1991-1062392-7.

\bibitem{LimaMoreira} D.~Lima and C.~G.~Moreira, \emph{Phase transitions on
the Markov and Lagrange dynamical spectra}, Ann.\ Inst.\ H.~Poincar\'e
C Anal.\ Non Lin\'eaire 38 (2021), 1429--1459,
doi:10.1016/\allowbreak j.anihpc.2020.11.007; arXiv:1801.04636.

\bibitem{Moreira} C.~G.~Moreira, \emph{Geometric properties of the Markov and
Lagrange spectra}, Ann.\ of Math.\ (2) 188 (2018), 145--170,
doi:10.4007/\allowbreak annals.2018.188.1.3.

\bibitem{MEA} C.~G.~Moreira, H.~Erazo, N.~Angelini,
\emph{Generalized Hausdorff measure of irrationals with Lagrange value exactly 3},
arXiv:2511.10328v3.

\bibitem{MU96} R.~D.~Mauldin and M.~Urba\'nski, \emph{Dimensions and measures in
infinite iterated function systems}, Proc.\ London Math.\ Soc.\ (3) 73 (1996),
105--154, doi:10.1112/\allowbreak plms/\allowbreak s3-73.1.105.

{\sloppy
\bibitem{MW} R.~D.~Mauldin and S.~C.~Williams,
\emph{Hausdorff dimension in graph directed constructions},
Trans.~Amer.~Math.~Soc.~309 (1988), no.~2, 811--829.
}

\bibitem{Phase1} Phase-1 archive, \emph{On the irrationality of $e+\pi$}
(formalization project), Lean~4 modules \texttt{EpiFormal/\allowbreak Classification.lean}
(Theorem~2 and Corollary~3) and \texttt{EpiFormal/\allowbreak Spectral\allowbreak Gap.lean}, 2026-11-15.
\end{thebibliography}

\end{document}
